% =============================================================================
%  MANUAL DO GERADOR DE MISSÕES UNOORIA
%  Identidade visual conforme o brandbook Unooria (abril/2026):
%    #0e1039 azul profundo (estrutura) · #8b85ac lavanda (inteligência)
%    #7fa5e2 azul claro (interface) · preto · branco
%  Tipografia: Garet é a fonte institucional; na composição deste documento
%  usamos Montserrat (geométrica, mesma família de desenho) para títulos e
%  Source Sans Pro para corpo de texto, por disponibilidade em TeX.
%  Compilar com:  tectonic -X compile manual-gerador-missoes.tex
% =============================================================================
\documentclass[11pt,a4paper]{article}

\usepackage[brazil]{babel}
\usepackage[default]{sourcesanspro}          % corpo de texto
\usepackage{fontspec}
\setsansfont{Montserrat-Regular.otf}[        % títulos: geométrica, no espírito da Garet
  BoldFont=Montserrat-Bold.otf,
  ItalicFont=Montserrat-Italic.otf,
  BoldItalicFont=Montserrat-BoldItalic.otf]
\setmonofont{SourceCodePro-Regular.otf}[BoldFont=SourceCodePro-Bold.otf,Scale=0.90]
\usepackage{microtype}
\usepackage{amssymb}
\usepackage{needspace}

\usepackage[a4paper,top=2.6cm,bottom=2.4cm,left=2.4cm,right=2.4cm,headheight=16pt,footskip=1.1cm]{geometry}
\usepackage[table,dvipsnames]{xcolor}
\usepackage{graphicx}
\usepackage{booktabs,tabularx,longtable,array,multirow}
\usepackage{enumitem}
\usepackage{fancyhdr}
\usepackage{titlesec}
\usepackage{titletoc}
\usepackage{caption}
\usepackage{tcolorbox}
\tcbuselibrary{skins,breakable}
\usepackage{tikz}
\usetikzlibrary{shapes.geometric,shapes.symbols,arrows.meta,positioning,fit,backgrounds,calc,decorations.pathreplacing,patterns}
\usepackage{listings}
\usepackage{eso-pic}
\usepackage[hidelinks,pdfusetitle]{hyperref}

% ------------------------------------------------------------------ PALETA ---
\definecolor{unNavy}{HTML}{0E1039}     % azul profundo — estrutura
\definecolor{unLavanda}{HTML}{8B85AC}  % lavanda — inteligência
\definecolor{unAzul}{HTML}{7FA5E2}     % azul claro — interface
\definecolor{unPreto}{HTML}{000000}
\definecolor{unBranco}{HTML}{FFFFFF}
% neutros derivados (tintas da própria paleta, não cores novas)
\definecolor{unPapel}{HTML}{F5F5FA}
\definecolor{unLinha}{HTML}{D9D7E4}
\definecolor{unTexto}{HTML}{2B2D45}
% cores funcionais de sinalização — não são cores de marca; existem só para
% reproduzir os estados que os relatórios da esteira usam (verde, âmbar, vermelho).
\definecolor{sinalVerde}{HTML}{2E6F5E}
\definecolor{sinalAmbar}{HTML}{B07A22}
\definecolor{sinalVermelho}{HTML}{A33B33}

% ------------------------------------------------------------------ TÍTULOS ---
\titleformat{\section}[hang]
  {\sffamily\Large\bfseries\color{unNavy}}
  {\makebox[1.15cm][l]{\color{unLavanda}\thesection}}{0pt}{}
  [{\vspace{2pt}\color{unLavanda}\titlerule[0.8pt]}]
\titleformat{\subsection}[hang]
  {\sffamily\large\bfseries\color{unNavy}}
  {\makebox[1.15cm][l]{\color{unAzul}\thesubsection}}{0pt}{}
\titleformat{\subsubsection}[hang]
  {\sffamily\normalsize\bfseries\color{unNavy}}
  {\makebox[1.15cm][l]{\color{unLavanda}\thesubsubsection}}{0pt}{}
\titlespacing*{\section}{0pt}{2.2em}{1.0em}
\titlespacing*{\subsection}{0pt}{1.5em}{0.6em}
\titlespacing*{\subsubsection}{0pt}{1.2em}{0.45em}

\titlecontents{section}[1.9em]
  {\vspace{2.5pt}\sffamily\bfseries\small\color{unNavy}}
  {\contentslabel[\thecontentslabel]{1.9em}}{}
  {\titlerule*[0.7pc]{.}\contentspage}
\titlecontents{subsection}[3.9em]
  {\footnotesize\color{unTexto}}
  {\contentslabel[\thecontentslabel]{2.4em}}{}
  {\titlerule*[0.7pc]{.}\contentspage}

% ------------------------------------------------------- CABEÇALHO E RODAPÉ ---
\pagestyle{fancy}
\fancyhf{}
\renewcommand{\headrulewidth}{0pt}
\renewcommand{\footrulewidth}{0pt}
\fancyhead[L]{\raisebox{-2pt}{\includegraphics[height=8pt]{marca/logotipo-navy.png}}\hspace{6pt}%
  \textcolor{unLinha}{\rule[1pt]{0.4pt}{9pt}}\hspace{6pt}%
  {\sffamily\scriptsize\color{unLavanda}Manual do Gerador de Missões}}
\fancyhead[R]{\sffamily\scriptsize\color{unLavanda}\nouppercase{\leftmark}}
\fancyfoot[L]{\sffamily\scriptsize\color{unLavanda}Unooria · uso interno}
\fancyfoot[R]{\sffamily\small\bfseries\color{unNavy}\thepage}
\renewcommand{\sectionmark}[1]{\markboth{#1}{}}

\fancypagestyle{parte}{\fancyhf{}\renewcommand{\headrulewidth}{0pt}}

% ---------------------------------------------------------------- AMBIENTES ---
\newtcolorbox{caixaConceito}[1]{
  enhanced,breakable,colback=unLavanda!8,colframe=unLavanda,coltitle=white,
  fonttitle=\sffamily\bfseries\small,title={#1},arc=2mm,boxrule=0pt,
  leftrule=2.4pt,toprule=0pt,rightrule=0pt,bottomrule=0pt,
  attach boxed title to top left={xshift=6mm,yshift=-2.6mm},
  boxed title style={colback=unLavanda,arc=1mm,boxrule=0pt},
  left=10pt,right=10pt,top=11pt,bottom=9pt}

\newtcolorbox{caixaNorma}[1]{
  enhanced,breakable,colback=unAzul!10,colframe=unAzul,coltitle=unNavy,
  fonttitle=\sffamily\bfseries\small,title={#1},arc=2mm,boxrule=0pt,
  leftrule=2.4pt,toprule=0pt,rightrule=0pt,bottomrule=0pt,
  attach boxed title to top left={xshift=6mm,yshift=-2.6mm},
  boxed title style={colback=unAzul,arc=1mm,boxrule=0pt},
  left=10pt,right=10pt,top=11pt,bottom=9pt}

\newtcolorbox{caixaGuardrail}[1]{
  enhanced,breakable,colback=sinalVermelho!7,colframe=sinalVermelho,coltitle=white,
  fonttitle=\sffamily\bfseries\small,title={#1},arc=2mm,boxrule=0pt,
  leftrule=2.4pt,toprule=0pt,rightrule=0pt,bottomrule=0pt,
  attach boxed title to top left={xshift=6mm,yshift=-2.6mm},
  boxed title style={colback=sinalVermelho,arc=1mm,boxrule=0pt},
  left=10pt,right=10pt,top=11pt,bottom=9pt}

\newtcolorbox{caixaPratica}[1]{
  enhanced,breakable,colback=unNavy!5,colframe=unNavy,coltitle=white,
  fonttitle=\sffamily\bfseries\small,title={#1},arc=2mm,boxrule=0pt,
  leftrule=2.4pt,toprule=0pt,rightrule=0pt,bottomrule=0pt,
  attach boxed title to top left={xshift=6mm,yshift=-2.6mm},
  boxed title style={colback=unNavy,arc=1mm,boxrule=0pt},
  left=10pt,right=10pt,top=11pt,bottom=9pt}

% caixa curta sem título, para destaques de uma frase
\newtcolorbox{destaque}{enhanced,breakable,colback=unPapel,colframe=unLavanda,
  arc=2mm,boxrule=0pt,leftrule=2.4pt,toprule=0pt,rightrule=0pt,bottomrule=0pt,
  left=10pt,right=10pt,top=8pt,bottom=8pt}

% --------------------------------------------------------------- COMANDOS ----
\newcommand{\cod}[1]{{\ttfamily\small #1}}
\newcommand{\Verde}[1]{\textcolor{sinalVerde}{\textbf{#1}}}
\newcommand{\Ambar}[1]{\textcolor{sinalAmbar}{\textbf{#1}}}
\newcommand{\Vermelho}[1]{\textcolor{sinalVermelho}{\textbf{#1}}}
% Os relatórios da esteira usam 🟢 🟡 🔴 ⚪. Nenhuma fonte de texto traz esses glifos e o
% XeTeX os descarta com um aviso fácil de perder — por isso são desenhados.
\newcommand{\bolinha}[1]{\tikz[baseline=-0.5ex]{\fill[#1] (0,0) circle (0.33em);}}
\newcommand{\statusVerde}{\bolinha{sinalVerde}}
\newcommand{\statusAmbar}{\bolinha{sinalAmbar}}
\newcommand{\statusVermelho}{\bolinha{sinalVermelho}}
\newcommand{\statusPendente}{\tikz[baseline=-0.5ex]{\draw[unLavanda,line width=0.5pt] (0,0) circle (0.30em);}}
\newcommand{\selo}[2]{\tikz[baseline=(x.base)]{\node[fill=#1,rounded corners=1.6pt,inner xsep=3.6pt,inner ysep=1.6pt,text=white,font=\sffamily\bfseries\scriptsize] (x) {#2};}}
\newcommand{\Evidencia}{\selo{unNavy}{EVIDÊNCIA}}
\newcommand{\Contrato}{\selo{unAzul}{CONTRATO}}
\newcommand{\Convencao}{\selo{unLavanda}{CONVENÇÃO}}
\newcommand{\simbolo}[1][0.9cm]{\includegraphics[width=#1]{marca/simbolo-lavanda.png}}

% abertura de parte: página navy, sangrada, com o símbolo da marca
\newcommand{\parteUnooria}[3]{%
  \clearpage
  \thispagestyle{parte}
  \AddToShipoutPictureBG*{\AtPageLowerLeft{\color{unNavy}\rule{\paperwidth}{\paperheight}}}
  \AddToShipoutPictureBG*{\AtPageLowerLeft{\hspace{0.58\paperwidth}\raisebox{0.11\paperheight}{%
    \includegraphics[width=0.50\paperwidth]{marca/simbolo-grafismo.png}}}}
  \null\vspace*{0.20\paperheight}
  {\sffamily\noindent\color{unAzul}\large\bfseries PARTE #1}\\[0.35em]
  {\color{unLavanda}\rule{3.2cm}{1.6pt}}\\[0.9em]
  {\sffamily\noindent\color{white}\fontsize{27}{31}\selectfont\bfseries #2}\\[0.9em]
  {\noindent\begin{minipage}{0.74\textwidth}\color{unLavanda}\large\itshape #3\end{minipage}}
  \vfill
  \noindent\includegraphics[width=3.1cm]{marca/logotipo-lavanda.png}
  \clearpage}

\newcommand{\parteApendice}[2]{%
  \clearpage
  \thispagestyle{parte}
  \AddToShipoutPictureBG*{\AtPageLowerLeft{\color{unNavy}\rule{\paperwidth}{\paperheight}}}
  \AddToShipoutPictureBG*{\AtPageLowerLeft{\hspace{0.58\paperwidth}\raisebox{0.11\paperheight}{%
    \includegraphics[width=0.50\paperwidth]{marca/simbolo-grafismo.png}}}}
  \null\vspace*{0.20\paperheight}
  {\color{unLavanda}\rule{3.2cm}{1.6pt}}\\[0.9em]
  {\sffamily\noindent\color{white}\fontsize{27}{31}\selectfont\bfseries #1}\\[0.9em]
  {\noindent\begin{minipage}{0.74\textwidth}\color{unLavanda}\large\itshape #2\end{minipage}}
  \vfill
  \noindent\includegraphics[width=3.1cm]{marca/logotipo-lavanda.png}
  \clearpage}

\captionsetup{font={small,sf},labelfont={bf,color=unNavy},textfont={color=unTexto},width=0.92\textwidth}

\lstset{
  basicstyle=\ttfamily\footnotesize\color{unTexto},
  backgroundcolor=\color{unPapel},
  frame=single,rulecolor=\color{unLinha},framesep=6pt,
  breaklines=true,columns=fullflexible,showstringspaces=false,
  xleftmargin=2pt,framexleftmargin=2pt,aboveskip=1.1em,belowskip=1.1em}

\renewcommand{\arraystretch}{1.22}
\arrayrulecolor{unLinha}
\setlist[itemize]{leftmargin=1.1em,itemsep=2pt,topsep=4pt}
\setlist[enumerate]{leftmargin=1.3em,itemsep=2pt,topsep=4pt}
\setlength{\parskip}{0.45em}
\setlength{\parindent}{0pt}
\color{unTexto}

\hypersetup{pdftitle={Manual do Gerador de Missões Unooria},
  pdfauthor={Unooria — Produto, Design e Engenharia},
  pdfsubject={Arquitetura multiagente, operação e fundamentação da esteira de missões}}

\begin{document}
\shorthandoff{"}   % aspas retas aparecem em trechos de código

% =============================================================================
% CAPA
% =============================================================================
\begin{titlepage}
\thispagestyle{empty}
\AddToShipoutPictureBG*{\AtPageLowerLeft{\color{unNavy}\rule{\paperwidth}{\paperheight}}}
% grafismo: símbolo da marca em tinta sobre o fundo, respeitando o safespace
\AddToShipoutPictureBG*{\AtPageLowerLeft{\hspace{0.40\paperwidth}\raisebox{0.045\paperheight}{%
  \includegraphics[width=0.66\paperwidth]{marca/simbolo-grafismo.png}}}}

\null\vspace*{0.055\paperheight}
\noindent\includegraphics[width=5.0cm]{marca/logotipo-claro.png}

\vspace{0.13\paperheight}

{\sffamily\noindent\color{unAzul}\normalsize\bfseries
 PLATAFORMA UNOORIA DE APRENDIZAGEM INTERATIVA}\\[0.5em]
{\color{unLavanda}\rule{4.2cm}{1.8pt}}

\vspace{0.9em}
{\sffamily\noindent\color{white}\fontsize{40}{44}\selectfont\bfseries
 Manual do\\ Gerador de Missões}

\vspace{1.1em}
{\noindent\color{unLavanda}\fontsize{14}{19}\selectfont\itshape
 A esteira de 72 passos que transforma um briefing de meia página\\
 em um PRD pronto para codificação e em uma missão que a\\ professora aprova --- o que ela faz, como usá-la e por que é assim.}

\vfill

{\sffamily\noindent\color{unAzul}\footnotesize\bfseries PARA QUEM ESTE MANUAL É}\\[0.3em]
{\noindent\color{white}\small
 Coordenação pedagógica e professoras consultoras · Product Owners · Designers instrucionais e de experiência ·
 Engenharia de software · Quem opera a esteira no dia a dia}

\vspace{1.0em}
{\color{unLavanda}\rule{\textwidth}{0.6pt}}
\vspace{0.6em}

\noindent
\begin{minipage}[t]{0.49\textwidth}
{\sffamily\footnotesize\color{unAzul}\bfseries EDIÇÃO}\\[0.25em]
{\color{white}\small 2.\textsuperscript{a} edição --- 25 de setembro de 2026\\
Repositório \cod{gerador-missoes} · rubrica de profundidade v1.21}
\end{minipage}\hfill
\begin{minipage}[t]{0.47\textwidth}
{\sffamily\footnotesize\color{unAzul}\bfseries ÂMBITO}\\[0.25em]
{\color{white}\small 1.\textsuperscript{o} ano do Fundamental ao 3.\textsuperscript{o} do Médio\\
BNCC Computação (Resolução CNE/CEB n.\textsuperscript{o} 1/2022) · LGPD, art. 14 · ECA}
\end{minipage}

\vspace{1.4em}
{\noindent\color{unLavanda}\scriptsize Unooria Educação e Tecnologia --- documento de trabalho, uso interno.}
\end{titlepage}

% =============================================================================
% FICHA TÉCNICA
% =============================================================================
\thispagestyle{empty}
\null\vspace*{1.2cm}

{\sffamily\color{unNavy}\Large\bfseries Ficha técnica desta edição}\\[0.3em]
{\color{unLavanda}\rule{3.2cm}{1.6pt}}

\vspace{1.2em}
\begin{tabularx}{\textwidth}{>{\sffamily\bfseries\color{unNavy}\footnotesize}p{4.6cm}X}
\toprule
O que este documento é & O manual único do Gerador de Missões. Nasce da fusão de dois documentos anteriores: um voltado à operação (como escrever a entrada, como ler o diagnóstico, o que cada portão cobra) e outro à fundamentação (por que a arquitetura é multiagente, o que sustenta a régua pedagógica, o que a lei exige). Os dois eram corretos e incompletos em direções opostas. \\
\addlinespace
Como foi verificado & Toda descrição de agente, portão, variável de ambiente e artefato foi conferida contra o código do repositório --- \cod{agents/}, \cod{config.py}, \cod{prompts/} --- e não contra os manuais anteriores. Onde os dois divergiam, o código decidiu. \\
\addlinespace
Arquitetura descrita & \textbf{72 passos numerados} na ordem de execução --- 37 agentes com modelo de linguagem e 35 verificações determinísticas em Python ---, oito ciclos de correção automática e \textbf{parecer humano entre as etapas}. Orquestração em Google ADK (\emph{Agent Development Kit}), modelo padrão \cod{gemini-3.8-flash}. \\
\addlinespace
O que mudou nesta edição & A 1.\textsuperscript{a} edição (19/09/2026) descrevia 55 passos e uma esteira que rodava sozinha do briefing ao PRD. Desde então entraram: o \textbf{parecer humano entre etapas} com login por papel, a camada de \textbf{profundidade pedagógica} (mapa conceitual, ideação e juiz do aha), a \textbf{memória entre execuções}, as \textbf{perguntas de produto} e o \textbf{recorte da aula} por disciplina. São 114 \emph{commits} de diferença. \\
\addlinespace
Identidade visual & Paleta e tipografia do \emph{brandbook} Unooria: azul profundo \cod{\#0e1039} (estrutura), lavanda \cod{\#8b85ac} (inteligência), azul claro \cod{\#7fa5e2} (interface). A fonte institucional é a Garet; nesta composição usamos Montserrat, da mesma família de desenho geométrico, por disponibilidade tipográfica. Verde, âmbar e vermelho aparecem apenas como \textbf{cores funcionais de sinalização}, para reproduzir os três estados que os relatórios da esteira emitem; não são cores de marca e nunca se aplicam ao logotipo. \\
\bottomrule
\end{tabularx}

\vspace{1.8em}
\begin{destaque}
{\sffamily\bfseries\color{unNavy}Uma convenção de leitura.} Ao longo do manual, cada regra da esteira aparece marcada pela natureza da sua fundamentação:\\[0.5em]
\Evidencia\ deriva de pesquisa publicada, com citação. \quad
\Contrato\ deriva de documento normativo --- BNCC, LGPD, manifesto da plataforma; não se pesquisa, se confere. \quad
\Convencao\ a equipe escolheu por consistência de produto, sem base empírica.\\[0.5em]
Convenção é legítima. Convenção apresentada como achado, não.
\end{destaque}

\vspace{1.4em}
\noindent\includegraphics[width=3.4cm]{marca/logotipo-navy.png}

\clearpage

% =============================================================================
% SUMÁRIO
% =============================================================================
\thispagestyle{empty}
{\sffamily\color{unNavy}\Large\bfseries Sumário}\\[0.3em]
{\color{unLavanda}\rule{3.2cm}{1.6pt}}
\vspace{0.6em}
\begingroup
\renewcommand{\contentsname}{}
\setcounter{tocdepth}{2}
\tableofcontents
\endgroup
\clearpage

% =============================================================================
% MAPA DE LEITURA
% =============================================================================
\section*{Por onde começar}
\addcontentsline{toc}{section}{Por onde começar}
\markboth{Por onde começar}{}

Este manual tem sete partes e não precisa ser lido na ordem. O caminho mais curto depende do que você faz com a esteira.

\vspace{0.6em}
\begin{center}
\begin{tikzpicture}[
  perfil/.style={rectangle,rounded corners=3pt,draw=unLavanda,fill=unPapel,
    text width=3.05cm,align=center,font=\sffamily\scriptsize\bfseries,
    minimum height=0.95cm,inner sep=4pt,text=unNavy},
  rota/.style={rectangle,rounded corners=3pt,fill=unNavy,text=white,
    text width=3.05cm,align=center,font=\sffamily\scriptsize,minimum height=1.45cm,inner sep=4pt},
  seta/.style={-{Stealth[length=2.2mm]},draw=unLavanda,thick}]
\node[perfil] (a) {Coordenação\\ e professoras};
\node[perfil,right=0.45cm of a] (b) {Product Owner\\ e produto};
\node[perfil,right=0.45cm of b] (c) {Design\\ instrucional e UX};
\node[perfil,right=0.45cm of c] (d) {Engenharia\\ e dados};
\node[rota,below=0.55cm of a] (a2) {Partes \textbf{I}, \textbf{IV} e \textbf{V}:\\ o que a missão entrega, o que sustenta a régua, o que a lei exige};
\node[rota,below=0.55cm of b] (b2) {Partes \textbf{II} e \textbf{VII}:\\ como escrever a entrada, como ler o diagnóstico, o que a operação ensinou};
\node[rota,below=0.55cm of c] (c2) {Partes \textbf{II} e \textbf{III}:\\ o contrato da \emph{spec}, o QA de entregáveis, o que Design pode e não pode decidir};
\node[rota,below=0.55cm of d] (d2) {Partes \textbf{III} e \textbf{VI}:\\ os portões determinísticos, o contrato de telemetria, o PRD e seu gate};
\draw[seta] (a)--(a2); \draw[seta] (b)--(b2); \draw[seta] (c)--(c2); \draw[seta] (d)--(d2);
\end{tikzpicture}
\end{center}

\vspace{0.3em}
Se você tem quinze minutos e precisa de um só capítulo, leia a Seção~\ref{sec:entrada} --- \textbf{como escrever a entrada}. É a parte do sistema com maior efeito sobre a qualidade do que sai, e a única que depende inteiramente de você.

\vspace{0.4em}
\begin{center}
\begin{tabularx}{\textwidth}{>{\sffamily\bfseries\color{unNavy}\footnotesize}p{2.05cm}>{\sffamily\bfseries\footnotesize\color{unNavy}}p{5.0cm}>{\raggedright\arraybackslash\footnotesize}X}
\toprule
\multicolumn{1}{l}{\sffamily\bfseries\footnotesize\color{unNavy}Parte} & \multicolumn{1}{l}{\sffamily\bfseries\footnotesize\color{unNavy}Título} & \multicolumn{1}{l}{\sffamily\bfseries\footnotesize\color{unNavy}Responde a} \\
\midrule
I & O que a esteira faz & Que documentos saem, para quem, e o que significa a esteira não abortar. \\
II & Como usar & O que colocar no briefing, como rodar, como ler o diagnóstico quando algo trava. \\
III & Como é construída & Por que 72 passos e não um prompt grande; portões, ciclos e o controle anti-alucinação. \\
IV & Por que funciona & A jornada de seis etapas, a carga cognitiva e a régua por idade --- com os números e as fontes. \\
V & O que protege & BNCC e o primado do letramento digital, LGPD art.~14, ética de acolhimento. \\
VI & Os agentes, um a um & O catálogo completo, com o que cada um faz e em que se apoia. \\
VII & O que a operação ensinou & As correções que só apareceram ao rodar a esteira inteira sobre material real. \\
\bottomrule
\end{tabularx}
\end{center}

\clearpage

\parteUnooria{I}{O que a esteira faz}{Meia página de briefing entra. Sete documentos saem --- e uma pessoa responde por cada etapa antes da seguinte começar.}

\section{Os documentos que saem}
\label{sec:artefatos}

Você escreve um briefing de meia página sobre a missão que quer, anexa o material do bimestre e aperta um botão. A esteira roda a etapa de \textbf{Produto} --- cerca de quinze minutos --- e \textbf{para, esperando o seu parecer}. Só depois de você dar o \emph{de acordo} ela segue para Design, e de novo para Dev. Ao fim, a pasta tem tudo o que é preciso para construir aquela missão, mais o relatório honesto do que não deu certo.

Essa parada é a mudança mais importante desde a primeira edição deste manual, e a Seção~\ref{sec:passagem} a descreve em detalhe. O resumo: a máquina julga o que é mecânico, e uma pessoa responde por cada etapa antes da seguinte começar.

\begin{center}
\begin{tikzpicture}[
  ent/.style={rectangle,rounded corners=3pt,draw=unLavanda,fill=unPapel,text=unNavy,
    text width=2.75cm,align=center,font=\sffamily\scriptsize,minimum height=1.05cm,inner sep=4pt},
  cx/.style={rectangle,rounded corners=4pt,fill=unNavy,text=white,
    text width=3.3cm,align=center,font=\sffamily\small\bfseries,minimum height=3.9cm,inner sep=6pt},
  said/.style={rectangle,rounded corners=3pt,draw=unAzul,fill=unAzul!12,text=unNavy,
    text width=3.35cm,align=left,font=\sffamily\scriptsize,minimum height=0.78cm,inner sep=4pt},
  seta/.style={-{Stealth[length=2.2mm]},draw=unLavanda,thick}]

\node[ent] (b) {\textbf{Briefing}\\ meia página};
\node[ent,below=0.3cm of b] (m) {\textbf{Material}\\ do bimestre\\ (\cod{.docx}, \cod{.pdf})};
\node[ent,below=0.3cm of m] (r) {\textbf{Restrições}\\ de sala};

\node[cx,right=1.0cm of m] (e) {72 passos\\[2pt]
  {\sffamily\scriptsize\mdseries 3 etapas\\ 8 ciclos de correção\\ 3 pareceres humanos}};

\node[said,right=1.0cm of e,yshift=1.95cm] (s1) {\textbf{one-pager-resumido}\\ o que o time de produto lê};
\node[said,below=0.16cm of s1] (s2) {\textbf{one-pager-pedagogico}\\ o documento de aprovação};
\node[said,below=0.16cm of s2] (s3) {\textbf{PRD.md}\\ 25 seções, para quem codifica};
\node[said,below=0.16cm of s3] (s4) {\textbf{diagnostico-execucao}\\ o que travou e de quem é};
\node[said,below=0.16cm of s4] (s5) {\textbf{08-prototipo.html}\\ as telas, navegáveis};
\node[said,below=0.16cm of s5] (s6) {\textbf{09-diagramas.html}\\ jornada e fluxo};
\node[said,below=0.16cm of s6] (s7) {\textbf{aprendizado} + \textbf{catálogo}\\ a memória entre execuções};

\draw[seta] (b.east) to[out=0,in=180] (e.west);
\draw[seta] (m.east) -- (e.west);
\draw[seta] (r.east) to[out=0,in=180] (e.west);
\foreach \s in {s1,s2,s3,s4,s5,s6,s7} \draw[seta] (e.east) to[out=0,in=180] (\s.west);
\end{tikzpicture}
\captionof{figure}{O contrato da ferramenta: o que ela exige de você e o que devolve.}
\end{center}

\vspace{0.4em}
\begin{center}
\begin{tabularx}{\textwidth}{>{\ttfamily\scriptsize\color{unNavy}\raggedright\arraybackslash}p{4.3cm}>{\raggedright\arraybackslash\footnotesize}X}
\toprule
\rmfamily\sffamily\bfseries\footnotesize\color{unNavy}Artefato & \sffamily\bfseries\footnotesize\color{unNavy}Para quem é \\
\midrule
produto/\allowbreak 00d-one-pager-\allowbreak resumido \newline (\cod{.md}, \cod{.docx}, \cod{.html}) & \textbf{O time de produto}, para a aprovação pedagógica. Formato curto: ficha, ideia, componente curricular do bimestre, socioemocionais, as seis etapas, o papel do professor, o contrato de evidência e o \emph{takeaway}. Nasceu da comparação com o documento que o time escrevia à mão --- ver Seção~\ref{sec:resumido}. \\
\addlinespace
one-pager-\allowbreak pedagogico \newline (\cod{.md}, \cod{.docx}, \cod{.html}) & Coordenação e a professora consultora. Linguagem de sala de aula, zero jargão de software. É o documento de aprovação: se ele não convence, a missão não vai para a escola. \\
\addlinespace
PRD.md & O agente --- ou a pessoa --- que vai codificar. 25 seções, com contratos de dados, eventos de telemetria, modelo relacional e plano de implementação em caminhos reais do monorepo. \\
\addlinespace
diagnostico-\allowbreak execucao.md & Você. Diz o que passou, o que passou com ressalva, o que travou, de quem é cada conserto e por onde começar. \textbf{É o primeiro arquivo a abrir.} \\
\addlinespace
design/\allowbreak 08-prototipo.html & Revisão humana. Arquivo único, sem dependências: abre com duplo clique e navega. É o que permite à professora simular a aula antes de aprovar. \\
\addlinespace
design/\allowbreak 09-diagramas.html & Quem quiser \emph{ver} a missão: a jornada por raias e o grafo de telas, remontados do Mermaid. \\
\addlinespace
aprendizado.md & O processo. O que esta execução ensinou, que padrões se repetiram e o que propor de mudança nas rubricas. \\
\addlinespace
catalogo-ficha.md & A memória entre execuções. A ficha desta missão entra no índice de todas as já geradas, e é esse índice que o micro-gate consulta para acusar missão repetida. \\
\bottomrule
\end{tabularx}
\end{center}

Mais quarenta e poucos documentos intermediários na mesma pasta, consultáveis quando algo não fecha: o mapa conceitual do material, os três mecanismos da ideação, a \emph{spec} da missão, a jornada, o fluxo, o dicionário de \emph{copy}, a camada de acessibilidade, a especificação de telas, o \emph{handoff} de dev, o \emph{brief} de assets, a matriz de evidência, as perguntas de produto e o relatório de cada portão. O Apêndice~\ref{ap:artefatos} traz a árvore completa.

\subsection{A regra que muda tudo: a esteira não desiste}

Um portão reprovado \textbf{não} interrompe a execução. Todos os documentos são produzidos até o fim; o que deriva de uma etapa reprovada sai marcado \cod{[PROVISÓRIO]}, e o diagnóstico final diz o que é firme e o que é rascunho.

Isso é deliberado, e custou uma reescrita do orquestrador. Na versão anterior, a primeira reprovação parava tudo: não saía PRD, não saía \emph{one-pager}, não saía relatório --- e quem precisava entender o erro ficava sem material para analisar.

\begin{destaque}
Um sistema que falha sem deixar evidência transfere para a pessoa o trabalho de reconstruir o que aconteceu. É o oposto do que uma esteira de produção deve fazer.
\end{destaque}

\begin{center}
\begin{tikzpicture}[
  cx/.style={rectangle,rounded corners=3pt,text width=6.6cm,align=left,inner sep=9pt,font=\sffamily\scriptsize}]
\newcommand{\tit}{\sffamily\small\bfseries}
\node[cx,draw=sinalVermelho,fill=sinalVermelho!6,text=unNavy] (a)
  {{\tit\color{sinalVermelho}Antes --- parada na primeira reprovação}\\[4pt]
   Portão do gate estrutural reprova no critério 3.2.\\
   $\rightarrow$ Execução abortada.\\
   $\rightarrow$ Nenhum PRD, nenhum \emph{one-pager}, nenhum relatório.\\
   $\rightarrow$ Quem lê precisa reabrir os \emph{logs} para saber o que houve.\\[3pt]
   \textbf{Resultado:} 50 chamadas pagas, zero documento.};
\node[cx,draw=sinalVerde,fill=sinalVerde!6,text=unNavy,right=0.5cm of a] (b)
  {{\tit\color{sinalVerde}Hoje --- a esteira segue e registra}\\[4pt]
   Portão do gate estrutural reprova no critério 3.2.\\
   $\rightarrow$ Ciclo de correção tenta consertar (teto de rodadas).\\
   $\rightarrow$ Não resolvido, a etapa segue com herança degradada.\\
   $\rightarrow$ Tudo que deriva dali sai \cod{[PROVISÓRIO]}; o PRD sai \statusVermelho.\\[3pt]
   \textbf{Resultado:} documentos completos e um diagnóstico que nomeia o dono do conserto.};
\end{tikzpicture}
\captionof{figure}{Por que a esteira deixa de abortar: a diferença não é de tolerância, é de evidência.}
\end{center}

\section{O que é uma missão Unooria}
\label{sec:missao}

Antes de falar da esteira, é preciso ser preciso sobre o produto que ela fabrica. Uma \textbf{missão} é uma experiência imersiva de letramento digital, de uma aula, em que o estudante usa a própria conta dentro de uma simulação fiel de um ambiente digital real e \textbf{descobre agindo} algo que ninguém contou para ele.

\begin{caixaConceito}{A metáfora do estúdio}
O Gerador de Missões não é uma máquina que escreve textos. É um \textbf{estúdio integrado de criação}, com três andares e uma sala de inspeção.

No primeiro andar, roteiristas e pedagogos (\textbf{Produto}) afiam a ideia até que o estudante viva um dilema real, e não uma aula expositiva disfarçada de jogo. No segundo, o estúdio de arte e cenografia (\textbf{Design}) projeta telas, fluxo, linguagem acessível e um protótipo que se abre no navegador. No terceiro, arquitetos e diretores técnicos (\textbf{Dev}) convertem tudo isso em especificação de engenharia, histórias de usuário e conformidade legal.

A sala de inspeção atravessa os três andares: a cada fronteira, um portão confere o trabalho contra um critério declarado --- e quem inspeciona nunca é quem produziu.
\end{caixaConceito}

\subsection{O framework de seis etapas}

Toda missão segue a mesma estrutura narrativa, e ela não é decorativa: é o que garante que a descoberta venha da ação, não do enunciado.

\begin{center}
\begin{tikzpicture}[
  et/.style={rectangle,rounded corners=3pt,draw=unNavy,fill=unPapel,text=unNavy,
    text width=2.25cm,align=center,font=\sffamily\scriptsize,minimum height=1.25cm,inner sep=4pt},
  pico/.style={rectangle,rounded corners=3pt,draw=unNavy,fill=unNavy,text=white,
    text width=2.25cm,align=center,font=\sffamily\scriptsize,minimum height=1.25cm,inner sep=4pt},
  seta/.style={-{Stealth[length=2.4mm]},draw=unLavanda,very thick}]
\node[et] (a) {\textbf{1. Desafiar}\\ a tarefa concreta};
\node[et,right=0.42cm of a] (b) {\textbf{2. Navegar}\\ explorar sem saber};
\node[et,right=0.42cm of b] (c) {\textbf{3. Decidir}\\ a escolha que importa};
\node[pico,below=0.75cm of c] (d) {\textbf{4. Resultado}\\ o ``ahá''};
\node[et,left=0.42cm of d] (e) {\textbf{5. Explicar}\\ por que aconteceu};
\node[et,left=0.42cm of e] (f) {\textbf{6. Refletir}\\ aplicar em cenário novo};
\draw[seta] (a)--(b); \draw[seta] (b)--(c);
\draw[seta] (c.south)--(d.north);
\draw[seta] (d)--(e); \draw[seta] (e)--(f);
\node[draw=unAzul,fill=unAzul!10,rounded corners=3pt,text width=5.1cm,align=center,
      font=\sffamily\scriptsize,text=unNavy,inner sep=6pt,above left=0.15cm and -1.2cm of a.north west] (nota)
  {O estudante que vive o \textbf{Navegar} não pode saber o que o \textbf{Resultado} vai revelar.\\ É o teste mais duro da esteira.};
\draw[-{Stealth[length=1.8mm]},draw=unAzul] (nota.south east) to[out=-60,in=110] (d.north west);
\end{tikzpicture}
\captionof{figure}{As seis etapas. O pico da experiência é o Resultado, e ele só funciona se a virada tiver sido guardada.}
\end{center}

\begin{center}
\begin{tabularx}{\textwidth}{>{\sffamily\bfseries\footnotesize\color{unNavy}}p{2.0cm}>{\footnotesize}p{4.3cm}X>{\footnotesize\itshape}p{2.6cm}}
\toprule
\sffamily\bfseries\footnotesize\color{unNavy}Etapa & \sffamily\bfseries\footnotesize\color{unNavy}O que o estudante faz & \sffamily\bfseries\footnotesize\color{unNavy}O que a etapa precisa garantir & \sffamily\normalfont\bfseries\footnotesize\color{unNavy}Apoio teórico \\
\midrule
Desafiar & Recebe uma tarefa concreta dentro da simulação: ``monte uma seleção de três vídeos sobre o tema para mostrar à turma''. & Um \emph{job-to-be-done}, não um relato. O estudante sabe \textbf{o que} vai fazer --- nunca \textbf{por quê}, porque o porquê é a descoberta. & Merrill (2002) \\
\addlinespace
Navegar & Explora o ambiente com os verbos nativos da plataforma simulada: rolar, curtir, salvar, comentar, buscar. & O \emph{gap} de informação. Nada no texto pode antecipar a virada --- há inclusive uma lista de palavras-sentinela que o portão procura. & Loewenstein (1994) \\
\addlinespace
Decidir & Executa a tarefa: escolhe, publica, envia, seleciona. & Escolha significativa: o estudante percebe que escolhe, a escolha muda algo perceptível, a missão a referencia depois e ele não pode desfazer depois de ver o efeito. & Dewey (1938) \\
\addlinespace
Resultado & Vê a consequência da própria decisão e compara com a de colegas reais da turma. & O ``ahá''. É mediado pela professora, com dados de comparação no painel --- não é uma tela de resposta certa. & Bruner (1961) \\
\addlinespace
Explicar & Formula com as próprias palavras o que aconteceu e por quê. & Que o conceito seja \textbf{nomeado depois} da experiência, nunca antes. & Vygotsky (1978) \\
\addlinespace
Refletir & Responde a uma provocação sobre uma situação nova, que não viveu na missão. & Aplicação, não confirmação. O teste é direto: alguém que \emph{não} fez a missão consegue responder só pela pergunta? Se sim, a provocação não serve. & Anderson (2001) \\
\bottomrule
\end{tabularx}
\end{center}

\subsection{Os cinco ambientes, e por que a lista é fechada}

A missão acontece dentro de uma simulação de plataforma real, e são cinco --- cada uma com uma seção própria no \cod{@unooria/design-system}:

\begin{center}
\begin{tikzpicture}[amb/.style={rectangle,rounded corners=3pt,draw=unLavanda,fill=unPapel,
   text=unNavy,text width=2.5cm,align=center,font=\sffamily\scriptsize,minimum height=1.0cm,inner sep=4pt}]
\node[amb] (a) {\textbf{Instagram}\\ \cod{instanooria}};
\node[amb,right=0.3cm of a] (b) {\textbf{YouTube}\\ \cod{unitube}};
\node[amb,right=0.3cm of b] (c) {\textbf{TikTok}\\ \cod{tiktooria}};
\node[amb,right=0.3cm of c] (d) {\textbf{WhatsApp}\\ \cod{whatsappmobile}};
\node[amb,right=0.3cm of d] (e) {\textbf{ChatGPT}\\ \cod{chatgpt}};
\end{tikzpicture}
\end{center}

Pedir uma missão ``no Roblox'' reprova de imediato, e a razão não é gosto: não existe seção de \emph{design system} para construir aquilo, então não há o que especificar, nem o que prototipar, nem o que codificar. A lista fechada é um \Contrato\ da plataforma, não uma preferência de produto. Nos anos iniciais, quando nenhum dos cinco ambientes serve à intenção pedagógica, a esteira admite uma sexta classe --- o mini-jogo --- decidida pelo agente de ambiente simulado (19).

\subsection{As duas camadas: o que é fixo e o que a professora parametriza}

Toda missão é projetada em duas camadas, e confundi-las é um erro caro de briefing.

\begin{center}
\begin{tikzpicture}[
  cam/.style={rectangle,rounded corners=4pt,text width=7.0cm,align=left,inner sep=10pt,font=\footnotesize}]
\node[cam,fill=unNavy,text=white] (a)
 {{\sffamily\small\bfseries 1. Engine imersiva --- estrutura fixa}\\[4pt]
  O framework das seis etapas, a mecânica da plataforma simulada, os verbos nativos, os \emph{checkpoints} com salvamento de estado e o banco curado padrão. Não muda de missão para missão, e não é parametrizável.};
\node[cam,fill=unLavanda!15,draw=unLavanda,text=unNavy,right=0.5cm of a] (b)
 {{\sffamily\small\bfseries 2. Parametrização e autoria assistida}\\[4pt]
  Na preparação da aula, a professora insere tema, contexto e foco curricular; a IA da plataforma gera rascunhos de posts, comentários e mídias; e \textbf{ela cura, edita e aprova} antes de liberar para a turma.};
\end{tikzpicture}
\end{center}

A consequência prática aparece já no primeiro agente: onde o conteúdo é customizável, a missão escreve \cod{[tema: parâmetro do professor]} em vez de fixar um assunto. Desenhar conteúdo genérico esperando personalizar depois é um defeito que a esteira detecta tarde e caro.

\subsection{O recorte da aula: a customização é transversal, não temática}
\label{sec:recorte}

O que a professora parametriza não é ``um tema''. É o \textbf{recorte da aula}, e ele tem cinco campos --- disciplina, tema, assuntos, contexto e foco. A diferença importa: significa que a professora de Português, a de Biologia e a de História podem preparar \textbf{a mesma missão} para conteúdos diferentes, porque a mecânica, as cardinalidades, a descoberta central e o currículo de letramento digital ficam fixos.

\begin{center}
\begin{tikzpicture}[
  fixo/.style={rectangle,rounded corners=3pt,fill=unNavy,text=white,text width=3.62cm,
    align=left,font=\sffamily\scriptsize,inner sep=8pt,minimum height=2.3cm},
  var/.style={rectangle,rounded corners=3pt,draw=unLavanda,fill=unPapel,text=unNavy,text width=3.12cm,
    align=left,font=\sffamily\scriptsize,inner sep=7pt,minimum height=2.3cm},
  seta/.style={-{Stealth[length=2.0mm]},draw=unLavanda,thick}]
\node[fixo] (f) {\textbf{Fixo na missão}\\[3pt] A mecânica. As seis etapas. A descoberta central. As cardinalidades. A habilidade de letramento digital.};
\node[var,right=0.42cm of f] (a) {\textbf{Aula de Português}\\[3pt] disciplina, tema e 3 a 6 assuntos da aula dela};
\node[var,right=0.3cm of a] (b) {\textbf{Aula de Ciências}\\[3pt] outro tema, outros assuntos, mesma mecânica};
\node[var,right=0.3cm of b] (c) {\textbf{Aula de História}\\[3pt] outro tema, outros assuntos, mesma mecânica};
\draw[seta] (f)--(a); \draw[seta] (a)--(b); \draw[seta] (b)--(c);
\end{tikzpicture}
\captionof{figure}{O recorte da aula. Toda missão sai com a tabela ``a mesma missão em três aulas'' como prova de que a transversalidade não é retórica.}
\end{center}

Duas decisões sustentam isso, e as duas são \Convencao\ declarada:

\begin{itemize}
  \item \textbf{A professora não declara habilidades BNCC da própria disciplina.} A missão responde pelo currículo de Computação; o conteúdo da aula dela entra como recorte, não como competência a comprovar. Isso evita que a esteira tenha de auditar uma matriz curricular que não é a dela.
  \item \textbf{Toda missão entrega a tabela das três aulas} --- Português, Ciências/Biologia e História. Se a mesma mecânica não sobrevive às três, ela não era transversal; era temática, e o marcador \cod{[LIMITE DE TRANSVERSALIDADE --- dono: produto]} registra isso em vez de esconder.
\end{itemize}

\subsection{A missão tem duas visões, e as duas são entregáveis}

Um erro que a esteira cometeu em três missões seguidas: entregar só a visão do estudante. A tela em que a professora prepara a aula --- escolhe o recorte, revisa o que a IA gerou, aprova e libera para a turma --- era descartada pela própria regra de priorizar o pico e o ``ahá''.

Hoje a \textbf{visão do professor} é entregável obrigatório, com cinco telas fixas (\cod{E0-PROF-01} a \cod{E0-PROF-05}) que aparecem na jornada, num segundo grafo de fluxo, na folha do Stitch, no espelho do Miro e no protótipo. Faltando, o QA de design e a banca cobram; sem tela correspondente no Stitch, a folha grava a lacuna com dono nomeado.

\parteUnooria{II}{Como usar}{A entrada é a única parte do sistema que depende inteiramente de você --- e é a que mais muda o resultado.}

\section{Como escrever a entrada}
\label{sec:entrada}

A esteira é boa em \textbf{estruturar}, \textbf{auditar} e \textbf{formalizar}. Ela é ruim em adivinhar o que você não disse --- e, por desenho, prefere marcar uma lacuna a inventar conteúdo. Toda a qualidade que ela pode acrescentar depende de haver matéria-prima para trabalhar.

\subsection{O que o primeiro agente precisa}

O agente 01 (Núcleo narrativo) é o único da esteira que lê os seus anexos. Ele produz quatro coisas, e precisa de matéria-prima para as quatro:

\begin{center}
\begin{tikzpicture}[
  q/.style={rectangle,rounded corners=3pt,draw=unNavy,fill=unPapel,text=unNavy,
    text width=3.28cm,align=left,font=\sffamily\scriptsize,minimum height=1.5cm,inner sep=7pt}]
\node[q] (a) {\textbf{A descoberta central}\\[2pt] O que o estudante vai saber depois que não sabia antes --- escrito como descoberta, não como objetivo de aula.};
\node[q,right=0.35cm of a] (b) {\textbf{A narrativa}\\[2pt] Quem ele é na simulação, o que está em jogo e por que a tarefa importa.};
\node[q,right=0.35cm of b] (c) {\textbf{O mecanismo}\\[2pt] A cadeia de ações que produz a virada, distribuída nas seis etapas.};
\node[q,right=0.35cm of c] (d) {\textbf{A identificação}\\[2pt] Ano, bimestre e ambiente --- os três entram no identificador da missão.};
\end{tikzpicture}
\end{center}

\subsection{Um briefing que funciona}

\begin{lstlisting}
Missao para o 5o ano do Fundamental, 3o bimestre, no ambiente do
YouTube (unitube).

Tema do bimestre (material Entre Redes anexado): como a recomendacao
de conteudo funciona e por que duas pessoas veem coisas diferentes.

O que quero que a crianca descubra: que o que ela ve no feed nao e
"o que existe", e sim o que o sistema aprendeu que ela costuma assistir.

Ideia de mecanismo: ela monta a playlist de um personagem a partir de
um banco de videos. Depois ve o feed que aquele personagem passa a
receber. Dois colegas com playlists diferentes recebem feeds
diferentes, e isso fica visivel na tela.

Restricoes: 50 minutos, aula unica, dois estudantes por tablet.
A professora precisa conseguir pausar a turma antes da revelacao.

Em aberto: nao sei se o banco deve ter 8 ou 12 videos.
\end{lstlisting}

Cinco coisas fazem esse briefing funcionar:

\begin{enumerate}
  \item \textbf{Ano e bimestre explícitos.} O identificador da missão codifica os dois (Seção~\ref{sec:id}); sem eles, o primeiro portão determinístico já reprova.
  \item \textbf{A descoberta está escrita como descoberta.} Compare: ``a criança vai aprender sobre algoritmos de recomendação'' (objetivo de aula) contra ``o que ela vê não é o que existe'' (descoberta). A segunda pode ser revelada por uma tela; a primeira, não.
  \item \textbf{Há um mecanismo, não só uma intenção.} ``Monta a playlist, vê o feed que resulta, compara com o colega'' é uma cadeia de ações. ``Trabalhar pensamento crítico sobre algoritmos'' não é --- e o gate estrutural tem um critério só para essa diferença.
  \item \textbf{As restrições de sala aparecem.} Tempo, dispositivo compartilhado e controle da professora mudam o desenho, e a esteira audita os três.
  \item \textbf{A dúvida está declarada.} Dizer ``não sei se 8 ou 12'' faz a esteira decidir e \emph{registrar} a decisão, com base declarada. Omitir faria ela decidir sem que você soubesse.
\end{enumerate}

\subsection{Briefings que produzem missão ruim}

\begin{center}
\begin{longtable}{p{6.4cm}p{8.5cm}}
\toprule
\sffamily\bfseries\footnotesize\color{unNavy}Entrada problemática & \sffamily\bfseries\footnotesize\color{unNavy}O que acontece \\
\midrule
\endfirsthead
\multicolumn{2}{l}{\sffamily\footnotesize\itshape\color{unLavanda}continuação} \\
\toprule
\sffamily\bfseries\footnotesize\color{unNavy}Entrada problemática & \sffamily\bfseries\footnotesize\color{unNavy}O que acontece \\
\midrule
\endhead
``Missão sobre \emph{fake news} para o Fundamental.'' &
Sem ano, sem bimestre, sem ambiente, sem mecanismo. A esteira inventa os quatro, e o que ela inventa é o mais genérico possível. \Vermelho{É o pior uso da ferramenta.} \\
\addlinespace
``O estudante navega por um feed com posts falsos e descobre que foram manipulados.'' &
O briefing \textbf{entrega a virada}. Há um portão dedicado a isso (agente 02) e ele vai reprovar --- mas o estrago já aconteceu: a ideia nasceu sem segredo, e refazer o núcleo custa rodadas do orçamento. \\
\addlinespace
``Uma missão que desenvolva pensamento crítico e cidadania digital.'' &
Intenção sem mecanismo. O teste do gate estrutural é literal: se você remover a frase, a missão continua igual? Então a frase não descreve nada construível. \\
\addlinespace
``O estudante avalia 15 perfis e classifica cada um como confiável ou não.'' &
Passa no teste de decisão significativa e reprova no volume e na monotonia: quinze decisões do mesmo tipo excedem o que a faixa sustenta. Vira \Ambar{ressalva} ou reprovação no agente 09. \\
\addlinespace
``Fazer no Roblox, que as crianças amam.'' &
Ambiente fora da lista fechada. Reprovação imediata --- não existe seção de \emph{design system} para construir. \\
\addlinespace
``A criança conversa com a IA e reflete sobre o que aprendeu.'' &
Reflexão genérica. O critério pede uma \emph{provocação concreta}: uma situação nova em que ela aplique o que viveu. O teste é direto --- alguém que não fez a missão consegue responder? \\
\addlinespace
Briefing de oito páginas com todas as telas já desenhadas. &
A esteira vira transcritora. Os portões detectam ``eco'' --- quando um agente repete a entrada em vez de construir --- mas você terá pago cinquenta chamadas para formatar o que já tinha. \\
\bottomrule
\end{longtable}
\end{center}

\subsection{Anexos valem mais que adjetivos}

Anexe o material do bimestre --- o \emph{Entre Redes} em \cod{.docx}, o PDF da coleção, o plano de aula. Ele é a fonte primária de código curricular e de contexto: a esteira lê os códigos da BNCC dali e os confere contra o anexo oficial do MEC. \textbf{Dois parágrafos de material real valem mais que uma página de descrição.}

\begin{caixaNorma}{O que a esteira aceita como anexo}
PDFs entram \textbf{em binário}, aproveitando a capacidade multimodal nativa do modelo --- diagramas, tabelas e \emph{layout} visual são interpretados, não apenas o texto. Documentos Word são convertidos para Markdown estruturado, preservando tabelas, listas e títulos. O limite do lote é de \textbf{10 arquivos} e \textbf{25\,MB} agregados.
\end{caixaNorma}

\subsection{O checklist antes de apertar o botão}

\begin{caixaPratica}{Cinco perguntas que evitam a maioria das reprovações}
\begin{itemize}[leftmargin=1.5em,label={\color{unLavanda}$\square$}]
  \item \textbf{Descoberta sem \emph{spoiler}.} O conceito é algo que o estudante descobre agindo, ou já está escrito no enunciado?
  \item \textbf{Ofício claro.} Qual é o papel que ele assume na narrativa, e esse papel justifica as ferramentas que ele usa?
  \item \textbf{Letramento digital ativo.} Em qual tela ele \emph{pratica} segurança, ética ou discernimento de informação --- e não apenas lê sobre isso?
  \item \textbf{Piso de realidade escolar.} Isso roda numa escola sem dispositivo próprio e com sinal instável em sala?
  \item \textbf{Roda de conversa.} O fechamento permite à professora aprofundar o debate com a turma depois da tela?
\end{itemize}
\end{caixaPratica}

\section{Rodar a esteira}

\subsection{Entrar}

O aplicativo exige login, e não existe modo aberto: sem as credenciais configuradas, ele não sobe. São \textbf{quatro contas fixas, uma por papel} --- \cod{produto}, \cod{design}, \cod{dev} e \cod{admin} ---, sem provedor de identidade externo. A sessão vale quatro horas contadas do login e sobrevive ao recarregar a página.

\begin{caixaNorma}{As contas são papéis, não pessoas}
O registro diz ``Design'', não qual designer. Isso é deliberado: o que precisa ficar registrado é \emph{quem respondia pela etapa}, e a conta \cod{admin} que age por outro papel aparece como ``admin (como Produto)'' --- nunca como o papel que ela imitou.

Quem age em cada etapa: \textbf{Produto} --- \cod{produto} e \cod{design}; \textbf{Design} --- \cod{design}; \textbf{Dev} --- \cod{dev}; o \cod{admin} em todas. Missão nova é criada por \cod{produto}, \cod{design} ou \cod{admin}.

O login é da aplicação, não da rede: os endereços temporários de download continuam acessíveis a quem tiver o link exato.
\end{caixaNorma}

\subsection{O parecer humano entre etapas}
\label{sec:passagem}

Esta é a mudança mais importante desde a primeira edição. A esteira \textbf{não roda mais do briefing ao PRD sozinha}: ela para ao fim de cada etapa e espera quem responde por ela.

\begin{center}
\begin{tikzpicture}[
  et/.style={rectangle,rounded corners=3pt,fill=unNavy,text=white,text width=1.75cm,
    align=center,font=\sffamily\scriptsize\bfseries,minimum height=0.9cm,inner sep=4pt},
  sem/.style={rectangle,rounded corners=3pt,draw=unAzul,fill=unAzul!12,text=unNavy,text width=1.6cm,
    align=center,font=\sffamily\scriptsize,minimum height=0.9cm,inner sep=4pt},
  hum/.style={rectangle,rounded corners=3pt,draw=unLavanda,fill=unLavanda!18,text=unNavy,text width=1.95cm,
    align=center,font=\sffamily\scriptsize,minimum height=0.9cm,inner sep=4pt},
  seta/.style={-{Stealth[length=2.0mm]},draw=unLavanda,thick},
  volta/.style={-{Stealth[length=2.0mm]},draw=sinalAmbar,thick,dashed}]
\node[et] (p) {PRODUTO};
\node[sem,right=0.3cm of p] (g1) {semáforo\\ dos gates};
\node[hum,right=0.3cm of g1] (h1) {parecer do\\ dono da etapa};
\node[hum,right=1.25cm of h1] (h2) {recebimento\\ da etapa seguinte};
\node[et,right=1.05cm of h2] (d) {DESIGN};
\draw[seta] (p)--(g1); \draw[seta] (g1)--(h1);
\draw[seta] (h1)-- node[font=\sffamily\scriptsize,above=1pt,text=sinalVerde]{de acordo} (h2);
\draw[seta] (h2)-- node[font=\sffamily\scriptsize,above=1pt,text=unNavy]{iniciar} (d);
\draw[volta] (h1.south) to[out=-90,in=-90,looseness=1.5] node[font=\sffamily\scriptsize,below,text=sinalAmbar]{refazer: só esta etapa} (p.south);
\draw[volta] (h2.north) to[out=90,in=90,looseness=1.5] node[font=\sffamily\scriptsize,above,text=sinalAmbar]{devolver} (h1.north);
\end{tikzpicture}
\captionof{figure}{A passagem entre duas etapas. O mesmo desenho se repete de Design para Dev; o \emph{de acordo} do Dev fecha a missão.}
\end{center}

São quatro momentos, nesta ordem:

\begin{enumerate}
  \item \textbf{O semáforo, que é da máquina.} Os gates da etapa, \statusVerde\ \statusAmbar\ \statusVermelho, como sempre. Ele não decide nada sozinho --- informa.
  \item \textbf{O parecer do dono da etapa.} \emph{De acordo} libera, e exige um comentário --- que não é formalidade: ele segue para os agentes das etapas seguintes, no bloco ``Notas de passagem''. \emph{Refazer} transforma o comentário em pedido de correção e refaz \textbf{só aquela etapa}, voltando ao parecer da mesma pessoa.
  \item \textbf{O recebimento pelo dono da etapa seguinte.} \emph{Iniciar} dispara a etapa dele. \emph{Devolver}, com comentário, reabre o parecer anterior --- e \textbf{não refaz nada sozinho}: quem recebeu a devolutiva decide entre refazer ou reafirmar o \emph{de acordo}.
  \item \textbf{O fechamento.} O \emph{de acordo} do Dev encerra a missão: diagnóstico, aprendizado e registro no catálogo rodam só aí.
\end{enumerate}

\begin{caixaGuardrail}{Gate vermelho e a exceção do admin}
Com um gate \statusVermelho, o botão \emph{De acordo} \textbf{não aparece}. Só a conta \cod{admin} libera, e a liberação entra no dossiê marcada como \textbf{exceção} --- a esteira não esquece que alguém passou por cima de um bloqueio.

As etapas rodam em segundo plano: dá para fechar a aba, e a missão volta para a fila \emph{Aguardando você} de quem dá o próximo parecer. Um batimento de 30 segundos detecta execução morta --- reinício ou \emph{deploy} no meio --- e restaura o backup. \textbf{Não faça \emph{deploy} com missão gerando.}
\end{caixaGuardrail}

\subsection{Perguntar a uma etapa, sem mudar nada}

Com o fluxo de pareceres ativo, o campo de mensagem da missão não fica travado: ele vira \textbf{conversa com as etapas}. Você escolhe Produto, Design ou Dev e pergunta --- ``por que este gate ficou vermelho?'', ``onde está a decisão sobre o banco de vídeos?'' --- e um agente de fora da esteira responde lendo os documentos daquela etapa.

\begin{itemize}
  \item A conversa \textbf{só lê}. Mudar a missão continua sendo só pelo \emph{Refazer} do cartão de parecer, que passa pelo dono.
  \item Ela fica no histórico da missão, mas \textbf{fora dos artefatos} --- não entra no ZIP, no manifesto nem no PRD.
  \item \textbf{Todos que abrem a missão veem, e qualquer papel pode perguntar.} Uma dúvida do dev sobre a decisão de produto é exatamente o tipo de coisa que precisa ficar visível.
  \item O botão \emph{Rascunhar o comentário do parecer} leva a conversa para o campo do comentário, para o dono revisar antes de assinar.
\end{itemize}

\subsection{Corrigir uma missão já gerada}

Cada mensagem numa conversa \emph{nova} cria uma missão do zero. Para corrigir uma que já existe, abra-a no histórico: o campo de mensagem passa a pedir correções nela.

\begin{enumerate}
  \item Você escreve o pedido, com anexos se precisar.
  \item Um agente de triagem, \textbf{fora da esteira}, decide de qual etapa recomeçar --- \emph{spec}, design, dev ou nenhuma --- e mostra a decisão com o motivo. Você confirma ou troca.
  \item A esteira roda \textbf{só dali em diante}, sobre o estado salvo. Refazer da \emph{spec} custa cerca de vinte minutos; do design, quinze; só o dev, sete.
  \item O pedido entra numa chave de estado que \textbf{todo agente lê}: os autores aplicam no que é seu e registram ``o que mudou nesta versão''; o QA e a banca conferem se foi atendido.
  \item O que seria perdido é copiado para uma pasta oculta antes de sobrescrever --- e há \textbf{desfazer a última correção}, que restaura documentos, estado e conversa.
\end{enumerate}

Com o fluxo de pareceres ativo, a correção livre pelo chat fica desligada, porque o \emph{Refazer} do cartão faz o mesmo passando pelo dono.

\subsection{Rodar uma etapa isolada}

Quando você está iterando sobre um \emph{prompt} ou investigando um defeito, rodar a esteira inteira é desperdício. Há duas formas:

\begin{itemize}
  \item \cod{PIPELINE\_STAGES} roda só as etapas pedidas sobre um estado já salvo;
  \item \cod{scripts/rerodar\_etapa.py} reexecuta Design, Dev ou só o diagnóstico sobre uma sessão salva --- inclusive uma que veio do \emph{bucket} do Cloud Run.
\end{itemize}

\begin{destaque}
\textbf{Rodar só a etapa de Dev sobre um estado salvo custa cerca de dois minutos, contra vinte da execução completa} --- e, o que importa mais, a entrada fica \emph{fixa} entre as tentativas, o que é a única forma de saber se a mudança que você fez teve efeito.
\end{destaque}

\subsection{Configuração}

Os padrões foram calibrados na operação real e raramente precisam mudar. As exceções são \cod{PIPELINE\_\allowbreak STAGES}, que muda a cada sessão de trabalho, e os três interruptores de rigor.

\begin{center}
\begin{longtable}{@{}p{4.9cm}p{1.15cm}p{9.2cm}@{}}
\toprule
\sffamily\bfseries\footnotesize\color{unNavy}Variável & \sffamily\bfseries\footnotesize\color{unNavy}Padrão & \sffamily\bfseries\footnotesize\color{unNavy}O que faz \\
\midrule
\endfirsthead
\toprule
\sffamily\bfseries\footnotesize\color{unNavy}Variável & \sffamily\bfseries\footnotesize\color{unNavy}Padrão & \sffamily\bfseries\footnotesize\color{unNavy}O que faz \\
\midrule
\endhead
\multicolumn{3}{@{}l}{\sffamily\bfseries\scriptsize\color{unLavanda}FLUXO E ACESSO} \\[2pt]
{\ttfamily\scriptsize PARECER\_HUMANO} & \footnotesize true & \footnotesize A esteira para ao fim de cada etapa e espera o parecer do dono. \\
{\ttfamily\scriptsize AUTH\_CREDENCIAIS} & \footnotesize --- & \footnotesize \textbf{Obrigatória.} Hashes das quatro contas. Sem ela o aplicativo não sobe. \\
{\ttfamily\scriptsize AUTH\_SEGREDO\_JWT} & \footnotesize --- & \footnotesize \textbf{Obrigatória.} Chave da sessão. Trocá-la derruba todas as sessões. \\
{\ttfamily\scriptsize AUTH\_HORAS\_SESSAO} & \footnotesize 4 & \footnotesize Validade absoluta da sessão, contada do login. \\
{\ttfamily\scriptsize PIPELINE\_STAGES} & \footnotesize todas & \footnotesize Roda só as etapas listadas, sobre um estado salvo. \\
\addlinespace
\multicolumn{3}{@{}l}{\sffamily\bfseries\scriptsize\color{unLavanda}RIGOR DOS PORTÕES} \\[2pt]
{\ttfamily\scriptsize ESTEIRA\_ESTRITA} & \footnotesize true & \footnotesize Reprovação marca herança degradada e o PRD sai \statusVermelho. Desligado, vira ressalva. \\
{\ttfamily\scriptsize PRODUTO\_MODO\_ESTRITO} & \footnotesize true & \footnotesize Gate de produto reprovado impede a liberação para as etapas seguintes. \\
{\ttfamily\scriptsize BNCC\_CHECK\_ESTRITO} & \footnotesize true & \footnotesize Currículo não conforme bloqueia Design e Dev. \\
{\ttfamily\scriptsize PROFUNDIDADE\_ESTRITA} & \footnotesize true & \footnotesize Profundidade reprovada na última rodada do núcleo fica \statusVermelho\ e o \emph{de acordo} de produto é recusado. Em \cod{false}, rebaixa a \statusAmbar\ --- para usar \textbf{enquanto os cortes da rubrica são calibrados}. \\
{\ttfamily\scriptsize KPIS\_CHECK\_ESTRITO} & \footnotesize true & \footnotesize Incoerência entre indicador, evento e tela impede o PRD. \\
\addlinespace
\multicolumn{3}{@{}l}{\sffamily\bfseries\scriptsize\color{unLavanda}TETOS DE CORREÇÃO} \\[2pt]
{\ttfamily\scriptsize GATE\_ESTRUTURAL\_MAX\_RODADAS} & \footnotesize 3 & \footnotesize Ciclo gate estrutural $\leftrightarrow$ revisor do \emph{one-pager}. \\
{\ttfamily\scriptsize SPEC\_MAX\_RODADAS} & \footnotesize 2 & \footnotesize Ciclo \emph{spec} $\leftrightarrow$ currículo, manifesto, evidência e perguntas. \\
{\ttfamily\scriptsize AVALIACOES\_MAX\_RODADAS} & \footnotesize 1 & \footnotesize Ciclo do parecer de adequação. Em 1, o parecer entra como ressalva. \\
{\ttfamily\scriptsize QA\_DESIGN\_MAX\_RODADAS} & \footnotesize 2 & \footnotesize Ciclo QA de entregáveis $\leftrightarrow$ retrabalho de design. \\
{\ttfamily\scriptsize HANDOFF\_MAX\_RODADAS} & \footnotesize 2 & \footnotesize Ciclo \emph{handoff} $\leftrightarrow$ contrato de telemetria. \\
{\ttfamily\scriptsize BANCA\_MAX\_RODADAS} & \footnotesize 2 & \footnotesize Ciclo banca $\leftrightarrow$ dono do achado. \\
{\ttfamily\scriptsize PRD\_MAX\_RODADAS} & \footnotesize 2 & \footnotesize Ciclo redator do PRD $\leftrightarrow$ consistência do documento. \\
{\ttfamily\scriptsize ORCAMENTO\_REEXECUCOES} & \footnotesize 8 & \footnotesize Teto \textbf{global} de agentes reexecutados na execução inteira. \\
\addlinespace
\multicolumn{3}{@{}l}{\sffamily\bfseries\scriptsize\color{unLavanda}MODELO} \\[2pt]
{\ttfamily\scriptsize GEMINI\_MODEL} & \footnotesize --- & \footnotesize Modelo padrão de todos os agentes: \cod{gemini-3.8-flash}. A variável \cod{GEMINI\_MODEL\_PRD} permite outro só para o redator do PRD. \\
{\ttfamily\scriptsize THINKING\_BUDGET} & \footnotesize 8192 & \footnotesize \emph{Tokens} de raciocínio por chamada. \\
{\ttfamily\scriptsize TEMPERATURA\_JUIZ}, {\ttfamily\scriptsize SEED\_JUIZ} & \footnotesize 0,1 e fixa & \footnotesize Os juízes rodam com temperatura baixa e semente fixa: dois julgamentos do mesmo texto têm de dar o mesmo veredito. \\
\bottomrule
\end{longtable}
\end{center}

\textbf{Qualquer teto em 1 desliga o ciclo correspondente}: a esteira avalia e reporta, sem corrigir. É útil para medir a qualidade bruta de um \emph{prompt}, e péssimo como configuração de produção.

\section{O identificador da missão}
\label{sec:id}

O código da experiência carrega ano e bimestre, e é ele que nomeia tabela, pasta e \emph{branch}. Por isso um portão determinístico confere que os campos concordam entre si --- divergência trava.

\begin{center}
\begin{tikzpicture}[
  peca/.style={rectangle,rounded corners=2pt,font=\ttfamily\large\bfseries,text=white,inner sep=5pt},
  rot/.style={font=\sffamily\scriptsize,text=unNavy,align=center,text width=2.4cm}]
\node[peca,fill=unLavanda] (p1) {exp};
\node[peca,fill=unNavy,right=1pt of p1] (p2) {05};
\node[peca,fill=unAzul,right=1pt of p2] (p3) {ef};
\node[peca,fill=unNavy,right=1pt of p3] (p4) {3};
\node[peca,fill=unLavanda,right=1pt of p4] (p5) {b};
\node[rot,below=0.7cm of p1] (r1) {prefixo fixo};
\node[rot,below=0.7cm of p2] (r2) {5.\textsuperscript{o} ano \\ \cod{series: ["5EF"]}};
\node[rot,below=0.7cm of p3] (r3) {etapa: \cod{ef} ou \cod{em}};
\node[rot,below=0.7cm of p4] (r4) {3.\textsuperscript{o} bimestre \\ \cod{unidade: 3}};
\node[rot,below=0.7cm of p5] (r5) {bimestre};
\foreach \i in {1,...,5} \draw[draw=unLavanda,thin] (p\i.south) -- (r\i.north);
\end{tikzpicture}
\captionof{figure}{Anatomia de \cod{exp05ef3b}. O portão confere que \cod{series} e \cod{unidade} concordam com o código.}
\end{center}

\begin{center}
\begin{tabularx}{0.9\textwidth}{>{\ttfamily\footnotesize\color{unNavy}}p{2.6cm}>{\footnotesize}p{5.0cm}X}
\toprule
\rmfamily\sffamily\bfseries\footnotesize\color{unNavy}Código & \sffamily\bfseries\footnotesize\color{unNavy}Significa & \sffamily\bfseries\footnotesize\color{unNavy}Campos que têm de concordar \\
\midrule
exp01ef1b & 1.\textsuperscript{o} ano EF, 1.\textsuperscript{o} bimestre & \cod{series: ["1EF"]} · \cod{unidade: 1} \\
exp05ef3b & 5.\textsuperscript{o} ano EF, 3.\textsuperscript{o} bimestre & \cod{series: ["5EF"]} · \cod{unidade: 3} \\
exp02em2b & 2.\textsuperscript{o} ano EM, 2.\textsuperscript{o} bimestre & \cod{series: ["2EM"]} · \cod{unidade: 2} \\
\bottomrule
\end{tabularx}
\end{center}

Missão transversal (sem série definida) e ferramenta de professor não têm ano nem bimestre: seguem no formato antigo \cod{exp-NN-<slug>}. No histórico da interface, os registros passaram a começar pelo ano e bimestre --- \cod{2EF-1B · O Detalhe que Muda Tudo} --- lidos do próprio identificador. Quando não há como deduzir, o título sai sem prefixo, em vez de inventar um.

\section{Como ler o diagnóstico}
\label{sec:diagnostico}

O \cod{diagnostico-execucao.md} é o primeiro arquivo a abrir, e tem sempre a mesma ordem. Lê-se em trinta segundos e decide se vale ler o resto.

\begin{center}
\begin{tikzpicture}[
  it/.style={rectangle,rounded corners=2pt,draw=unLinha,fill=unPapel,text=unNavy,
    text width=13.6cm,align=left,font=\footnotesize,inner sep=6pt},
  num/.style={circle,fill=unNavy,text=white,font=\sffamily\scriptsize\bfseries,inner sep=2.6pt,minimum size=0.52cm}]
\node[it] (a) {\textbf{Veredito e contagem} --- impedimentos, ressalvas e decisões humanas, em números.};
\node[it,below=0.1cm of a] (b) {\textbf{Como cada portão se saiu} --- tabela com resultado e rodadas gastas.};
\node[it,below=0.1cm of b] (c) {\textbf{O que reprovou, na ordem} --- a sequência dos fatos, não a lista de culpados.};
\node[it,below=0.1cm of c] (d) {\textbf{O que precisa ser consertado, por dono} --- agrupado por responsável.};
\node[it,below=0.1cm of d] (e) {\textbf{O que só uma pessoa decide} --- o que nenhuma rodada de correção resolve.};
\node[it,below=0.1cm of e] (f) {\textbf{O que saiu provisório} --- os documentos que dependem de algo reprovado.};
\node[it,below=0.1cm of f] (g) {\textbf{Documentos ausentes} --- sintoma de agente que não produziu nada.};
\node[it,below=0.1cm of g] (h) {\textbf{Por onde começar} --- a ação que destrava mais coisa de uma vez.};
\foreach \n/\i in {1/a,2/b,3/c,4/d,5/e,6/f,7/g,8/h} \node[num,left=0.18cm of \i.west] {\n};
\end{tikzpicture}
\captionof{figure}{As oito seções do diagnóstico, sempre nessa ordem.}
\end{center}

\begin{caixaPratica}{As três severidades, e o que cada uma pede de você}
\begin{description}[leftmargin=2.9cm,style=sameline,font=\sffamily\bfseries\color{unNavy}]
  \item[\Vermelho{Impedimento}] O que impede construir. Tem conserto certo e, quase sempre, dono dentro da esteira. Se sobrou, é porque o teto de rodadas ou o orçamento global estourou.
  \item[\Ambar{Ressalva}] Risco aceito e registrado. Não bloqueia, mas \textbf{nenhuma ressalva sobrevivente pode ser omitida do \emph{one-pager}}, mesmo com veredito verde --- há uma regra explícita contra isso.
  \item[Decisão humana] O que a esteira não tem autoridade para resolver: escolha de produto, licença de imagem, homologação de autoria, validação pedagógica. Não é defeito da missão; é informação que faltou de quem forneceu o insumo.
\end{description}
\end{caixaPratica}

\section{Três recomendações de uso}

\begin{description}[leftmargin=1.2cm,style=nextline,font=\sffamily\bfseries\color{unNavy}]
  \item[Comece pelo diagnóstico, não pelo PRD.] Ele diz em trinta segundos se vale ler o resto, e dá a ordem de ataque.
  \item[Trate \cod{[PROVISÓRIO]} como rascunho, nunca como entrega.] O marcador existe porque o documento foi construído sobre algo que um portão reprovou --- ele pode estar certo, mas ninguém verificou.
  \item[Rode uma etapa isolada quando estiver iterando.] \cod{PIPELINE\_STAGES} sobre um estado salvo custa minutos em vez de uma execução inteira, e mantém a entrada fixa entre as tentativas.
\end{description}

\parteUnooria{III}{Como é construída}{Cinquenta e quatro agentes não são um capricho de arquitetura. São a resposta a um modo de falha medido.}

\section{Por que muitos agentes, e não um prompt grande}
\label{sec:decomposicao}

A pergunta aparece sempre, e a resposta não é ``porque especialização é boa''.

\subsection{Densidade de instrução degrada o cumprimento}

A etapa de Produto sozinha carrega cerca de \textbf{74 mil \emph{tokens} de instrução e 51 critérios tabelados}. No \emph{benchmark} IFScale, que mede exatamente isso, os melhores modelos de fronteira chegam a 68\% de acurácia quando precisam seguir 500 instruções ao mesmo tempo~\cite{ifscale2025}, com viés sistemático em favor das instruções apresentadas primeiro. Liu et al.~\cite{liu2024} documentam o efeito irmão em contexto longo: o que está no meio é usado pior que o que está nas pontas.

Há evidência interna, e ela é mais convincente que a externa. O \emph{prompt} do agente 09 abre justificando a própria existência:

\begin{destaque}
\itshape ``Esta avaliação existe porque a rubrica original errava sistematicamente aqui: os critérios de faixa etária competiam por atenção com $\sim$30 critérios de outra natureza.''
\end{destaque}

A versão agregada falhava de modo reprodutível; a decomposta corrigiu. Não é teoria importada --- é um defeito que se observou e se consertou.

\subsection{O critério de decomposição}

Não se divide por etapa cronológica. Divide-se por \textbf{decisão encapsulada}, no sentido clássico de Parnas~\cite{parnas1972}: cada agente responde a uma pergunta que pode ser respondida --- e errada --- independentemente das outras, e emite um envelope JSON que um portão consegue ler.

\begin{center}
\begin{tabularx}{0.93\textwidth}{>{\sffamily\bfseries\footnotesize\color{unNavy}}p{5.0cm}X}
\toprule
\sffamily\bfseries\footnotesize\color{unNavy}Agente & \sffamily\bfseries\footnotesize\color{unNavy}A pergunta que ele responde \\
\midrule
05 Gate estrutural & A missão tem estrutura construível? \\
08 Avaliação pedagógica & Cumpre currículo e protege o estudante? \\
09 Experiência do estudante & Alguém dessa idade consegue viver isso? \\
30 QA de entregáveis & A construção é fiel ao que foi especificado? \\
48 Consistência do PRD & O documento final se contradiz? \\
\bottomrule
\end{tabularx}
\end{center}

\subsection{Redundância não é cobertura}

Três critérios do agente 09 eram testes idênticos a três critérios do agente 05 --- mesma pergunta, mesmo modo de falha. Foram removidos. Knight e Leveson~\cite{knight1986} mostraram, em 27 implementações independentes submetidas a um milhão de testes, que falhas correlacionadas anulam o ganho da redundância. Quando os dois ``avaliadores independentes'' são \emph{prompts} diferentes do \textbf{mesmo modelo}, sobre o \textbf{mesmo texto}, a correlação é ainda maior.

\subsection{O contraponto honesto}

A literatura recente é cética quanto a sistemas multiagente: sob orçamento fixo de \emph{tokens}, um agente único frequentemente iguala ou supera arquiteturas multiagente~\cite{tran2026}. Vale registrar isso num manual da casa, e não escondê-lo.

A condição de contorno declarada pelos próprios autores é que o multiagente vence \emph{``quando a utilização efetiva de contexto de um agente único está degradada''} --- que é exatamente o regime desta esteira, com 74 mil \emph{tokens} de instrução só na primeira etapa. E Cemri et al.~\cite{cemri2025}, sobre mais de 1.600 execuções anotadas, mostram que decompor cria uma classe nova de falha: \textbf{contratos de interface}. É por isso que cada fronteira aqui tem envelope JSON verificado por teste --- a decomposição não sai de graça, e o preço dela é conhecido.

\section{As três etapas e o mapa dos agentes}
\label{sec:etapas}

\begin{center}
\begin{tikzpicture}[
  ag/.style={rectangle,rounded corners=2.5pt,draw=unLavanda!70,fill=unPapel,text=unNavy,
    text width=2.42cm,align=center,font=\sffamily\scriptsize,minimum height=0.92cm,inner sep=3pt},
  det/.style={ag,draw=unAzul,fill=unAzul!13,font=\sffamily\scriptsize\itshape},
  fim/.style={ag,draw=unNavy,fill=unNavy,text=white,font=\sffamily\scriptsize\bfseries},
  seta/.style={-{Stealth[length=2.0mm]},draw=unLavanda,thick},
  faixa/.style={font=\sffamily\scriptsize\bfseries,text=unNavy}]

\node[ag] (p1) {00b--00d\\ memória, mapa do material, ideação};
\node[ag,right=0.3cm of p1] (p2) {01--03c\\ núcleo, \emph{spoiler} e profundidade};
\node[ag,right=0.3cm of p2] (p3) {04--07\\ \emph{one-pager} e gate estrutural};
\node[ag,right=0.3cm of p3] (p4) {08--13\\ avaliações, gate final, \emph{spec}};
\node[det,right=0.3cm of p4] (p5) {14--16b\\ currículo, evidência, perguntas};

\node[ag,below=1.15cm of p1] (d1) {17--20\\ lacunas e ambiente};
\node[ag,right=0.3cm of d1] (d2) {21--25\\ ofício, jornada, fluxo, \emph{copy}, CAA};
\node[ag,right=0.3cm of d2] (d3) {26--28\\ telas e protótipo};
\node[det,right=0.3cm of d3] (d4) {29--32\\ \emph{lint}, artefatos e QA};
\node[ag,right=0.3cm of d4] (v1) {33--35\\ manifesto, \emph{handoff}, assets};

\node[det,below=1.15cm of d1] (v2) {36--38\\ telemetria e lacunas};
\node[ag,right=0.3cm of v2] (v3) {39--44\\ banca de 5 dimensões};
\node[ag,right=0.3cm of v3] (v4) {45--47\\ despacho e PRD};
\node[det,right=0.3cm of v4] (v5) {48--48b\\ consistência do PRD e dos KPIs};
\node[fim,right=0.3cm of v5] (fim) {49--52\\ \emph{one-pager}, diagnóstico, catálogo};

\draw[seta] (p1)--(p2); \draw[seta] (p2)--(p3); \draw[seta] (p3)--(p4); \draw[seta] (p4)--(p5);
\draw[seta] (p5.south) to[out=-90,in=90,looseness=0.9] (d1.north);
\draw[seta] (d1)--(d2); \draw[seta] (d2)--(d3); \draw[seta] (d3)--(d4); \draw[seta] (d4)--(v1);
\draw[seta] (v1.south) to[out=-90,in=90,looseness=0.9] (v2.north);
\draw[seta] (v2)--(v3); \draw[seta] (v3)--(v4); \draw[seta] (v4)--(v5); \draw[seta] (v5)--(fim);

\node[faixa,above=0.08cm of p1.north west,anchor=south west] {PRODUTO};
\node[faixa,above=0.08cm of d1.north west,anchor=south west] {DESIGN};
\node[faixa,above=0.08cm of v2.north west,anchor=south west] {DEV \ \color{unLavanda}(e transversal)};

\node[draw=unAzul,fill=unAzul!13,rounded corners=2pt,font=\sffamily\scriptsize,
      text=unNavy,inner sep=4pt,below=0.55cm of v3.south] {caixa em azul claro = passo \textbf{determinístico}, sem modelo de linguagem};
\end{tikzpicture}
\captionof{figure}{Os 72 passos, numerados na ordem de execução. A Parte~VI descreve cada um.}
\end{center}

\subsection{Produto (00b--16b): a missão pode ser construída?}

Decide \emph{o que} a missão é, e é de longe a etapa mais densa --- 32 dos 72 passos. Ela começa \textbf{antes} do núcleo narrativo, com três passos que a primeira edição deste manual não tinha:

\begin{itemize}
  \item a \textbf{memória} (00b) recupera do índice as missões já geradas mais próximas do ano do briefing, para a nova não repetir uma anterior;
  \item o \textbf{mapa conceitual} (00c) lê o material do bimestre e extrai o conceito limiar, os conceitos de apoio, as crenças que o bimestre derruba e o vocabulário --- \emph{antes} de alguém inventar mecanismo;
  \item a \textbf{ideação} (00d) produz três mecanismos distintos para aquele conceito, pontuados na rubrica de profundidade. O núcleo desenvolve o escolhido.
\end{itemize}

Só então vem o núcleo narrativo, que agora passa por \textbf{dois} juízes em vez de um: o micro-gate de \emph{spoiler} e o juiz de profundidade (Parte~IV). Depois, o \emph{one-pager} de 18 seções, os critérios do gate estrutural, as avaliações em paralelo, a \textbf{\emph{spec}} --- o contrato de 12 blocos --- e, no fim, quatro verificações que a \emph{spec} precisa passar: currículo, manifesto, matriz de evidência e perguntas de produto.

\begin{caixaConceito}{Por que a conformidade curricular é conferida ainda em Produto}
A ficha do manifesto (agente 15) e a conformidade BNCC (14) poderiam rodar em Dev, junto com o resto da verificação técnica. Rodam em Produto de propósito: um campo errado ali custa \textbf{uma reescrita da \emph{spec}}; o mesmo campo errado descoberto em Dev custa \textbf{o descarte do Design inteiro}, porque tudo que Design construiu se apoiou na \emph{spec} errada.
\end{caixaConceito}

\subsection{Design (17--32): como isso se constrói}

Fecha as lacunas da \emph{spec}, decide o ambiente simulado, escreve o \emph{briefing} de ofício, e então constrói: jornada e fluxo em \textbf{diagramas Mermaid}, \emph{copy} em duas passadas, camada de acessibilidade, especificação de telas e um \textbf{protótipo HTML navegável} que abre com duplo clique. Fecha com o QA de entregáveis, que audita fidelidade à \emph{spec} --- e cujos achados voltam só para quem tem \emph{ticket}.

A regra que governa esta etapa inteira cabe numa linha: \textbf{Design preenche ausência, nunca altera presença.} Onde a \emph{spec} decidiu, Design constrói; onde a \emph{spec} calou, Design decide e registra a decisão com base declarada. Se a decisão pertence a produto, vira um pedido de re-\emph{spec}, não uma escolha silenciosa.

\subsection{Dev (33--49): isso é implementável?}

Confere a ficha do manifesto, produz o \emph{handoff} e o \emph{brief} de assets, valida o contrato de telemetria, e submete tudo a uma banca de cinco dimensões em paralelo. Consolida, escreve o PRD e, por último, o \emph{one-pager} pedagógico --- que é o último documento justamente porque precisa dizer a verdade sobre tudo o que veio antes.

\subsection{Transversal (50--52): o que aconteceu, e o que fica}

O diagnóstico determinístico monta o dossiê a partir dos envelopes que os portões gravaram --- roda \textbf{sempre}, inclusive quando tudo o mais falhou, e é essa garantia que sustenta a decisão de não abortar. O agente de aprendizado escreve a leitura em prosa.

O terceiro é novo e resolve um problema antigo: \textbf{o registro no catálogo} (52). Até ele existir, duas missões parecidas só eram percebidas por acaso, ou nunca --- o agente que inventava a missão não via as anteriores. Agora a ficha entra no índice automaticamente ao fim da esteira, com o \emph{status} derivado dos gates. Missão bloqueada \textbf{entra de propósito}: o que se quer evitar é reinventá-la, não escondê-la. \textbf{Promover é ato humano} --- nenhum agente marca uma missão como aprovada.

\section{Os oito ciclos de correção}
\label{sec:ciclos}

Quando um portão reprova, quem decide repetir é o \textbf{orquestrador, em código} --- não uma instrução em linguagem natural dentro do \emph{prompt}, que o modelo pode ou não seguir.

\begin{center}
\begin{tikzpicture}[node distance=0.6cm and 1.25cm,
  cx/.style={rectangle,rounded corners=3pt,draw=unLavanda,fill=unPapel,text=unNavy,
    text width=2.3cm,align=center,font=\sffamily\scriptsize,minimum height=0.95cm,inner sep=4pt},
  port/.style={cx,draw=unAzul,fill=unAzul!13,font=\sffamily\scriptsize\itshape},
  ok/.style={cx,draw=unNavy,fill=unNavy,text=white},
  seta/.style={-{Stealth[length=2.2mm]},draw=unLavanda,thick},
  volta/.style={-{Stealth[length=2.2mm]},draw=sinalAmbar,thick,dashed}]
\node[cx] (a) {\textbf{autor}\\ produz};
\node[port,right=of a] (g) {\textbf{portão}\\ julga};
\node[cx,right=of g] (r) {\textbf{corretor}\\ conserta};
\node[ok,right=1.9cm of r] (s) {aprovado:\\ segue};
\draw[seta] (a)--(g);
\draw[seta] (g)-- node[font=\sffamily\scriptsize,above,text=sinalVermelho] {reprova} (r);
\draw[seta] (g.north) to[out=50,in=130] node[font=\sffamily\scriptsize,above,text=sinalVerde] {aprova} (s.north);
\draw[volta] (r.south) to[out=-90,in=-90] node[font=\sffamily\scriptsize,below,text=sinalAmbar] {reavalia, até o teto de rodadas} (g.south);
\end{tikzpicture}
\captionof{figure}{O padrão comum aos oito ciclos: quem julga não conserta, quem conserta não julga.}
\end{center}

\begin{center}
\begin{longtable}{p{2.3cm}p{4.0cm}p{8.3cm}}
\toprule
\sffamily\bfseries\footnotesize\color{unNavy}Ciclo & \sffamily\bfseries\footnotesize\color{unNavy}Portão $\rightarrow$ corretor & \sffamily\bfseries\footnotesize\color{unNavy}O que resolve \\
\midrule
\endfirsthead
\toprule
\sffamily\bfseries\footnotesize\color{unNavy}Ciclo & \sffamily\bfseries\footnotesize\color{unNavy}Portão $\rightarrow$ corretor & \sffamily\bfseries\footnotesize\color{unNavy}O que resolve \\
\midrule
\endhead
Núcleo & 02 e 02b $\rightarrow$ 01 & \emph{Spoiler}, falta de segredo, escolha sem consequência --- \textbf{e profundidade}: o núcleo só sai do ciclo quando passa nos \emph{dois} juízes. É o ciclo que mais mudou desde a primeira edição. \\
\addlinespace
Gate estrutural & 05 $\rightarrow$ 07 & Os critérios de estrutura, mais o \emph{lint} de contrato de plataforma (04b), que roda a cada rodada. O gate escreve o texto de correção pronto para colar --- é o ciclo mais barato da esteira. \\
\addlinespace
Avaliações & 10b $\rightarrow$ 10c & Os ajustes que o parecer de adequação exigiu. \textbf{Vem desligado} (teto 1): o parecer entra como ressalva e o documento não se move. \\
\addlinespace
\emph{Spec} & 14, 15, 15b e 15e $\rightarrow$ 13 & Código da BNCC inexistente ou de outro ano; campo do manifesto fora do vocabulário; habilidade sem evidência declarada; pergunta de produto sem ação que a responda. Conserto mecânico. \\
\addlinespace
QA de design & 30 e 29a $\rightarrow$ dono do \emph{ticket} & Cada achado tem dono declarado; só os endereçados rodam de novo, e o protótipo é refeito depois deles. O \emph{lint} do protótipo (29a) entra aqui porque vê o que o QA não vê: ele lê o arquivo HTML, não o texto. \\
\addlinespace
\emph{Handoff} & 36 $\rightarrow$ 34 & Contrato de eventos, papéis, canal de sala. A correção vem escrita na mensagem do portão. \\
\addlinespace
Banca & 39--44 $\rightarrow$ dono do achado & Inclui donos da \textbf{etapa anterior}: um erro de jornada detectado em Dev volta para o agente de jornada. \\
\addlinespace
PRD & 48 e 48a $\rightarrow$ 47 & Consistência interna do entregável final: prefixo de tabela, caminho do módulo, requisito P0 sem rastro, evento sem tabela leitora, tarefa não atômica, identificador inválido --- e a coerência entre cada KPI, o evento que o alimenta e a tela em que ele dispara. \\
\bottomrule
\end{longtable}
\end{center}

\begin{caixaGuardrail}{Dois freios, e uma coisa que não volta}
\textbf{Dois freios.} Cada ciclo tem teto próprio, e existe um orçamento global de reexecuções para a execução inteira (\cod{ORCAMENTO\_REEXECUCOES}, padrão 8). Estourado qualquer um dos dois, a esteira para de tentar, marca pendência humana e segue até o fim.

\textbf{O que não volta.} Achado cujo dono está na etapa de \emph{Produto} não reabre a \emph{spec}: reabri-la invalidaria tudo que o Design construiu sobre ela. Vira pendência nomeada no diagnóstico, para decisão humana.
\end{caixaGuardrail}

\section{Os portões, e o que fazer quando reprovam}
\label{sec:portoes}

Há dois tipos, e a diferença entre eles é a espinha dorsal do projeto.

\begin{center}
\begin{tikzpicture}[
  cx/.style={rectangle,rounded corners=4pt,text width=7.1cm,align=left,inner sep=10pt,font=\footnotesize}]
\node[cx,fill=unAzul!13,draw=unAzul,text=unNavy] (a)
 {{\sffamily\small\bfseries\color{unNavy}Determinístico --- é função Python}\\[5pt]
  Não gasta \emph{token}, não erra, não muda de opinião entre duas execuções e não pode ser convencido. Compara literalmente contra um documento: a lista oficial da BNCC, o vocabulário fechado do manifesto, o contrato real da API de telemetria.\\[4pt]
  \textbf{Toda regra que \emph{pode} ser mecânica foi tirada dos avaliadores e trazida para cá.}};
\node[cx,fill=unLavanda!13,draw=unLavanda,text=unNavy,right=0.45cm of a] (b)
 {{\sffamily\small\bfseries\color{unNavy}Avaliador --- é agente com rubrica}\\[5pt]
  Julga o que não se pode conferir mecanicamente: se há segredo, se a escolha importa, se a linguagem cabe na idade. Avaliadores baseados em modelos de linguagem têm vieses documentados de posição, verbosidade e leniência~\cite{zheng2023}.\\[4pt]
  \textbf{Por isso a rubrica sempre traz teste prático e contraexemplo}~\cite{jonsson2007} --- é o que sustenta concordância entre avaliações.};
\end{tikzpicture}
\end{center}

\begin{center}
\begin{longtable}{p{4.3cm}p{1.9cm}p{8.4cm}}
\toprule
\sffamily\bfseries\footnotesize\color{unNavy}Portão & \sffamily\bfseries\footnotesize\color{unNavy}Tipo & \sffamily\bfseries\footnotesize\color{unNavy}Reprova quando \\
\midrule
\endfirsthead
\toprule
\sffamily\bfseries\footnotesize\color{unNavy}Portão & \sffamily\bfseries\footnotesize\color{unNavy}Tipo & \sffamily\bfseries\footnotesize\color{unNavy}Reprova quando \\
\midrule
\endhead
Micro-gate (02) & rubrica & O briefing ou o núcleo entrega a virada; a escolha não muda nada; a missão repete uma já gerada. \\
Cobertura conceitual (01c) & \textbf{código} & O conceito limiar não vive no Decidir nem no ``ahá''; conceito do material sem destino; anatomia do ``ahá'' incompleta; trecho citado que não existe no material. \\
Profundidade (02b e 03b) & rubrica \newline + \textbf{código} & Nota abaixo de 3 em conceito vivido, aposta comprometida ou ápice do ``ahá''. O juiz dá as notas; \textbf{quem conta o veredito é código}, não a prosa do juiz. \\
Contrato de plataforma (04b) & \textbf{código} & Canal de sala por WebSocket, estado no navegador, SAEB em ano não avaliado, número sem origem declarada. \\
Gate estrutural (05) & rubrica & Falta mecanismo, há caminho sem tratamento, o ambiente está fora da lista. \\
Gate final (11) & rubrica & Os três pareceres avaliaram versões diferentes; correção que ficou só no relatório. \\
Currículo BNCC (14) & \textbf{código} & Código inexistente, de outro ano, ou sem letramento digital. \\
Manifesto (15 e 33) & \textbf{código} & Campo do bloco 0 ausente ou fora do vocabulário fechado. \\
Matriz de evidência (15b) & \textbf{código} & Habilidade sem momento em tela, sem medida ligada, ou com grau ``pontuado'' sem dois sinais nomeados e gabarito no servidor. \\
Perguntas de produto (15c--15e) & \textbf{código} \newline + rubrica & Pergunta sem decisão, sem corte ou sem ação declarada na tabela de eventos; e o que regex não vê --- se a pergunta muda mesmo uma decisão nossa, e se é sobre o produto e não sobre a criança. \\
\emph{Lint} do protótipo (29a) & \textbf{código} & O que o QA não enxerga por ler só o texto: o arquivo HTML em si. \\
QA de design (30) & rubrica & Construção infiel à \emph{spec}; diagrama ausente ou inválido; visão do professor faltando. \\
Telemetria (36) & \textbf{código} & Evento fora do contrato real da API; medida sem tabela de leitura; WebSocket. \\
Rastro de lacunas (38) & \textbf{código} & Marcadores de lacuna abertos nos documentos, por dono. \\
Banca (39--44) & rubrica & Qualquer critério bloqueante das cinco dimensões. \\
Coerência dos KPIs (48a) & \textbf{código} & KPI sem pergunta, sem decisão, sem corte, ou cujo evento não dispara na tela que ele nomeia. \\
Pendências e consistência do PRD (48) & \textbf{código} & Marcadores deixados por falta de informação; e o documento se contradizendo --- dois caminhos para o mesmo módulo, prefixo de tabela fora do padrão, requisito P0 sem teste nem tarefa, evento sem tabela leitora, tarefa que cobre tudo, identificador inválido. Reprovado, devolve ao redator (48b). \\
\bottomrule
\end{longtable}
\end{center}

\section{O que impede o sistema de alucinar}
\label{sec:anti}

Decompor em 72 passos resolve densidade de instrução e cria um problema novo: contrato de interface~\cite{cemri2025}. Quatro regras de engenharia de contexto sustentam as fronteiras.

\begin{enumerate}
  \item \textbf{Isolamento total de contexto.} Todos os agentes, com a única exceção do 01, rodam com \cod{include\_\allowbreak contents="none"}: não herdam o histórico da conversa. Recebem exclusivamente o que foi interpolado no próprio \emph{prompt}, a partir do estado da sessão. Sem isso, a poluição de contexto se acumula e cada agente passa a ``lembrar'' de versões antigas dos documentos.
  \item \textbf{Interpolação com degradação declarada.} A variável \cod{\{chave\}} falha imediatamente se a chave não existir --- é o comportamento correto quando a dependência é obrigatória. A variante \cod{\{chave?\}} degrada com elegância, e é o que permite rodar a etapa de Dev isolada, sobre um estado que não tem Design.
  \item \textbf{Contrato de envelope JSON.} Todo relatório de portão encerra com o mesmo envelope, e é ele --- não a prosa do relatório --- que o orquestrador lê para decidir o que fazer:
\begin{lstlisting}
{
  "status": "aprovado | apto_com_ressalvas | bloqueado",
  "veredito": "LIBERADA PARA DESIGN/DEV",
  "criterios_reprovados": [],
  "bloqueios": []
}
\end{lstlisting}
  \item \textbf{Decisões de produto protegidas.} Decisão estratégica que depende de premissa de negócio ausente no briefing é marcada como \cod{[DECISÃO PENDENTE --- responsável de produto]}. Os agentes são \textbf{proibidos de inventá-la}, e o diagnóstico a endereça a uma pessoa. É o oposto do comportamento padrão de um modelo de linguagem, que é preencher a lacuna com o plausível.
\end{enumerate}

\section{As integrações, e por que nenhuma é caminho crítico}

Dois servidores MCP participam da etapa de Design, com papéis distintos. Nenhum dos dois é obrigatório: fora do ar, cada um registra a própria indisponibilidade no artefato e a esteira entrega a missão assim mesmo --- porque o que vale é sempre o texto versionado.

\begin{center}
\begin{tikzpicture}[
  cx/.style={rectangle,rounded corners=4pt,text width=7.1cm,align=left,inner sep=10pt,font=\footnotesize}]
\node[cx,fill=unPapel,draw=unLavanda,text=unNavy] (a)
 {{\sffamily\small\bfseries\color{unNavy}Stitch --- as telas}\\[5pt]
  O modelo \textbf{usa ferramentas}: é ele que decide quais telas gerar. Sem chave de acesso, a esteira segue e registra \cod{STITCH: indisponível}; a especificação de telas em Markdown continua sendo a fonte.};
\node[cx,fill=unPapel,draw=unLavanda,text=unNavy,right=0.45cm of a] (b)
 {{\sffamily\small\bfseries\color{unNavy}Miro --- o espelho dos diagramas}\\[5pt]
  Aqui o modelo \textbf{não usa ferramenta nenhuma}. Um passo determinístico publica os mesmos blocos Mermaid como board navegável, fora do turno do agente. O board é \textbf{espelho, nunca fonte}: editá-lo não altera a missão, e a ausência dele não é achado de QA.};
\end{tikzpicture}
\end{center}

A razão de o Miro ser publicado por código, e não pelo agente, é concreta e custou uma execução: deixar um modelo construir o board item a item foi o que estourou o teto de contexto quando os diagramas viviam no FigJam. O Figma, aliás, saiu da esteira pelo mesmo tipo de razão --- não há cliente oficial que o sirva neste ambiente, e \emph{board} que ninguém consegue gerar não é entregável.

\parteUnooria{IV}{Por que funciona}{Uma rubrica que não diz de onde tira os números acaba defendida por antiguidade.}

\section{Descoberta ativa, e não instrucionismo}

A esteira foi construída sobre a aprendizagem pela ação de Dewey~\cite{dewey1938}, a aprendizagem pela descoberta de Bruner~\cite{bruner1961} e o construcionismo de Papert~\cite{papert1980}. Na prática, isso se traduz numa inversão de ordem que toda missão precisa respeitar --- e que o micro-gate (agente 02) verifica antes de qualquer outra coisa.

\begin{center}
\begin{tikzpicture}[
  et/.style={rectangle,rounded corners=2.5pt,text width=2.55cm,align=center,
    font=\sffamily\scriptsize,minimum height=0.95cm,inner sep=4pt},
  seta/.style={-{Stealth[length=2.0mm]},thick},
  rot/.style={font=\sffamily\scriptsize\bfseries}]
% instrucionismo
\node[et,fill=sinalVermelho!8,draw=sinalVermelho,text=unNavy] (i1) {a regra é enunciada};
\node[et,fill=sinalVermelho!8,draw=sinalVermelho,text=unNavy,right=0.35cm of i1] (i2) {o estudante lê};
\node[et,fill=sinalVermelho!8,draw=sinalVermelho,text=unNavy,right=0.35cm of i2] (i3) {responde a um quiz};
\node[et,fill=sinalVermelho!8,draw=sinalVermelho,text=unNavy,right=0.35cm of i3] (i4) {esquece};
\draw[seta,draw=sinalVermelho] (i1)--(i2); \draw[seta,draw=sinalVermelho] (i2)--(i3); \draw[seta,draw=sinalVermelho] (i3)--(i4);
\node[rot,text=sinalVermelho,anchor=south west] at (i1.north west) {INSTRUCIONISMO};
% descoberta
\node[et,fill=unPapel,draw=unNavy,text=unNavy,below=1.35cm of i1] (d1) {age dentro da simulação};
\node[et,fill=unPapel,draw=unNavy,text=unNavy,right=0.35cm of d1] (d2) {vê a consequência da escolha};
\node[et,fill=unNavy,draw=unNavy,text=white,right=0.35cm of d2] (d3) {chega ao ``ahá''};
\node[et,fill=unPapel,draw=unNavy,text=unNavy,right=0.35cm of d3] (d4) {a professora nomeia o conceito};
\draw[seta,draw=unLavanda] (d1)--(d2); \draw[seta,draw=unLavanda] (d2)--(d3); \draw[seta,draw=unLavanda] (d3)--(d4);
\node[rot,text=unNavy,anchor=south west] at (d1.north west) {UNOORIA};
\end{tikzpicture}
\captionof{figure}{A mesma aula, nas duas ordens. Na segunda, o conceito é nomeado \emph{depois} da experiência --- e por uma pessoa, não por uma tela.}
\end{center}

O detalhe que se perde com facilidade é o último quadro. A missão \textbf{não fecha o conceito sozinha}: ela produz a experiência e entrega à professora os dados de comparação para que a formalização aconteça na roda de conversa. Missão que explica tudo na tela final devolve o estudante ao instrucionismo, com uma camada de jogo por cima.

\section{A profundidade da missão: a anatomia do ``ahá''}
\label{sec:profundidade}

Em 24 de setembro de 2026 o time pedagógico devolveu as missões geradas com uma crítica que nenhum gate tinha detectado: elas eram \textbf{rasas}. Não trabalhavam os conceitos do material do bimestre, e o ``ahá'' não tinha ápice. Pior: o responsável de produto vinha liberando essas missões \emph{confiando no semáforo} --- que estava verde, porque media outra coisa.

\begin{caixaGuardrail}{A causa, conferida no código}
A esteira tinha 73 critérios sobre estrutura, plataforma, currículo e carga --- e \textbf{um} sobre o ``ahá'': ``dá para resumir numa frase?''. Qualquer missão rasa passa nesse teste.

E o núcleo narrativo \emph{recebia} o material do bimestre, mas nenhuma regra pedia que ele o \emph{usasse}. Duas omissões que ninguém tinha notado porque o semáforo verde respondia a uma pergunta diferente da que o time pedagógico fazia.
\end{caixaGuardrail}

A resposta foi uma camada inteira, que hoje é a parte mais nova do sistema: três passos antes do núcleo, duas tabelas dentro dele e um juiz próprio no ciclo.

\subsection{Ler o material antes de inventar}

A ordem importa, e antes ela estava errada: o núcleo inventava o mecanismo e depois procurava no material algo que o justificasse. Hoje há dois passos antes dele.

\begin{center}
\begin{tikzpicture}[
  cx/.style={rectangle,rounded corners=3pt,text width=4.42cm,align=left,inner sep=9pt,
    font=\footnotesize,text=unNavy},
  seta/.style={-{Stealth[length=2.2mm]},draw=unLavanda,thick}]
\node[cx,fill=unAzul!12,draw=unAzul] (a)
 {{\sffamily\small\bfseries 00c · Mapa conceitual}\\[4pt]
  Lê o material do bimestre e extrai: o \textbf{conceito limiar}, os conceitos de apoio \emph{com o trecho citado}, as \textbf{crenças de entrada} que o bimestre derruba e o vocabulário do Explicar.};
\node[cx,fill=unLavanda!15,draw=unLavanda,right=0.35cm of a] (b)
 {{\sffamily\small\bfseries 00d · Ideação}\\[4pt]
  \textbf{Três} mecanismos distintos para aquele conceito limiar, cada um com a anatomia do ``ahá'' escrita e pontuado na rubrica. O mecanismo que veio no briefing entra como um dos candidatos.};
\node[cx,fill=unNavy,text=white,right=0.35cm of b] (c)
 {{\sffamily\small\bfseries 01 · Núcleo}\\[4pt]
  Desenvolve \textbf{o escolhido}, e escreve duas tabelas novas: os conceitos vividos em tela e a anatomia do ``ahá'' em seis passos.};
\draw[seta] (a)--(b); \draw[seta] (b)--(c);
\end{tikzpicture}
\end{center}

Uma verificação determinística (01c) confere o núcleo contra o mapa: o conceito limiar vive no Decidir ou no ``ahá''; todo conceito do material tem destino; a anatomia tem os seis passos; e \textbf{o trecho citado pelo mapa existe mesmo no material} --- porque citação inventada é o modo de falha mais fácil de um modelo que precisa fundamentar.

\subsection{Os seis passos do ``ahá''}

Um ``ahá'' genuíno não é informação nova. É uma crença que o estudante carregava, pôs à prova com uma escolha dele, e viu cair pela consequência dessa escolha.

\begin{center}
\begin{longtable}{p{0.55cm}p{3.0cm}p{5.6cm}p{5.4cm}}
\toprule
\multicolumn{1}{l}{\sffamily\bfseries\footnotesize\color{unLavanda}\#} & \multicolumn{1}{l}{\sffamily\bfseries\footnotesize\color{unNavy}Passo} & \multicolumn{1}{l}{\sffamily\bfseries\footnotesize\color{unNavy}O teste} & \multicolumn{1}{l}{\sffamily\bfseries\footnotesize\color{unNavy}A falha vista em execução} \\
\midrule
\endfirsthead
\toprule
\multicolumn{1}{l}{\sffamily\bfseries\footnotesize\color{unLavanda}\#} & \multicolumn{1}{l}{\sffamily\bfseries\footnotesize\color{unNavy}Passo} & \multicolumn{1}{l}{\sffamily\bfseries\footnotesize\color{unNavy}O teste} & \multicolumn{1}{l}{\sffamily\bfseries\footnotesize\color{unNavy}A falha vista em execução} \\
\midrule
\endhead
\sffamily\bfseries\color{unLavanda}1 & \sffamily\bfseries\footnotesize\color{unNavy}Crença de entrada & \footnotesize Está nomeada, e alguém daquela idade de fato acredita nisso? & \footnotesize Nenhuma crença declarada: a missão supõe que o estudante já sabe, e o ``ahá'' só confirma. \\
\addlinespace
\sffamily\bfseries\color{unLavanda}2 & \sffamily\bfseries\footnotesize\color{unNavy}Aposta comprometida & \footnotesize A escolha é dele, é sobre o conceito limiar, e \textbf{trava} antes da revelação. \emph{Teste de remoção:} tirando o Decidir, o ``ahá'' muda? Se não muda, é informação, não descoberta. & \footnotesize ``A Foto do Mural'': nenhuma aposta, só observar. ``O Rastro do Boato'': a escolha era sobre confiabilidade da fonte, e a descoberta, sobre distribuição. \\
\addlinespace
\sffamily\bfseries\color{unLavanda}3 & \sffamily\bfseries\footnotesize\color{unNavy}Consequência sentida & \footnotesize Dois estudantes que escolheram diferente veem coisas diferentes? A consequência é mostrada pelo ambiente, e não dita pelo professor ou por um painel? & \footnotesize Panorama igual para qualquer escolha; consequência que só existia na fala oral do professor; painel com ``94\% e 6\%''. \\
\addlinespace
\sffamily\bfseries\color{unLavanda}4 & \sffamily\bfseries\footnotesize\color{unNavy}Virada & \footnotesize \emph{Teste da crença:} leia o passo 1 e depois a virada. Ela \textbf{desmente} a crença, ou só acrescenta informação? A reação esperada é ``caramba, eu achava que\ldots'', não ``ah, faz sentido''. & \footnotesize Resultado que confirma o que o Desafiar já tinha dito --- \emph{spoiler} estrutural que passou pelo micro-gate. \\
\addlinespace
\sffamily\bfseries\color{unLavanda}5 & \sffamily\bfseries\footnotesize\color{unNavy}Nomeação & \footnotesize O termo do material aparece? \textbf{Só depois} da virada? Usa o que o próprio estudante fez? Compara ao menos duas escolhas da turma? & \footnotesize Card de texto genérico; conceito explicado antes da vivência; termo técnico sem âncora no que aconteceu. \\
\addlinespace
\sffamily\bfseries\color{unLavanda}6 & \sffamily\bfseries\footnotesize\color{unNavy}Transferência & \footnotesize Um estudante que não fez a missão responde só pela pergunta? A resposta exige o mecanismo (``isso acontece porque\ldots'')? & \footnotesize Refletir binário com resposta óbvia; distrator implausível; ``pode ter \emph{fake news}'' aceito como resposta. \\
\bottomrule
\end{longtable}
\end{center}

Duas regras de contorno, que evitam que a rubrica produza o dano oposto:

\begin{itemize}
  \item \textbf{Discriminação, não cinismo.} Em missão sobre desinformação, a lição não pode virar ``desconfie de tudo'': há sempre ao menos um item confiável que o estudante deve \emph{manter}, e a medida separa o acerto no confiável do acerto no manipulado. Jogos de letramento midiático às vezes só aumentam o ceticismo geral, e isso não é aprendizado.
  \item \textbf{Anos iniciais são diferentes.} Do 2.\textsuperscript{o} ao 5.\textsuperscript{o} ano, ``viver primeiro e explicar depois'' não tem a mesma evidência que nos anos finais, e aprender com a própria previsão errada depende de controle inibitório ainda em formação. Por isso a discrepância ``você escolheu X, aconteceu Y'' aparece \textbf{explícita na tela}, e um andaime breve antes do Decidir é permitido --- desde que não entregue a virada.
\end{itemize}

\subsection{O juiz, e por que ele não decide sozinho}

Um agente próprio dá nota de 1 a 4 em seis dimensões --- conceito do material vivido, aposta comprometida, ápice do ``ahá'', explicar depois de viver, transferência e demanda à altura do ano --- sempre com \textbf{trecho citado} e instrução de reescrita. Nota abaixo de 3 nas três primeiras reprova.

\begin{destaque}
\textbf{Quem conta o veredito é código, não o juiz.} Um passo determinístico relê as notas que o juiz deu e recalcula o resultado, ignorando o que ele escreveu na conclusão. É a mesma desconfiança que sustenta o resto da esteira: avaliadores baseados em modelos tendem à leniência~\cite{zheng2023}, e um juiz que dá nota 2 e conclui ``aprovado com ressalvas'' é um caso observado, não hipotético.
\end{destaque}

\subsection{O que esta rubrica não faz}

\begin{itemize}
  \item \textbf{Não cria número novo.} Os tetos de opções, duração, gestos e abstração continuam vindo da fonte única (Seção~\ref{sec:regua}).
  \item \textbf{Profundidade não é carga.} Nota 4 nunca exige mais telas, mais opções ou mais minutos --- a tela do pico continua sendo a mais enxuta da missão. Confundir as duas coisas produziria missões densas e cansativas em nome da profundidade.
  \item \textbf{Não decide pelo time.} Quando o mecanismo vem do responsável de produto, o juiz pontua e aponta; \textbf{trocar de mecanismo é decisão humana}.
\end{itemize}

\begin{caixaConceito}{O que ainda falta: a calibragem}
Os cortes --- reprovar abaixo de 3 nas três primeiras dimensões --- são \Convencao, e estão declarados assim. Falta calibrá-los contra três a cinco pares ``versão gerada $\rightarrow$ versão revisada pelo time pedagógico'', que é o material que diria onde a régua deve de fato ficar.

Enquanto isso não acontece, \cod{PROFUNDIDADE\_ESTRITA=false} rebaixa o bloqueio a ressalva sem mexer nas notas nem nas rodadas. É o interruptor certo para usar durante a calibragem --- e o errado para deixar ligado depois dela.
\end{caixaConceito}

\section{Carga cognitiva, e o que a esteira faz com ela}

A teoria da carga cognitiva~\cite{sweller1988,sweller2019} separa três demandas que competem pela mesma memória de trabalho:

\begin{center}
\begin{tikzpicture}[
  cg/.style={rectangle,rounded corners=4pt,text width=4.45cm,align=left,inner sep=9pt,font=\footnotesize}]
\node[cg,fill=unNavy,text=white] (a)
 {{\sffamily\small\bfseries Intrínseca}\\[4pt] A complexidade real do que se quer ensinar --- entender que uma recomendação é inferida do próprio comportamento. \\[4pt] \textbf{Não se reduz sem reduzir o aprendizado.}};
\node[cg,fill=sinalVermelho!8,draw=sinalVermelho,text=unNavy,right=0.35cm of a] (b)
 {{\sffamily\small\bfseries Estranha}\\[4pt] O esforço desperdiçado com interface confusa, texto longo, botão escondido, instrução ambígua. \\[4pt] \textbf{É o alvo de metade dos critérios de design da esteira.}};
\node[cg,fill=unAzul!13,draw=unAzul,text=unNavy,right=0.35cm of b] (c)
 {{\sffamily\small\bfseries Germânica}\\[4pt] O esforço que constrói esquema duradouro: comparar, explicar, aplicar a um caso novo. \\[4pt] \textbf{É o que a etapa Refletir existe para provocar.}};
\end{tikzpicture}
\end{center}

Disso sai uma régua concreta que os agentes de design aplicam: a memória de trabalho lida com \textbf{dois a três elementos novos e interagentes} por vez~\cite{sweller2019} --- e elemento já automatizado, como o ícone de coração ou o botão de play, \emph{não conta}. Derivando por faixa: cerca de 1 a 2 gestos novos por ato nos anos iniciais, 2 a 3 no 3.\textsuperscript{o}--5.\textsuperscript{o}, 3 a 4 nos anos finais e no Médio. A derivação é explicitamente marcada como extrapolação --- ninguém mediu ``elementos de interface''.

\section{A régua por idade}
\label{sec:regua}

Os tetos que a esteira aplica vivem num arquivo só: \cod{prompts/references/literatura-cognitiva.md}. \textbf{Nenhum agente carrega régua própria}; quando um critério precisa de um número, ele vem de lá.

\begin{caixaConceito}{Por que a fonte tem de ser única}
Porque números soltos em \emph{prompts} divergem sem que ninguém perceba. Foi o que aconteceu quando um avaliador novo chegou com ``3 a 4 elementos'' e a rubrica de outro agente tinha uma tabela por faixa dizendo outra coisa. Os dois estavam em produção, os dois julgavam a mesma missão, e nenhum dos dois estava errado segundo o próprio \emph{prompt}.
\end{caixaConceito}

Cada número no arquivo carrega uma marca: \textbf{[DADO DIRETO]} quando a fonte mediu exatamente aquilo, naquela faixa; \textbf{[EXTRAPOLAÇÃO]} quando a rubrica deriva um limite prático de um princípio geral; \textbf{[CONVENÇÃO]} quando a equipe escolheu por consistência de produto. Extrapolação justifica ressalva com mais frequência que bloqueio.

\subsection{Quantas opções o estudante compara de uma vez}

\begin{center}
\begin{tabularx}{\textwidth}{>{\sffamily\bfseries\footnotesize\color{unNavy}}p{3.9cm}>{\centering\arraybackslash\bfseries\footnotesize}p{1.5cm}X}
\toprule
\sffamily\bfseries\footnotesize\color{unNavy}Faixa & \sffamily\bfseries\footnotesize\color{unNavy}Teto & \sffamily\bfseries\footnotesize\color{unNavy}De onde vem \\
\midrule
1.\textsuperscript{o}--2.\textsuperscript{o} EF (6--7 anos) & 2 & \emph{Span} complexo medido: a criança sustenta entre 1,6 e 2,9 itens \emph{enquanto processa}~\cite{gathercole2004,cowan2016}. \\
3.\textsuperscript{o}--5.\textsuperscript{o} EF (8--10 anos) & 3 & \emph{Span} complexo de 1,9 a 3,7 itens. \\
6.\textsuperscript{o}--9.\textsuperscript{o} EF (11--14 anos) & 3--4 & Memória visual de trabalho $k = 2{,}4$ aos 11--13 anos. \\
Ensino Médio (15--17 anos) & 4 & $k = 2{,}72$ aos 16 anos, ainda \textbf{abaixo} do adulto jovem (3,16)~\cite{isbell2015}. \\
\bottomrule
\end{tabularx}
\end{center}

\begin{destaque}
\textbf{Opção visível não é opção lembrada.} Se todas as alternativas ficam na tela enquanto o estudante decide, o gargalo deixa de ser memória e passa a ser comparação atencional --- e o teto sobe em \textbf{+1}. Se ele precisa lembrar de opções que já saíram da tela, vale o teto da tabela. Essa distinção era a maior imprecisão da versão anterior da régua.
\end{destaque}

\subsection{Quanto tempo dura cada ato}

\begin{center}
\begin{minipage}{\textwidth}\centering
\begin{tikzpicture}[xscale=0.62]
\draw[draw=unLinha] (0,0) -- (0,-4.6);
\foreach \x in {0,4,8,12,16} {
  \draw[draw=unLinha] (\x,0.25) -- (\x,-4.6);
  \node[font=\sffamily\scriptsize,text=unLavanda,above] at (\x,0.25) {\x\ min};}
\foreach \y/\lab/\a/\b in {
  -0.9/{1.\textsuperscript{o}--2.\textsuperscript{o} EF}/6/8,
  -1.9/{3.\textsuperscript{o}--5.\textsuperscript{o} EF}/8/10,
  -2.9/{6.\textsuperscript{o}--9.\textsuperscript{o} EF}/10/12,
  -3.9/{Ensino Médio}/12/15} {
  \fill[unLavanda!35,rounded corners=1.5pt] (0,\y-0.22) rectangle (\a,\y+0.22);
  \fill[unNavy,rounded corners=1.5pt] (\a,\y-0.22) rectangle (\b,\y+0.22);
  \node[font=\sffamily\scriptsize,text=unNavy,anchor=east] at (-0.25,\y) {\lab};
  \node[font=\sffamily\scriptsize\bfseries,text=white,anchor=east] at (\b-0.15,\y) {\a--\b};}
\end{tikzpicture}
\captionof{figure}{Teto de duração por ato. A escala é deliberadamente achatada: em 52 salas observadas, a atenção cai monotonicamente com a duração e o efeito de série foi equívoco --- \textbf{a duração importa mais que a idade}~\cite{godwin2016}.}
\end{minipage}
\end{center}

O teto que mais importa, porém, é o \textbf{agregado}: no máximo \textbf{25 minutos de tempo-aluno} numa aula de 50. Ele é novo, e existe por aritmética: três atos no teto antigo somavam 45 minutos e não sobrava aula. Some-se que apenas 64\% do tempo de aula é atividade acadêmica no Brasil~\cite{bruns2018} --- descontando instrução inicial e distribuição de dispositivo, a janela real fica entre 22 e 27 minutos.

E há um fator que pesa mais que o relógio: \textbf{o formato}. Atividade individual ou em pequeno grupo sustenta atenção bem melhor que turma inteira (razão de chances até 1,62), e o critério de duração passou a considerar isso.

\subsection{O que se pode pedir que o estudante entenda}

Missão de letramento digital quase sempre pede raciocínio sobre um sistema invisível --- um algoritmo, uma inferência, um viés. Não há uma idade única; há três portas.

\begin{center}
\begin{tikzpicture}[
  porta/.style={rectangle,rounded corners=3pt,text width=3.28cm,align=center,
    font=\sffamily\scriptsize,minimum height=1.6cm,inner sep=6pt},
  seta/.style={-{Stealth[length=2.2mm]},draw=unLavanda,thick}]
\node[porta,fill=unAzul!13,draw=unAzul,text=unNavy] (a)
  {\textbf{Descrever}\\[3pt] ``o app me mostra mais vídeo de gato''\\[3pt] {\bfseries a partir de 8 anos}};
\node[porta,fill=unLavanda!20,draw=unLavanda,text=unNavy,right=0.4cm of a] (b)
  {\textbf{Explicar o mecanismo}\\[3pt] dados $\rightarrow$ inferência $\rightarrow$ recomendação\\[3pt] {\bfseries 10--11 anos}};
\node[porta,fill=unNavy,draw=unNavy,text=white,right=0.4cm of b] (c)
  {\textbf{Criticar a intenção}\\[3pt] e o viés, \emph{com apoio}\\[3pt] {\bfseries 12--13 anos}};
\node[porta,fill=unPapel,draw=unNavy,text=unNavy,right=0.4cm of c] (d)
  {\textbf{Criticar sozinho}\\[3pt] sem andaime\\[3pt] {\bfseries 14+ \ \itshape(extrapolação)}};
\draw[seta] (a)--(b); \draw[seta] (b)--(c); \draw[seta] (c)--(d);
\end{tikzpicture}
\captionof{figure}{As três portas do raciocínio sobre sistemas invisíveis~\cite{starks2024,rozendaal2011}. É o núcleo do letramento digital, e até a revisão de setembro não tinha régua nenhuma.}
\end{center}

\subsection{O piso da escola pública brasileira}

Nenhuma missão é avaliada contra o laboratório ideal. O cenário de piso é declarado, com fonte --- e mudou na última revisão, porque a suposição da equipe estava errada em duas direções ao mesmo tempo.

\begin{center}
\begin{tabularx}{\textwidth}{>{\sffamily\bfseries\footnotesize\color{unNavy}}p{4.7cm}>{\bfseries\footnotesize\centering\arraybackslash}p{3.0cm}X}
\toprule
\sffamily\bfseries\footnotesize\color{unNavy}Parâmetro & \sffamily\bfseries\footnotesize\color{unNavy}Piso realista & \sffamily\bfseries\footnotesize\color{unNavy}Observação \\
\midrule
Alunos por turma, anos iniciais & 25 & A média oficial é 22,3 --- afirmar ``30 a 40'' \textbf{infla} a dificuldade. \\
Alunos por turma, finais e Médio & 32 & 26,6 no Censo; 32,5 na turma de referência do TALIS. \\
Dispositivo & nenhum fornecido pela escola & Só 47\% das escolas disponibilizam computador ao aluno. O piso não é ``compartilhar'', é \textbf{não ter}. \\
Internet & sinal ausente em sala & 96\% das escolas têm internet, mas 39\% das estaduais relatam que o sinal não chega à sala. \\
Leitura, 5.\textsuperscript{o} ano & 39\% abaixo do adequado & Ou seja, \textbf{60,7\% estão adequados}: tratar os anos iniciais como majoritariamente defasados contraria o Saeb. \\
Leitura, 9.\textsuperscript{o} ano e Médio & 62\% e 65\% abaixo & Aqui a defasagem se confirma. \\
Professor & sem formação recente em tecnologia & 46\% não tiveram; apenas cerca de um terço decide sobre o uso de recurso digital. \\
\bottomrule
\end{tabularx}
\end{center}

\begin{caixaConceito}{O teste de piso, em uma frase}
Uma missão precisa funcionar numa turma de 25 a 32 estudantes, com dispositivo que a escola não forneceu, com sinal que pode não chegar à sala, conduzida por uma professora que não teve formação recente em tecnologia --- e ainda assim caber em 25 minutos de tempo-aluno. Tudo que a esteira exige de tela limpa, texto curto e ausência de download pesado deriva daqui, e não de gosto estético.
\end{caixaConceito}

\section{O registro negativo: o que não sustenta a régua}

Esta é a seção mais incomum da fonte única, e a mais importante. Sem ela, uma fundamentação refutada tende a reaparecer na revisão seguinte, porque \textbf{o argumento de autoridade sobrevive ao artigo que o derrubou}.

\begin{center}
\begin{longtable}{p{5.1cm}p{10.2cm}}
\toprule
\sffamily\bfseries\footnotesize\color{unNavy}Não use & \sffamily\bfseries\footnotesize\color{unNavy}Por quê \\
\midrule
\endfirsthead
\toprule
\sffamily\bfseries\footnotesize\color{unNavy}Não use & \sffamily\bfseries\footnotesize\color{unNavy}Por quê \\
\midrule
\endhead
\textbf{\emph{Choice overload}} como base do teto de opções &
Meta-análise de 63 condições ($N = 5.036$): efeito médio praticamente nulo. Em adultos, o ponto de inflexão fica entre 12 e 24 opções --- não 4. Em crianças, um estudo de 2023 não encontrou o efeito \textbf{nem com 53 opções}, e uma revisão de 2024 indica que crianças sofrem \emph{menos} que adultos com arranjos grandes. A lógica ``mais nova, menos opções por \emph{overload}'' está invertida~\cite{castelo2023}. \\
\addlinespace
\textbf{\emph{Decision fatigue}} como base do teto de volume &
Dois estudos multilaboratório de replicação registrada com efeito nulo ($d = 0{,}04$); o caso dos ``juízes com fome'' se explica por ordem de apresentação~\cite{hagger2016}. O sinal real é monotonia, não esgotamento. \\
\addlinespace
\textbf{``Idade da criança = minutos de atenção''} &
Não tem origem primária publicada. Circula em três variantes (1$\times$, 2--3$\times$, 5$\times$ a idade) que diferem por um fator de cinco. O análogo adulto --- ``10 a 15 minutos de atenção em aula'' --- foi formalmente desmontado. \\
\addlinespace
\textbf{Piaget} para vetar abstração por idade &
Idades fixas e universais são justamente a parte refutada. Use ciclos de prioridade cognitiva~\cite{demetriou2022} e dependência de suporte~\cite{fischer1980}. \\
\addlinespace
\textbf{Bloom ou DOK} como régua de adequação etária &
São instrumentos de alinhamento entre currículo e avaliação, ortogonais à série --- Webb prevê DOK 3 e 4 em todos os anos. SOLO~\cite{biggs1982} é a única das três defensável para calibrar a resposta esperada. \\
\addlinespace
\textbf{Classificar a tarefa pelo verbo do enunciado} &
Refutado sobre 940 itens: o verbo não é \emph{proxy} confiável do processo cognitivo. O teste de alinhamento de verbo continua útil como sinal, \textbf{nunca} como veredito automático. \\
\bottomrule
\end{longtable}
\end{center}

\subsection{O que mudou na revisão de setembro de 2026}

A revisão não ajustou números por gosto: trocou a fundamentação de quatro critérios, porque a pesquisa mostrou que as fontes citadas não sustentavam o que a rubrica afirmava.

\begin{center}
\begin{longtable}{p{3.9cm}p{3.0cm}p{8.0cm}}
\toprule
\sffamily\bfseries\footnotesize\color{unNavy}Mudança & \sffamily\bfseries\footnotesize\color{unNavy}Antes & \sffamily\bfseries\footnotesize\color{unNavy}Por quê \\
\midrule
\endfirsthead
\toprule
\sffamily\bfseries\footnotesize\color{unNavy}Mudança & \sffamily\bfseries\footnotesize\color{unNavy}Antes & \sffamily\bfseries\footnotesize\color{unNavy}Por quê \\
\midrule
\endhead
Teto de opções no Médio & 4 a 5 & A capacidade visual de trabalho aos 16 anos é 2,72 itens, ainda abaixo do adulto. Tratar o Médio como capacidade adulta superestimava. \\
\addlinespace
Base do teto de opções & \emph{choice overload} & O efeito não replicou. O que sustenta os tetos é \emph{span} complexo de memória de trabalho --- outro construto. \\
\addlinespace
Volume de decisões & obrigatório, com teto & Virou desejável e marcado como convenção; o sinal real é monotonia. \\
\addlinespace
Tetos de duração & 12 / 15 / 20 min & Observação de 52 salas: queda monotônica com a duração, sem minuto mágico. A escala foi achatada. \\
\addlinespace
Teto agregado da missão & não existia & Três atos no teto antigo somavam mais que uma aula inteira. \\
\addlinespace
Formato da atividade & ignorado & Individual ou em pequeno grupo sustenta atenção bem melhor que turma inteira. \\
\addlinespace
Justificativa de abstração & Piaget & A parte citada era a refutada. Substituída por Demetriou e Fischer. \\
\addlinespace
Raciocínio sobre sistema invisível & não existia & É o núcleo do letramento digital e não tinha régua. Entrou com três portas. \\
\addlinespace
Cenário da escola pública & suposição & Substituído por dados de INEP, CETIC e Saeb. \\
\bottomrule
\end{longtable}
\end{center}

\section{O papel do professor}

A falha mais comum em EdTech é desenhar como se a professora não existisse --- ou pior, como se ela fosse um obstáculo a contornar. Aqui ela é o centro da experiência, e isso é auditado por critério próprio na avaliação pedagógica.

\begin{center}
\begin{tikzpicture}[
  mom/.style={rectangle,rounded corners=3pt,text width=3.24cm,align=left,
    font=\sffamily\scriptsize,inner sep=8pt,minimum height=2.5cm,text=unNavy},
  seta/.style={-{Stealth[length=2.2mm]},draw=unLavanda,thick}]
\node[mom,fill=unPapel,draw=unLavanda] (a)
  {\textbf{Antes}\\[3pt] Customiza tema, contexto e exemplos locais. A IA gera rascunhos; \textbf{ela cura, edita e aprova}. Nada entra na sala sem aprovação explícita.};
\node[mom,fill=unPapel,draw=unLavanda,right=0.3cm of a] (b)
  {\textbf{Durante}\\[3pt] Acompanha no painel quem avançou e quem travou, e \textbf{controla o ritmo} --- inclusive pausando a turma antes da revelação.};
\node[mom,fill=unNavy,text=white,right=0.3cm of b] (c)
  {\textbf{No pico}\\[3pt] Conduz a comparação entre colegas reais. O ``ahá'' emerge de diferenças genuínas na turma, não de uma tela de resposta.};
\node[mom,fill=unPapel,draw=unLavanda,right=0.3cm of c] (d)
  {\textbf{Depois}\\[3pt] Conduz a roda de conversa com o roteiro de perguntas abertas. \textbf{A missão digital não encerra a aula: ela acende o debate.}};
\draw[seta] (a)--(b); \draw[seta] (b)--(c); \draw[seta] (c)--(d);
\end{tikzpicture}
\captionof{figure}{Os quatro momentos em que a autonomia docente é exercida --- e auditada.}
\end{center}

Há um quinto ponto, menos nobre e igualmente decisivo: \textbf{ritmos diferentes}. Toda missão precisa prever o que fazem os estudantes que terminam cedo, sob pena de eles darem \emph{spoiler} para o resto da turma e destruírem a virada para todo mundo.

\parteUnooria{V}{O que protege}{Desenhar para crianças impõe uma responsabilidade que nenhum software comum tem. Aqui ela está escrita em código.}

\section{Alinhamento curricular e o primado do letramento digital}
\label{sec:bncc}

Em outubro de 2022, o Conselho Nacional de Educação homologou a \textbf{Resolução CNE/CEB n.\textsuperscript{o} 1/2022}, que instituiu as normas sobre Computação na Educação Básica como complemento oficial à BNCC~\cite{cne2022}. A norma organiza o ensino de computação em três eixos:

\begin{center}
\begin{tikzpicture}[
  eixo/.style={rectangle,rounded corners=4pt,text width=4.45cm,align=left,inner sep=9pt,font=\footnotesize}]
\node[eixo,fill=unPapel,draw=unLavanda,text=unNavy] (a)
 {{\sffamily\small\bfseries Pensamento computacional}\\[4pt] Resolver problemas por algoritmos, decomposição, abstração e reconhecimento de padrões.\\[5pt] {\color{sinalAmbar}\bfseries avalia e alerta}};
\node[eixo,fill=unPapel,draw=unLavanda,text=unNavy,right=0.35cm of a] (b)
 {{\sffamily\small\bfseries Mundo digital}\\[4pt] Compreender a infraestrutura: \emph{hardware}, \emph{software}, codificação de dados, inteligência artificial.\\[5pt] {\color{sinalAmbar}\bfseries avalia e alerta}};
\node[eixo,fill=unNavy,text=white,right=0.35cm of b] (c)
 {{\sffamily\small\bfseries Cultura digital}\\[4pt] Uso ético, crítico, responsável e reflexivo das tecnologias digitais na sociedade.\\[5pt] {\color{unAzul}\bfseries É A ÚNICA QUE BLOQUEIA}};
\end{tikzpicture}
\end{center}

\begin{caixaGuardrail}{A regra de ouro do gate curricular}
Na esteira Unooria, a \textbf{única} habilidade da BNCC com poder de bloquear a missão é o cumprimento \textbf{ativo} de letramento digital --- o eixo Cultura Digital.

Habilidades de componentes tradicionais (Língua Portuguesa, História, Matemática), socioemocionais ou do eixo Pensamento Computacional recebem nota, sugestão e alerta \statusAmbar, mas \textbf{nunca travam a esteira}. Sem letramento digital comprovado na prática, a missão recebe \statusVermelho.

Esta é uma \textbf{decisão de governança de produto}, não um achado de pesquisa --- e está declarada como tal na própria rubrica, para que ninguém a defenda como se fosse ciência.
\end{caixaGuardrail}

A decisão se apoia em três razões, e vale enunciá-las porque elas são contestáveis --- e devem ser contestadas com argumento, não com antiguidade:

\begin{enumerate}
  \item \textbf{Contra o tecnologismo vazio.} Muitas plataformas de programação infantil ensinam a arrastar blocos para mover um boneco e param aí. Criar um robô sem saber discernir uma notícia falsa, ou sem entender o valor dos próprios dados, não forma cidadão para este século~\cite{buckingham2007,unesco2018}.
  \item \textbf{Urgência de proteção.} Crianças brasileiras passam horas conectadas. A escola tem dever cívico e legal de prepará-las contra assédio digital, vazamento de imagem, golpe e manipulação algorítmica. Letramento digital é, antes de tudo, competência de autoproteção~\cite{livingstone2017}.
  \item \textbf{Vivência, não teoria.} O critério exige que o estudante \textbf{pratique} o letramento digital na interface. ``Ler uma dica de segurança'' reprova. Ele precisa identificar, configurar, reportar ou avaliar --- com o dedo na tela.
\end{enumerate}

\subsection{Como o verificador de BNCC funciona}

Modelos de linguagem inventam códigos curriculares com naturalidade e confiança. Por isso a esteira \textbf{não confia na memória do modelo}: o agente 14 é código Python puro, que consulta a lista oficial do MEC em \cod{prompts/references/bncc-computacao.json}.

\begin{center}
\begin{tikzpicture}[
  reg/.style={rectangle,rounded corners=3pt,draw=unAzul,fill=unAzul!10,text=unNavy,
    text width=4.42cm,align=left,font=\sffamily\scriptsize,inner sep=8pt,minimum height=1.9cm},
  seta/.style={-{Stealth[length=2.2mm]},draw=unLavanda,thick}]
\node[reg] (a) {\textbf{1. Existência real}\\[3pt] Todo código citado (ex.: \cod{EF06CO06}) tem de existir na lista oficial. Código inventado bloqueia.};
\node[reg,right=0.35cm of a] (b) {\textbf{2. Ano compatível}\\[3pt] O código tem de corresponder ao ano da missão. Missão de 3.\textsuperscript{o} ano não cita habilidade exclusiva de 9.\textsuperscript{o}.};
\node[reg,right=0.35cm of b] (c) {\textbf{3. Letramento digital presente}\\[3pt] No Fundamental, a missão sem ao menos uma habilidade de Cultura Digital é reprovada.};
\end{tikzpicture}
\captionof{figure}{As três regras do verificador determinístico de currículo. Nenhuma delas passa por um modelo de linguagem.}
\end{center}

A escolha de fazer isso em código, e não por avaliador, é deliberada: é regra mecânica, e avaliador baseado em modelo tem viés documentado de leniência~\cite{zheng2023}. Um avaliador aceitaria um código plausível; uma comparação literal, não.

\subsection{Quando nenhuma habilidade do ano serve}

Há um caso que a primeira edição deste manual não cobria, e que o time encontrou na prática: a missão é boa, o material é bom, e \textbf{nenhuma habilidade de Cultura Digital daquele ano descreve a ação que o estudante pratica}. Aconteceu com uma missão de 6.\textsuperscript{o} ano sobre filtragem algorítmica --- a Cultura Digital do 6.\textsuperscript{o} ano simplesmente não fala disso.

A decisão do time foi \textbf{bloquear e devolver}, não acomodar. A esteira grava um marcador de decisão pendente com dono nomeado, e a missão não segue.

\begin{caixaConceito}{Por que não afrouxar aqui}
Afrouxar teria duas saídas, e as duas são piores. Aceitar uma habilidade de outro eixo transformaria o único critério curricular bloqueante em decorativo. Aceitar uma habilidade de Cultura Digital que \emph{não} descreve a ação produziria um documento que afirma cumprir o currículo sem cumprir --- e alguém adiante confiaria nessa afirmação.

O que a esteira faz em vez disso é nomear a decisão: mudar a missão de ano, ancorar em outro componente, ou aceitar que a missão não é de letramento digital. Essas são escolhas de produto, e a máquina não tem autoridade sobre nenhuma delas.
\end{caixaConceito}

\subsection{O contrato de evidência: o que o número sobre uma criança pode afirmar}

Declarar a habilidade não basta; é preciso declarar \textbf{como ela é medida}. A \emph{spec} traz uma tabela de medidas que um portão determinístico lê, e cada medida tem um grau:

\begin{center}
\begin{tabularx}{\textwidth}{>{\ttfamily\footnotesize\color{unNavy}}p{2.2cm}>{\raggedright\arraybackslash\footnotesize}X}
\toprule
\multicolumn{1}{l}{\sffamily\bfseries\footnotesize\color{unNavy}Grau} & \sffamily\bfseries\footnotesize\color{unNavy}O que significa, e o que ele autoriza afirmar \\
\midrule
scored & Dois ou mais sinais independentes. Tem número, e o número entra nas médias. \\
indicative & Um sinal só. Vira texto, e \textbf{fica fora das médias} --- porque um sinal não sustenta uma nota. \\
zero & O estudante teve a chance e não fez. Entra na média. \\
none & \textbf{A oportunidade não chegou até ele.} Sai da conta. \\
\bottomrule
\end{tabularx}
\end{center}

\begin{caixaGuardrail}{A distinção entre \cod{zero} e \cod{none} é uma regra de proteção}
Tratar ``não teve chance'' como ``escolheu não fazer'' faz uma criança \textbf{parecer ter recusado o que nunca lhe foi oferecido}. E o número será lido por um adulto que decide sobre ela.

A esteira não confia na declaração: das 25 medidas marcadas como \cod{scored} nas primeiras execuções, \textbf{só 10 nomeavam os dois sinais} que o próprio grau exige. Hoje um portão determinístico rebaixa ou reprova cada uma que não nomeia --- o grau inflado era invisível justamente por parecer rigor.
\end{caixaGuardrail}

\section{Inteligência artificial com humano no controle}

A literatura sobre segurança em IA é clara quanto aos riscos de alucinação, viés e manipulação quando crianças interagem diretamente com modelos generativos abertos~\cite{bommasani2021}. A diretriz da plataforma é inegociável e tem três partes:

\begin{center}
\begin{tabularx}{\textwidth}{>{\sffamily\bfseries\footnotesize\color{unNavy}}p{4.6cm}X}
\toprule
\sffamily\bfseries\footnotesize\color{unNavy}Salvaguarda & \sffamily\bfseries\footnotesize\color{unNavy}O que significa na prática \\
\midrule
Geração na preparação, nunca ao vivo sem filtro & A IA é ferramenta de produtividade \textbf{para a professora}, no planejamento. Textos, personagens e posts simulados são revisados, ajustados e aprovados por ela antes de chegarem à turma. \\
\addlinespace
A criança interage com conteúdo curado & Em sala, o estudante vê um banco fechado e aprovado. Isso elimina por construção a exposição a resposta imprópria, discurso de ódio ou conteúdo de risco. \\
\addlinespace
Personagem com IA opera fechado & Quando a missão usa um personagem interativo, ele roda sob restrições rígidas de sistema, que o impedem de sair do tópico pedagógico ou emitir opinião própria. \\
\bottomrule
\end{tabularx}
\end{center}

\section{Privacidade: o artigo 14 da LGPD}

O tratamento de dados de menores no Brasil é regulado pelo \textbf{art.~14 da Lei n.\textsuperscript{o} 13.709/2018}~\cite{lgpd2018}, lido em conjunto com o Estatuto da Criança e do Adolescente~\cite{eca1990}:

\begin{destaque}
\itshape ``O tratamento de dados pessoais de crianças e adolescentes deverá ser realizado em seu melhor interesse, considerados os direitos fundamentais à vida, à privacidade e ao pleno desenvolvimento.'' \normalfont(art.~14, §\,1.\textsuperscript{o})
\end{destaque}

A esteira converte isso em três exigências verificadas --- duas delas por código, na dimensão de Dados e LGPD da banca:

\begin{enumerate}
  \item \textbf{Nada de captura abusiva de mídia.} É vedado exigir que o estudante grave vídeo do próprio rosto, foto de casa ou áudio aberto sem justificativa pedagógica indispensável e consentimento dos responsáveis.
  \item \textbf{Isolamento por turma e escola.} Todo acesso a banco de dados roda sob política estrita de isolamento (\cod{FORCE ROW LEVEL SECURITY}). Dado de criança de uma escola não pode, sob nenhuma hipótese de falha de software, ser visto por outra instituição.
  \item \textbf{Telemetria de progresso, não de pessoa.} Os eventos registram o percurso na missão --- que etapa foi concluída, quantas tentativas --- e nunca localização geográfica ou identificador de publicidade. O contrato de evento é conferido literalmente pelo agente 36: verbo, objeto, resultado e contexto, com tabela de leitura declarada.
\end{enumerate}

\section{Ética de acolhimento e não-rotulação}

\begin{center}
\begin{tikzpicture}[
  cx/.style={rectangle,rounded corners=4pt,text width=7.1cm,align=left,inner sep=10pt,font=\footnotesize}]
\node[cx,fill=sinalVermelho!7,draw=sinalVermelho,text=unNavy] (a)
 {{\sffamily\small\bfseries\color{sinalVermelho}Proibido}\\[5pt]
  Classificar o estudante --- ``aluno lento'', ``perfil desatento'', qualquer categorização psicológica, moral ou ideológica. Também é proibido fazê-lo sentir-se culpado ou incapaz por uma escolha dentro da simulação.};
\node[cx,fill=unPapel,draw=unNavy,text=unNavy,right=0.45cm of a] (b)
 {{\sffamily\small\bfseries\color{unNavy}O que a missão faz no lugar}\\[5pt]
  Julga \textbf{a escolha, nunca a pessoa}, e mostra a consequência dentro da história, com humor e empatia: ``Opa! O computador travou. Que pista deixamos passar?'' O erro é informação sobre o sistema, não sobre a criança.};
\end{tikzpicture}
\end{center}

Essa regra tem um efeito de projeto que costuma passar despercebido: ela \textbf{proíbe a tela de resposta certa}. Se o sistema declarasse quem acertou, ele estaria classificando --- e a comparação entre colegas, que é o coração do momento ``ahá'', viraria ranking. Por isso a revelação é conduzida pela professora com dados de comparação, e não por uma tela de pontuação.

\parteUnooria{VI}{Os agentes, um a um}{O que cada um faz, e em que se apoia para decidir. Sem essa coluna, uma rubrica acaba defendida por antiguidade.}

\section{Como ler este catálogo}
\label{ap:agentes}

Esta parte existe por uma razão específica: uma regra que não diz de onde veio é indefensável quando alguém a contesta --- e indefensável quando alguém deveria contestá-la e não o faz. Por isso cada agente aparece aqui com três informações: o que faz, e sobre que natureza de fundamentação a sua regra se apoia.

\begin{center}
\begin{tikzpicture}[
  nat/.style={rectangle,rounded corners=4pt,text width=4.45cm,align=left,inner sep=9pt,font=\footnotesize,text=unNavy}]
\node[nat,fill=unNavy!6,draw=unNavy] (a)
 {\Evidencia\\[5pt] Deriva de pesquisa publicada, com citação. Pode ser derrubada por pesquisa melhor --- e já foi, quatro vezes.};
\node[nat,fill=unAzul!12,draw=unAzul,right=0.35cm of a] (b)
 {\Contrato\\[5pt] Deriva de documento normativo: BNCC, LGPD, manifesto da plataforma. Não se pesquisa: confere-se.};
\node[nat,fill=unLavanda!15,draw=unLavanda,right=0.35cm of b] (c)
 {\Convencao\\[5pt] A equipe escolheu por consistência de produto. Legítimo, desde que declarado --- e sujeito a mudar sem pesquisa nova.};
\end{tikzpicture}
\end{center}

\begin{destaque}
São \textbf{72 posições numeradas}: 37 agentes com modelo de linguagem e \textbf{35 verificações determinísticas} --- código Python, sem modelo, marcadas em azul claro nas tabelas. Quase metade da esteira não usa modelo nenhum, e isso é o desenho, não uma economia: toda regra que \emph{pode} ser mecânica foi tirada dos avaliadores.

O revisor do \emph{one-pager} aparece duas vezes, como 07 (movido pelo gate estrutural) e como 10c (movido pelo parecer de adequação); a ficha do manifesto é conferida em Produto (15) e de novo em Dev (33), contra o mesmo vocabulário.
\end{destaque}

\section{Produto (00b--16b)}

\begin{center}
\begin{longtable}{p{2.6cm}p{5.2cm}p{7.2cm}}
\toprule
\multicolumn{1}{l}{\sffamily\bfseries\footnotesize\color{unNavy}Passo} & \multicolumn{1}{l}{\sffamily\bfseries\footnotesize\color{unNavy}O que faz} & \multicolumn{1}{l}{\sffamily\bfseries\footnotesize\color{unNavy}Em que se apoia} \\
\midrule
\endfirsthead
\toprule
\multicolumn{1}{l}{\sffamily\bfseries\footnotesize\color{unNavy}Passo} & \multicolumn{1}{l}{\sffamily\bfseries\footnotesize\color{unNavy}O que faz} & \multicolumn{1}{l}{\sffamily\bfseries\footnotesize\color{unNavy}Em que se apoia} \\
\midrule
\endhead
\textcolor{unAzul}{\textbf{00b}} Missões já geradas & Lê o índice e entrega ao núcleo as missões mais próximas do ano do briefing, para a nova não repetir uma anterior. & \Contrato\ interno. A semelhança é lexical e só \emph{recupera e ordena} --- quem reprova é o micro-gate. \\
\addlinespace
\textbf{00c} Mapa conceitual do material & Conceito limiar, conceitos núcleo e de apoio com trecho citado, crenças de entrada que o bimestre derruba, vocabulário do Explicar. & \Convencao\ com fundamentação em conceitos limiares. Existe porque o núcleo \emph{recebia} o material e nada pedia que o usasse --- ver Seção~\ref{sec:profundidade}. \\
\addlinespace
\textbf{00d} Ideação de mecanismos & Três mecanismos distintos para o conceito limiar, cada um com a anatomia do ``ahá'', pontuados na rubrica de profundidade. & \Convencao. O mecanismo que veio no briefing entra como candidato, não como decisão --- e trocá-lo continua sendo ato humano. \\
\addlinespace
\textbf{01} Núcleo narrativo & Descoberta, conceitos vividos em tela, anatomia do ``ahá'', narrativa e as seis etapas, a partir do briefing, dos anexos, do mapa e do mecanismo escolhido. & \Convencao\ --- o \emph{framework} de seis etapas é desenho de produto. A exigência de \emph{gap} de informação tem lastro em curiosidade epistêmica~\cite{loewenstein1994}. \\
\addlinespace
\textcolor{unAzul}{\textbf{01b}} Missões já geradas, pelo núcleo & Refiltra o índice pelo ambiente, eixo e descoberta do núcleo recém-escrito. É esta lista que o micro-gate julga. & \Contrato\ interno. \\
\addlinespace
\textcolor{unAzul}{\textbf{01c}} Cobertura conceitual & Núcleo $\times$ mapa: o limiar vive no Decidir ou no ``ahá'', todo conceito tem destino, a anatomia tem os seis passos, e \textbf{o trecho citado existe mesmo no material}. & \Contrato\ interno. A última regra existe porque citação inventada é o modo de falha mais fácil de um modelo obrigado a fundamentar. \\
\addlinespace
\textbf{02} Micro-gate do núcleo & \emph{Spoiler}, \emph{gap} de informação, escolha significativa, \emph{job-to-be-done} e semelhança com missão anterior. & \Convencao\ derivada de 01. Separar quem escreve de quem critica reduz autopreferência~\cite{zheng2023}. \\
\addlinespace
\textbf{02b} Juiz de profundidade & Nota de 1 a 4 em seis dimensões, com trecho citado e instrução de reescrita. & \Evidencia\ e \Convencao\ declarada: a anatomia do ``ahá'' tem lastro em previsão e conflito cognitivo; \textbf{os cortes estão por calibrar}. \\
\addlinespace
\textcolor{unAzul}{\textbf{03}} Veredito do micro-gate & Registra se o núcleo passou em \emph{spoiler}, \emph{gap}, escolha e \emph{job}. & \Contrato\ interno. \\
\addlinespace
\textcolor{unAzul}{\textbf{03b}} Veredito de profundidade & \textbf{Reconta o veredito nas notas do juiz}, não no que ele escreveu na conclusão. Funde a cobertura conceitual. & \Contrato\ interno, contra a leniência de avaliadores~\cite{zheng2023}. \\
\addlinespace
\textcolor{unAzul}{\textbf{03c}} Saída do loop do núcleo & Micro-gate \emph{e} profundidade aprovados encerram o ciclo; qualquer um reprovado devolve o núcleo com os dois relatórios. & \Contrato\ interno. \\
\addlinespace
\textbf{04} Seções de apoio & Completa o \emph{one-pager} canônico (seções 2 e 5--18) e o \emph{brief} de escopo. & \Convencao\ --- formato de produto. A separação em relação a 01 é deliberada: densidade de instrução degrada o cumprimento~\cite{ifscale2025}. \\
\addlinespace
\textcolor{unAzul}{\textbf{04b}} Contrato de plataforma & \emph{Lint} a cada rodada do gate: canal de sala por consulta periódica (nunca WebSocket), estado no servidor (nunca no navegador), SAEB só nos anos avaliados, número com origem declarada. & \Contrato\ da plataforma. Nasceu de erros reais do \emph{one-pager} gerado, fechados na origem em vez de corrigidos a cada vez. \\
\addlinespace
\textbf{05} Gate estrutural & Critérios de estrutura, ambiente, logística, evidências, viabilidade e infraestrutura. Escreve a correção pronta para colar. & \Convencao\ com critérios de \Evidencia: teto de decisões e duração vêm da fonte única. Rubrica com teste prático e contraexemplo sustenta concordância~\cite{jonsson2007}. \\
\addlinespace
\textcolor{unAzul}{\textbf{06}} Veredito do gate & Lê o envelope JSON e decide repetir, seguir ou bloquear. & \Contrato\ interno. \\
\addlinespace
\textbf{07} Revisor do \emph{one-pager} & Aplica o complemento que o gate escreveu; não inventa decisão de produto. & \Convencao. Separa julgar de consertar. \\
\addlinespace
\textbf{08} Avaliação pedagógica & Currículo, papel do professor, \emph{guardrails} e a relação entre o que a interface permite e o que ela protege. & \Contrato: BNCC Computação~\cite{cne2022}. A regra de que só letramento digital bloqueia é decisão de governança, declarada como tal. \\
\addlinespace
\textbf{09} Experiência do estudante & Carga cognitiva, faixa etária, linguagem, duração e agência. & \Evidencia\ --- o passo mais ancorado da esteira~\cite{gathercole2004,cowan2016,godwin2016,betts2006,moreno2002,kalyuga2003}. Todos os números vêm da fonte única. \\
\addlinespace
\textbf{10} Adequação pedagógica & Parecer em seis eixos com ajustes acionáveis. Entra como ressalva. & \Evidencia~\cite{biggs1982,osterhaus2021,starks2024}; \Contrato\ para a BNCC; \Convencao\ para o placar. Não carrega régua própria. \\
\addlinespace
\textcolor{unAzul}{\textbf{10b}} Veredito do parecer & Parecer sem ajuste encerra o ciclo; com ajuste, devolve ao revisor. & \Contrato\ interno. \\
\addlinespace
\textbf{10c} Revisor movido pelo parecer & O mesmo agente 07, aplicando os ajustes que o parecer exigiu. Só roda com o teto acima de 1. & \Convencao. \\
\addlinespace
\textbf{11} Gate final de produto & Coerência entre os três pareceres e veredito de liberação. & \Convencao. Existe porque avaliadores paralelos podem julgar \emph{versões diferentes} do mesmo documento~\cite{cemri2025}. \\
\addlinespace
\textcolor{unAzul}{\textbf{12}} Veredito final & Lê o envelope do gate final. & \Contrato\ interno. \\
\addlinespace
\textbf{13} \emph{Spec} da missão & Os 12 blocos do contrato contra o qual Design e Dev constroem. & \Contrato\ --- manifesto da plataforma, isolamento por linha, contrato de eventos. \\
\addlinespace
\textcolor{unAzul}{\textbf{14}} Conformidade BNCC & Código existe, é do ano, e há letramento digital. Confere também \textbf{o que a \emph{spec} diz sobre a habilidade}, não só o código. & \Contrato, verificado em código~\cite{cne2022}. \\
\addlinespace
\textcolor{unAzul}{\textbf{15}} Ficha do manifesto & Dez campos do bloco 0 contra o vocabulário fechado. & \Contrato. Roda em Produto porque o conserto ali custa uma reescrita da \emph{spec}, não o descarte do Design. \\
\addlinespace
\textcolor{unAzul}{\textbf{15b}} Matriz de evidência & Grau \cod{scored} com dois sinais nomeados, toda habilidade com momento em tela e medida ligada, gabarito no servidor. & \Contrato. Das 25 medidas \cod{scored} das primeiras execuções, só 10 nomeavam dois sinais. \\
\addlinespace
\textcolor{unAzul}{\textbf{15c}} Perguntas de produto & Quantidade fechada, decisão com limiar, corte do vocabulário, e a ação que responde cada pergunta declarada na tabela de eventos. & \Contrato\ interno. \\
\addlinespace
\textbf{15d} Juiz das perguntas & O que regex não vê: a pergunta muda uma decisão nossa, é sobre o produto e não sobre a criança, e não é uma transversal repintada. & \Convencao. \\
\addlinespace
\textcolor{unAzul}{\textbf{15e}} Veredito do juiz & Reprovação em pergunta declarada pelo responsável de produto bloqueia; em proposta, segue como ressalva para ele decidir. & \Contrato\ interno. \\
\addlinespace
\textcolor{unAzul}{\textbf{16}} Veredito da \emph{spec} & Currículo, manifesto, evidência e perguntas conformes encerram o ciclo. & \Contrato\ interno. \\
\addlinespace
\textbf{16b} \emph{One-pager} resumido & O formato curto do time de produto: ficha, ideia, componente curricular, socioemocionais, seis etapas, professor, contrato de evidência e \emph{takeaway}. Sai em três formatos. & \Convencao\ de comunicação, com regras explícitas de registro --- ver Seção~\ref{sec:resumido}. \\
\bottomrule
\end{longtable}
\end{center}

\section{Design (17--32)}

\begin{center}
\begin{longtable}{p{2.6cm}p{5.2cm}p{7.2cm}}
\toprule
\sffamily\bfseries\footnotesize\color{unNavy}Agente & \sffamily\bfseries\footnotesize\color{unNavy}O que faz & \sffamily\bfseries\footnotesize\color{unNavy}Em que se apoia \\
\midrule
\endfirsthead
\toprule
\sffamily\bfseries\footnotesize\color{unNavy}Agente & \sffamily\bfseries\footnotesize\color{unNavy}O que faz & \sffamily\bfseries\footnotesize\color{unNavy}Em que se apoia \\
\midrule
\endhead
\textbf{17} Fechamento de lacunas & Varre a \emph{spec} célula a célula, registra decisões com base declarada e pedidos de re-\emph{spec}, e consolida. & \Convencao. Preenche ausência, nunca altera presença --- a regra existe para que Design não redecida produto. \\
\addlinespace
\textcolor{unAzul}{\textbf{18}} Separação e consolidação & Copia a \emph{spec} consolidada, preserva a original e separa o relatório de preenchimento. & \Contrato\ interno. \\
\addlinespace
\textbf{19} Ambiente simulado & Condicional: decide o ambiente quando a \emph{spec} cala --- app clonado, app novo ou mini-jogo. Dormente, emite só a ficha que os demais consomem. & \Contrato\ (lista fechada de cinco seções do \emph{design system}) mais \Convencao\ para o mini-jogo dos anos iniciais. \\
\addlinespace
\textcolor{unAzul}{\textbf{20}} \emph{Spec} com o bloco do ambiente & Copia a ficha do ambiente para a \emph{spec} e registra o estado. & \Contrato\ interno. \\
\addlinespace
\textbf{21} Briefing de ofício & Aposta, forma do arco, pico, \emph{beat} memorável, mapa de \emph{beats} e não-negociáveis. & \Convencao\ de ofício narrativo. \textbf{É a parte menos fundamentável da esteira, e está declarada assim}: são heurísticas de dramaturgia, não achados. \\
\addlinespace
\textbf{22} Jornada · \textbf{23} Fluxo & Raias por etapa, cenários de dispositivo e inventário de telas; grafo de estados com condições, gates e cinco testes de fluxo. Ambos em Mermaid. & \Convencao. O diagrama existe porque representação externa reduz a carga de manter a estrutura na cabeça~\cite{sweller2019}. \\
\addlinespace
\textbf{24} · \textbf{27} \emph{Copy} (duas passadas) & Dicionário de \emph{strings} por tela, calibrado à faixa; depois revisão dentro da tela já construída. & \Evidencia\ para a calibragem de linguagem por faixa (fonte única) e \Convencao\ para o formato de duas passadas. \\
\addlinespace
\textbf{25} Camada CAA & Símbolos, modos alternativos de resposta e pranchas de comunicação. Só roda quando a faixa ou a \emph{spec} pedem. & \Contrato\ de acessibilidade. A regra de que o símbolo vem \emph{junto} do texto, nunca no lugar dele, é redundância deliberada e benéfica para quem não decodifica com fluência~\cite{knoop2018}. \\
\addlinespace
\textbf{26} Telas & Especificação por tela e wireframes em escala de cinza sobre a seção do \emph{design system} do ambiente. & \Contrato\ (tokens do \emph{design system}) mais \Evidencia\ para o teto de elementos novos por tela~\cite{sweller2019}. \\
\addlinespace
\textbf{28} Protótipo & Arquivo HTML único, sem dependências, com as telas navegáveis, o painel da etapa e as cadências do pico. & \Convencao. É o artefato que a revisão humana abre; a exigência de arquivo único e offline é operacional. \\
\addlinespace
\textcolor{unAzul}{\textbf{29}} Artefatos e links & Separa o HTML do relatório, colhe os links das telas e monta a folha de diagramas. & \Contrato\ interno. \\
\addlinespace
\textcolor{unAzul}{\textbf{29a}} \emph{Lint} do protótipo & Confere o arquivo HTML em si --- o que o QA de entregáveis não vê, porque lê só o texto. Reprovado, devolve ao protótipo no retrabalho. & \Contrato\ interno. Existe porque o QA auditava a \emph{descrição} do protótipo e dava por bom um arquivo que não abria como descrito. \\
\addlinespace
\textcolor{unAzul}{\textbf{29b}} Espelho no Miro & Publica no Miro os mesmos diagramas Mermaid, como board navegável. & \Convencao. Publicação determinística do que já foi escrito; o texto continua sendo a fonte e a integração indisponível não interrompe a esteira. \\
\addlinespace
\textbf{30} QA de entregáveis & Fidelidade à \emph{spec}, \emph{diff} de \emph{strings}, currículo, faixa, fluxo, \textbf{visão do professor}, briefing, ambiente e CAA. & \Convencao, com uma regra de \Evidencia: é proibido reportar número que não foi medido --- o \emph{diff} de \emph{strings} só conta se houve medição. \\
\addlinespace
\textcolor{unAzul}{\textbf{31}} Veredito de design & Lê o envelope do QA. & \Contrato\ interno. \\
\addlinespace
\textcolor{unAzul}{\textbf{32}} Retrabalho de design & Reexecuta só os agentes com \emph{ticket}, na ordem de precedência, e refaz o protótipo depois. & \Contrato\ interno. O despacho é seletivo porque redundância com falhas correlacionadas não aumenta cobertura~\cite{knight1986}. \\
\bottomrule
\end{longtable}
\end{center}

\section{Dev (33--49) e transversal (50--52)}

\begin{center}
\begin{longtable}{p{2.6cm}p{5.2cm}p{7.2cm}}
\toprule
\sffamily\bfseries\footnotesize\color{unNavy}Agente & \sffamily\bfseries\footnotesize\color{unNavy}O que faz & \sffamily\bfseries\footnotesize\color{unNavy}Em que se apoia \\
\midrule
\endfirsthead
\toprule
\sffamily\bfseries\footnotesize\color{unNavy}Agente & \sffamily\bfseries\footnotesize\color{unNavy}O que faz & \sffamily\bfseries\footnotesize\color{unNavy}Em que se apoia \\
\midrule
\endhead
\textcolor{unAzul}{\textbf{33}} Ficha do manifesto & Bloco 0 contra as listas fechadas: produto, séries, unidade, público, \emph{viewport}, permissões, \emph{capabilities}, seção do DS, liberação por turma. Não conforme impede o PRD. & \Contrato\ da plataforma. \\
\addlinespace
\textbf{34} \emph{Handoff} para dev & Modelo de insumo preenchido, \emph{stories} com critérios fechados, eventos, escopo V1/V2 e hipóteses. & \Contrato\ da plataforma: assinatura de evento, papéis, canal de sala. \\
\addlinespace
\textbf{35} \emph{Brief} de assets & Inventário, especificação por tipo, critérios de reprovação e ordem de produção. & \Convencao\ de produção. \\
\addlinespace
\textcolor{unAzul}{\textbf{36}} Contrato de telemetria & Evento = verbo · objeto · resultado · contexto; medida com tabela de leitura; sem WebSocket; seis papéis; quota declarada. & \Contrato. Comparação literal contra o contrato real da API --- não há julgamento envolvido. \\
\addlinespace
\textcolor{unAzul}{\textbf{37}} Veredito do \emph{handoff} & Telemetria conforme encerra o ciclo; reprovada devolve o \emph{handoff} com o relatório. & \Contrato\ interno. \\
\addlinespace
\textcolor{unAzul}{\textbf{38}} Rastro de lacunas & Marcadores em \emph{spec}, design, \emph{handoff} e assets, \textbf{por dono}: o que faltou de insumo, produto, design ou dev impede o PRD; o que depende da plataforma vira pendência humana \statusAmbar. & \Contrato\ interno. A categoria \emph{humano} foi criada depois da primeira execução completa --- ver Seção~\ref{sec:execucoes}. \\
\addlinespace
\textbf{39--43} Banca de cinco dimensões & Produto e narrativa; design e imersão; pedagogia e BNCC; dados, telemetria e LGPD; técnica e arquitetura. Em paralelo, com critérios numerados por dimensão. & \Contrato\ em quatro das cinco; \Evidencia\ em pedagogia. A separação existe porque um avaliador único com 70 critérios erra de forma sistemática~\cite{ifscale2025} e porque falhas correlacionadas não se compensam~\cite{knight1986}. \\
\addlinespace
\textbf{44} Consolidação da banca & Veredito único \statusVerde\,/\,\statusAmbar\,/\,\statusVermelho: \textbf{a pior das cinco dimensões}, com os bloqueantes lidos do arquivo de critérios. & \Contrato: regra determinística, sem bom senso no meio. A consolidação deixou de ser feita por um modelo justamente porque ele arredondava para cima. \\
\addlinespace
\textcolor{unAzul}{\textbf{45}} Veredito de dev & Lê o envelope da banca. & \Contrato\ interno. \\
\addlinespace
\textcolor{unAzul}{\textbf{46}} Retrabalho da banca & Reexecuta o dono de cada achado --- inclusive quando o dono está em Design. Achado de Produto vira pendência nomeada. & \Contrato\ interno. \\
\addlinespace
\textbf{47} PRD padrão ouro & As 25 seções que alimentam quem codifica, com o contexto arquitetural da plataforma. & \Contrato\ da plataforma e da LGPD (dado de criança e adolescente). \\
\addlinespace
\textcolor{unAzul}{\textbf{48}} Pendências e consistência do PRD & Marcadores deixados pelo redator, mais a consistência interna: prefixo de tabela único, requisito P0 com teste e tarefa, evento com tabela leitora, caminho único por módulo, tarefa atômica. & \Contrato. Determinístico porque é comparação literal --- e porque o PRD é o documento com mais restrições simultâneas da esteira. \\
\addlinespace
\textcolor{unAzul}{\textbf{48a}} Coerência KPI $\leftrightarrow$ evento $\leftrightarrow$ tela & Todo indicador com pergunta, decisão, corte e um evento que dispara \emph{na tela que ele nomeia}; os sinais universais das perguntas transversais presentes na tabela de eventos. & \Contrato\ interno. Um indicador que ninguém consegue calcular é pior que indicador nenhum: ele promete um número que não existe. \\
\addlinespace
\textcolor{unAzul}{\textbf{48b}} Veredito do PRD & PRD consistente encerra o ciclo; reprovado devolve ao redator, até o teto. & \Contrato\ interno. Foi o último ciclo a ganhar volta --- ver Seção~\ref{sec:terceira}. \\
\addlinespace
\textbf{49} \emph{One-pager} pedagógico & A missão inteira em linguagem de sala de aula, com o papel do professor, os links do protótipo e as perguntas de aprovação. Sai em Markdown e \cod{.docx}. & \Convencao\ de comunicação, com uma regra que vale citar: \textbf{nenhuma ressalva sobrevivente pode ser omitida}, mesmo com veredito verde. Avaliadores tendem à leniência~\cite{zheng2023}; a regra existe contra isso. \\
\addlinespace
\textcolor{unAzul}{\textbf{50}} Diagnóstico da execução & O dossiê: como cada portão se saiu, o que trava, de quem é cada conserto, o que saiu provisório e que documento faltou. Sem modelo de linguagem. & \Contrato\ interno. \textbf{Roda sempre}, inclusive quando tudo o mais falhou --- é a peça que sustenta a decisão de não abortar. \\
\addlinespace
\textbf{51} Aprendizado & A leitura em prosa do diagnóstico, mais padrões aprendidos, lacunas com dono e candidatos a critério. & \Convencao. \\
\addlinespace
\textcolor{unAzul}{\textbf{52}} Registro no catálogo & A memória entre execuções: registra a missão no índice, mostra as mais próximas e acusa identificador repetido. \emph{Status} derivado dos gates. & \Contrato\ interno. \textbf{Promover é ato humano} --- nenhum agente marca uma missão como aprovada. Missão bloqueada entra no índice de propósito: o que se quer evitar é reinventá-la. \\
\bottomrule
\end{longtable}
\end{center}

\parteUnooria{VII}{O que a operação ensinou}{Quatro execuções completas sobre material real expuseram defeitos que nenhum teste unitário pegaria --- porque não eram erros de código.}

\section{O refino dos entregáveis}

Uma revisão dedicada ao \emph{one-pager} pedagógico e ao PRD perguntou o que a literatura recomendaria acrescentar a cada um. A resposta, nos dois casos, foi a mesma: \textbf{nada de conteúdo novo}. Os dois já estavam acima da média do gênero. O que faltava era \textbf{ordem, navegação e verificação}.

Um terceiro documento nasceu depois dessa revisão, e a razão dele é diferente das outras duas --- é sobre \emph{quem lê}.

\subsection{O \emph{one-pager} resumido: quando o documento certo é o curto}
\label{sec:resumido}

Em 21 de setembro comparamos o \emph{one-pager} que o time de produto escrevia à mão --- cerca de 2.550 palavras --- com o que a esteira gerava: cerca de 8.700. O veredito não foi que um era melhor.

\begin{center}
\begin{tikzpicture}[
  cx/.style={rectangle,rounded corners=4pt,text width=7.0cm,align=left,inner sep=10pt,font=\footnotesize}]
\node[cx,fill=unPapel,draw=unLavanda,text=unNavy] (a)
 {{\sffamily\small\bfseries\color{unNavy}O gerado, de 8.700 palavras}\\[5pt]
  Superior em \textbf{cobertura}: passa por todos os gates, documenta cada decisão, sustenta auditoria.\\[4pt]
  E é o documento que ninguém lê inteiro para aprovar uma missão.};
\node[cx,fill=unNavy,text=white,right=0.45cm of a] (b)
 {{\sffamily\small\bfseries O manual, de 2.550 palavras}\\[5pt]
  Superior em \textbf{leitura} --- e trazia quatro coisas que o gerado não tinha: componente curricular do bimestre, diferenciação frente às missões já feitas, contrato de evidência por habilidade e o \emph{takeaway}.};
\end{tikzpicture}
\end{center}

A decisão foi ter os dois, e não escolher: o resumido é o que o time de produto de fato lê para aprovar; o canônico é insumo de auditoria. Os quatro itens ausentes entraram --- e o contrato de evidência, que nasceu aqui, virou portão determinístico (Parte~V).

\paragraph{A segunda reclamação foi sobre tom, e o diagnóstico vale como método.} Com o formato certo, o responsável de produto reprovou a \emph{linguagem}: ``muito acadêmica e difícil de ler''. A causa não estava no conteúdo. O \emph{prompt} regulava fidelidade, códigos, tabelas e tamanho --- e \textbf{não dizia nada sobre como soar}. Sem regra de registro, o modelo cai no português pedagógico formal e, pior, imita o registro da própria instrução que recebeu.

A correção foram duas seções novas no \emph{prompt}: regras de linguagem com tabela de substituição de jargão, e calibragem de tom com cinco pares de exemplo extraídos do documento que o time considerava padrão-ouro. A estrutura de itens foi preservada intacta --- só os textos mudaram.

\begin{destaque}
\textbf{Reclamação de leitura é problema de regra de registro no \emph{prompt}, não de conteúdo.} O diagnóstico é direto: procure no \emph{prompt} se há alguma regra sobre \emph{como soar}. Se não houver, é isso.
\end{destaque}

\subsection{\emph{One-pager}: o problema era onde as coisas estavam}

A pesquisa sobre adoção de material didático converge num ponto: a professora decide por \textbf{praticidade} --- o material diz concretamente o que fazer, cabe na turma que ela tem, e compensa o esforço~\cite{doyle1977}. E em adoção de tecnologia escolar, a preocupação mais intensa é \emph{operacional}, não teórica~\cite{cbam2015}. Três consequências:

\begin{itemize}
  \item \textbf{A decisão pedida subiu para o topo.} Ela aparecia a 80\% do documento. Quem lê precisa saber na primeira tela o que está sendo pedido; se parar ali, ainda sabe o que fazer.
  \item \textbf{Entrou ``o que sua sala precisa ter''}, em lista fechada --- aparelhos, projetor, internet, material impresso, minutos de preparo. Antes isso estava disperso em prosa, e era justamente a primeira pergunta que a professora faz.
  \item \textbf{As perguntas de aprovação viraram instrumento}, com opção de resposta e assinatura, em vez de convite a opinião livre. Material sem expectativa clara de implementação não é usado.
\end{itemize}

E um defeito concreto que só a execução real expôs: o documento dizia ``o protótipo pode ser aberto no navegador'' \textbf{sem nenhum endereço}, e o teste que deveria pegar isso não bloqueava. O protótipo é o único item que deixa a professora simular a aula antes de aprovar --- o teste passou a ser bloqueante. Saíram também as URLs longas de servidor de arquivo coladas no corpo do texto: são ruído de leitura e vencem em poucos dias.

\subsection{PRD: o problema era que nada verificava o que ele afirmava}

Aqui a literatura foi mais direta ainda. Especificação bem escrita ajuda pouco; especificação \textbf{verificada contra o repositório} ajuda de forma mensurável~\cite{speckit2026}, e o que reduz alucinação em geração de código é símbolo e caminho verificáveis, não prosa mais detalhada~\cite{dehallucinator2024}. Somado a isso, documento longo não é vantagem: a recuperação piora no meio de contextos grandes~\cite{liu2024}.

As 25 seções ficaram como estavam. O que entrou foi um \textbf{gate determinístico de consistência interna}, porque o gabarito já exigia essas regras em prosa e ninguém as conferia. Ele barra:

\begin{center}
\begin{tabularx}{\textwidth}{>{\sffamily\bfseries\footnotesize\color{unNavy}}p{5.4cm}X}
\toprule
\sffamily\bfseries\footnotesize\color{unNavy}O que ele barra & \sffamily\bfseries\footnotesize\color{unNavy}Por que isso quebra quem codifica \\
\midrule
Mais de um prefixo de tabela no mesmo documento & O prefixo nomeia o esquema inteiro. Dois prefixos significam dois bancos possíveis, e a escolha vira sorteio. \\
Requisito P0 fora da matriz e sem tarefa & Um requisito bloqueante sem rastro é um requisito que ninguém vai implementar, e ninguém vai notar que não implementou. \\
Evento com leitura em fila, sem tabela própria & Contradiz o contrato real de telemetria: o dado ficaria sem onde ser lido. \\
Tarefa que declara cobrir ``todos os requisitos'' & Não é decomposição. Uma tarefa assim não pode ser estimada, revisada nem dada como concluída. \\
Nome de módulo que não é identificador válido & Quebra no primeiro \emph{import}. \\
Dois caminhos diferentes para o mesmo diretório & Metade do código vai para um lugar, metade para o outro. \\
\bottomrule
\end{tabularx}
\end{center}

Nenhum desses é lacuna de informação --- são \textbf{contradições internas}. Um PRD pode não ter marcador de pendência nenhum e ainda assim fazer quem codifica escolher um caminho ao acaso; a partir daí, o resto do documento deixa de valer contra o que foi escrito.

\Needspace*{10\baselineskip}
\section{O que a primeira e a segunda execução ensinaram}
\label{sec:execucoes}

Rodar a esteira inteira sobre um material real expôs três defeitos que nenhum teste unitário pegaria, porque não são erros de código: são lacunas de contrato. Valem como princípio, não como conserto pontual.

\begin{enumerate}
  \item \textbf{Ato que só uma pessoa pratica não é lacuna de construção.} Homologar autoria, indexar no material editorial, licenciar um acervo gráfico --- nenhum agente produz isso. Os gates de rastro tratavam esses marcadores como impedimento, o que reprovava a missão por algo que nenhuma rodada poderia resolver. Pior: como o texto do marcador não nomeava uma etapa da esteira, o achado caía na etapa padrão de quem o leu, e uma homologação editorial aparecia endereçada ao time de desenvolvimento. Criou-se a categoria \textbf{humano}: severidade \emph{decisão humana}, não bloqueio, e dono correto no dossiê.
  \item \textbf{O prefixo da tabela é o identificador inteiro, sem separador interno.} Com o identificador novo carregando ano e bimestre, referências que ainda ensinavam o formato antigo levaram um avaliador a pedir um prefixo híbrido, que não corresponde a nenhum dos dois padrões. A regra passou a estar escrita no gabarito da \emph{spec}, com contraexemplo.
  \item \textbf{Insumo incompleto deve produzir pendência nomeada, não reprovação genérica.} Quando o material de origem não traz autoria, revisão ou indexação curricular, isso é informação ausente de quem forneceu o insumo --- e o relatório precisa dizer exatamente isso, com o dono, em vez de devolver a missão como malfeita. \textbf{A régua não baixa; o endereço do problema é que muda.}
\end{enumerate}

\textbf{Uma quarta, achada na segunda execução: citação do formato não é pendência.} Um agente escreveu, como afirmação de conformidade, a frase ``nenhuma pendência \cod{[ABERTO --- <dono>]} sobrevive nesta \emph{spec}''. O varredor de lacunas leu a citação como pendência real e ainda propagou o \emph{placeholder} como se fosse o responsável, produzindo no dossiê a linha sem sentido ``comece por: \cod{[ABERTO --- <dono>]}''. A guarda contra citação de formato existia, mas não cobria \emph{placeholder} entre sinais de menor e maior --- que nunca é dono real.

\subsection{O efeito medido}

As correções foram verificadas rodando a mesma missão duas vezes, com a mesma entrada.

\begin{center}
\begin{tikzpicture}[yscale=0.52,xscale=1.0]
\foreach \y in {0,1,...,14} \draw[draw=unLinha,very thin] (-0.2,\y) -- (11.2,\y);
% antes
\fill[unLavanda!45,rounded corners=2pt] (0.4,0) rectangle (1.5,5);
\fill[unLavanda!45,rounded corners=2pt] (4.4,0) rectangle (5.5,13);
\fill[unLavanda!45,rounded corners=2pt] (8.4,0) rectangle (9.5,1);
% depois
\fill[unNavy,rounded corners=2pt] (1.7,0) rectangle (2.8,1);
\fill[unNavy,rounded corners=2pt] (5.7,0) rectangle (6.8,14);
\fill[unNavy,rounded corners=2pt] (9.7,0) rectangle (10.8,1);
\node[font=\sffamily\scriptsize\bfseries,text=unNavy] at (0.95,5.9) {5};
\node[font=\sffamily\scriptsize\bfseries,text=white] at (2.25,0.6) {1};
\node[font=\sffamily\scriptsize\bfseries,text=unNavy] at (4.95,13.9) {13};
\node[font=\sffamily\scriptsize\bfseries,text=white] at (6.25,13.0) {14};
\node[font=\sffamily\scriptsize\bfseries,text=unNavy] at (8.95,1.9) {1};
\node[font=\sffamily\scriptsize\bfseries,text=white] at (10.25,0.6) {1};
\node[font=\sffamily\scriptsize,text=unNavy,align=center] at (1.6,-1.4) {Impedimentos};
\node[font=\sffamily\scriptsize,text=unNavy,align=center] at (5.6,-1.4) {Ressalvas};
\node[font=\sffamily\scriptsize,text=unNavy,align=center] at (9.6,-1.4) {Decisões humanas};
\draw[draw=unLinha] (-0.2,0) -- (11.2,0);
\node[font=\sffamily\scriptsize,text=unLavanda,anchor=west] at (0.4,15.4) {\textcolor{unLavanda}{$\blacksquare$} antes \quad \textcolor{unNavy}{$\blacksquare$} depois};
\end{tikzpicture}
\captionof{figure}{Mesma missão, antes e depois das correções. Os quatro impedimentos que sumiram eram atos humanos que a esteira cobrava de si mesma.}
\end{center}

\begin{destaque}
\textbf{Nenhuma das correções baixou a régua.} O que mudou foi o endereço do problema, não a exigência --- e a contagem de ressalvas até subiu, porque impedimento falso virou ressalva verdadeira.
\end{destaque}

\section{O que a terceira e a quarta execução ensinaram}
\label{sec:terceira}

A terceira execução não trouxe mais nenhum falso impedimento: os cinco achados eram todos defeitos reais do documento final, e vinham do mesmo lugar --- o portão de consistência do PRD. Isso expôs um problema estrutural que as execuções anteriores escondiam, porque antes esse portão não existia.

\subsection{Um portão que julga precisa poder devolver}

O PRD é o último documento da esteira, e o seu portão era o único que \textbf{detectava sem poder corrigir}: apontava o defeito e seguia para o \emph{one-pager} com o defeito intacto, deixando para uma pessoa aquilo que tem correção certa e mecânica. Os outros seis ciclos já devolviam o achado ao autor. O do PRD passou a devolver também, com o mesmo padrão e teto próprio de rodadas.

\begin{caixaConceito}{O princípio, que vale muito além deste caso}
\textbf{Gate sem ciclo de volta transfere para a pessoa um trabalho que a máquina sabe fazer.} Um achado determinístico com correção única --- prefixo de tabela fora do padrão, identificador inválido, requisito bloqueante sem rastro --- não é uma decisão: é uma reescrita. O que deve chegar ao humano é a pendência de insumo e a decisão de escopo, nunca o erro de contrato que o próprio documento permite localizar.
\end{caixaConceito}

\subsection{Por que o redator errou com a regra escrita no \emph{prompt}}

As três regras violadas estavam todas escritas de forma explícita no gabarito do redator, com contraexemplo. Não é caso de instrução ausente, e sim de \textbf{densidade de instrução}: o cumprimento cai à medida que o número de regras cresce no mesmo \emph{prompt}, e o PRD é o documento com mais restrições simultâneas da esteira. Acrescentar ênfase ao \emph{prompt} não resolveria --- por isso a resposta foi um verificador determinístico \emph{fora} do modelo, mais o ciclo que devolve o relatório para reescrita.

\subsection{Um achado que não faz sentido é pior que nenhum achado}

A quarta execução derrubou os impedimentos de cinco para um, e o que sobrou era \textbf{falso}: o verificador acusava um requisito P0 ``sem rastro'' que estava declarado, estava na matriz e estava coberto por duas tarefas. A regra contava ocorrências do identificador e exigia três; mas um PRD bem escrito cita o requisito \emph{duas} vezes --- na declaração e na matriz --- porque é a matriz que nomeia a tarefa. E a cobertura costuma vir por intervalo (\cod{RF-001..RF-014}), que não repete o identificador do meio.

O mais instrutivo é o que essa cegueira \textbf{escondia}. Ao passar a ler intervalos, o verificador enxergou cinco tarefas cobrindo os catorze requisitos cada uma --- ou seja, o plano de implementação não estava decomposto, defeito bem mais sério que o falso positivo apontado. A regra de atomicidade existia, mas só reconhecia a ausência de decomposição quando escrita por extenso (``todos os RFs''); escrita como intervalo, passava batido.

\begin{caixaPratica}{Dois princípios de escrita de verificador determinístico}
\textbf{Contar ocorrências é uma aproximação de rastreabilidade, não rastreabilidade.} O teste certo não é quantas vezes o identificador aparece, e sim se existe uma linha que ligue o requisito a uma tarefa. Trocar a contagem pela relação corrigiu o falso positivo e, no mesmo movimento, \emph{apertou} a régua: o impedimento inventado virou cinco ressalvas verdadeiras.

\textbf{Achado que não carrega a evidência obriga quem lê a refazer a investigação.} Um dos achados citava a linha defeituosa cortada em 90 caracteres, e o corte caía justamente antes do identificador do evento: o relatório dizia que havia um evento errado sem dizer qual. O achado passou a nomear o identificador e a carregar a linha inteira. Relatório assim é indistinguível de nenhum relatório.
\end{caixaPratica}

\Needspace*{12\baselineskip}
\section{A devolutiva que mudou a arquitetura}

As correções descritas até aqui vieram de rodar a esteira e ler o que ela produzia. A de 24 de setembro veio de fora: \textbf{o time pedagógico devolveu as missões}. E ela é a mais instrutiva das quatro, porque o sistema não tinha errado --- ele estava respondendo corretamente a uma pergunta que não era a certa.

\begin{center}
\begin{tikzpicture}[
  cx/.style={rectangle,rounded corners=4pt,text width=4.42cm,align=left,inner sep=9pt,font=\footnotesize,text=unNavy},
  seta/.style={-{Stealth[length=2.2mm]},draw=unLavanda,thick}]
\node[cx,fill=unPapel,draw=unLavanda] (a)
 {{\sffamily\small\bfseries O que o semáforo media}\\[4pt] Estrutura, plataforma, currículo, carga. \textbf{73 critérios}, todos cumpridos. Verde.};
\node[cx,fill=unPapel,draw=unLavanda,right=0.35cm of a] (b)
 {{\sffamily\small\bfseries O que o time media}\\[4pt] A missão trabalha os conceitos do material? O ``ahá'' tem ápice? \textbf{Um critério} cobria isso, e qualquer missão rasa passava nele.};
\node[cx,fill=sinalVermelho!8,draw=sinalVermelho,right=0.35cm of b] (c)
 {{\sffamily\small\bfseries O efeito}\\[4pt] O responsável de produto liberava confiando no verde. \textbf{O semáforo não estava errado --- estava respondendo a outra pergunta.}};
\draw[seta] (a)--(b); \draw[seta] (b)--(c);
\end{tikzpicture}
\end{center}

Daí saíram as duas mudanças estruturais desta edição, e elas se sustentam mutuamente.

\subsection{Medir o que o time mede}

A primeira foi a camada de profundidade (Seção~\ref{sec:profundidade}): mapa conceitual, ideação, anatomia do ``ahá'' e um juiz próprio. O princípio é simples de enunciar e fácil de esquecer:

\begin{destaque}
\textbf{Um gate que não mede o que a revisão humana mede não protege --- ele autoriza.} O verde não era um erro de cálculo: era uma afirmação verdadeira sobre uma pergunta irrelevante, lida como se respondesse à pergunta que importava.
\end{destaque}

\subsection{Pôr uma pessoa na passagem}

A segunda foi o parecer humano entre etapas (Seção~\ref{sec:passagem}). Ela responde a uma objeção óbvia à primeira: se o gate automático pode estar medindo a coisa errada, então nenhuma régua nova resolve o problema de fundo --- ela só o move.

Por isso a esteira deixou de correr sozinha do briefing ao PRD. Hoje cada etapa termina com uma pessoa que \textbf{responde por ela}, com comentário obrigatório que segue para as etapas seguintes. A máquina continua fazendo o que faz bem --- conferir o que é mecânico, em escala, sem se cansar --- e parou de ser a última palavra sobre o que é bom.

\begin{caixaConceito}{Três decisões de desenho que valem além deste caso}
\textbf{O comentário do \emph{de acordo} é obrigatório e circula.} Não é registro de auditoria: ele entra no contexto dos agentes seguintes. Aprovação silenciosa não transporta o que a pessoa sabia.

\textbf{Devolver não refaz nada sozinho.} Quem recebe uma etapa e discorda devolve \emph{o parecer} para quem o deu --- e é essa pessoa que decide entre refazer e reafirmar. A alternativa, disparar o retrabalho automaticamente, transfere a decisão para quem não respondeu por ela.

\textbf{A exceção fica marcada.} Com gate vermelho não há \emph{de acordo}; só a conta de administração libera, e a liberação entra no dossiê como exceção. O sistema não impede que se passe por cima de um bloqueio --- ele impede que isso passe despercebido.
\end{caixaConceito}

\subsection{E uma coisa que a esteira passou a saber}

Até 22 de setembro, o gerador não tinha memória entre execuções: o agente que inventava a missão não via nenhuma das anteriores, e duas missões parecidas só eram percebidas por acaso. Hoje há um índice, alimentado automaticamente ao fim de cada execução, que o micro-gate consulta.

A decisão de desenho que vale registrar é a mais contraintuitiva: \textbf{missão bloqueada entra no índice de propósito}. O que se quer evitar é reinventar uma missão que já foi tentada --- e uma tentativa que falhou é exatamente o que ninguém deve repetir sem saber.

\section{A consequência de método}

Há uma mudança que vale mais que qualquer número corrigido: a fonte única da régua passou a ter uma seção de \textbf{registro negativo} --- o que \emph{não} sustenta a régua e não deve voltar (Parte~IV). Sem ela, uma fundamentação refutada tende a reaparecer na revisão seguinte, porque o argumento de autoridade sobrevive ao artigo que o derrubou.

O mesmo vale para este manual. Ele descreve um sistema que muda, e a única forma honesta de descrevê-lo é dizendo, a cada regra, de onde ela veio --- para que a próxima pessoa saiba o que precisa refutar para mudá-la.

\parteApendice{Apêndices}{A árvore de artefatos, as 25 seções do PRD, o glossário e as fontes.}

\appendix
\renewcommand{\thesection}{\Alph{section}}
\setcounter{section}{0}

\section{A árvore de artefatos de uma execução}
\label{ap:artefatos}

Toda execução persiste uma árvore isolada em \cod{output/<data-hora>-<identificador>/}.

\begin{lstlisting}
output/20260925-101500-exp02ef1b/
|-- PRD.md                          # o PRD padrao ouro (25 secoes)
|-- one-pager-pedagogico.{md,docx,html}   # o documento de validacao humana
|-- diagnostico-execucao.md         # o dossie da execucao
|-- aprendizado.md                  # padroes, lacunas e candidatos a criterio
|-- catalogo-ficha.md               # a ficha desta missao no indice geral
|-- briefing-original.md            # o insumo que voce enviou
|-- estado.json                     # o estado completo da sessao
|-- execucao.json                   # progresso da etapa em segundo plano
|-- conversa.json                   # historico, pareceres e passagens
|-- manifesto-artefatos.json        # metadados de todos os arquivos gerados
|
|-- produto/
|   |-- 00-material-anexado.md      # o material do bimestre, como a esteira o leu
|   |-- 00a-nucleo-narrativo.md     # descoberta, conceitos, anatomia do aha
|   |-- 00b-micro-gate.md           # spoiler, gap, escolha, job e semelhanca
|   |-- 00-one-pager.{md,docx,html} # one-pager canonico (18 secoes)
|   |-- 00c-one-pager-original.md   # a versao antes do revisor
|   |-- 00d-one-pager-resumido.{md,docx,html}  # o formato do time de produto
|   |-- 00e-contrato-produto.md     # lint de contrato de plataforma
|   |-- 00f-brief-escopo.md         # o brief de escopo, separado
|   |-- 01-gate-estrutural.md       # os criterios e o envelope
|   |-- 02-avaliacao-pedagogica.md  # BNCC, professor e guardrails
|   |-- 03-experiencia-aluno.md     # carga cognitiva, faixa e agencia
|   |-- 03b-adequacao-pedagogica.md # parecer em seis eixos
|   |-- 04-gate-final.md            # liberacao formal de produto
|   |-- 05-spec-missao.md           # o contrato de construcao
|   |-- 05a-spec-missao-gate0.md    # a spec antes do fechamento de lacunas
|   |-- 05b-conformidade-bncc.md    # bloco 11 x lista oficial do MEC
|   |-- 05c-matriz-evidencia.md     # grau, sinais e gabarito por habilidade
|   |-- 05d-perguntas-de-produto.md # as perguntas, a decisao e o corte
|   +-- 05e-juiz-perguntas-de-produto.md
|
|-- design/
|   |-- 00-fechamento-lacunas.md    # decisoes e pedidos de re-spec
|   |-- 00a-relatorio-lacunas.md    # o mesmo, em PT-BR, para produto
|   |-- 00b-ambiente-simulado.md    # a ficha do ambiente e a classe
|   |-- 01-briefing-oficio.md       # arco, pico e beat memoravel
|   |-- 02-jornada-telas.md         # raias x 6 etapas + visao do professor
|   |-- 03-fluxo-arestas.md         # grafo do estudante + grafo do professor
|   |-- 04-copy-strings.md          # dicionario de strings (passada 1)
|   |-- 04b-camada-caa.md           # simbolos, modos e pranchas
|   |-- 05-wireframes-cobertura.md  # cobertura das telas e do DS
|   |-- 06-copy-strings-revisada.md # revisao em contexto de tela
|   |-- 06b-lint-prototipo.md       # o que o QA nao ve por ler so o texto
|   |-- 07-qa-entregaveis.md        # auditoria de fidelidade
|   |-- 08-prototipo.html           # <- abre com duplo clique
|   |-- 09-diagramas.html           # jornada e fluxo remontados
|   |-- 10-stitch-telas.md          # tabela de telas e link do projeto
|   +-- 11-miro-boards.md           # link do board espelho
|
+-- dev/
    |-- 00-ficha-manifesto.md       # bloco 0 x listas fechadas
    |-- 01-handoff-dev.md           # stories fechadas e telemetria
    |-- 01b-brief-assets.md         # o que produzir e o que reprova
    |-- 01c-telemetria-check.md     # contrato de evento e arquitetura
    |-- 01d-rastro-lacunas.md       # o que faltou, por dono
    |-- 02a..02e-banca-*.md         # as cinco dimensoes
    |-- 03-banca-consolidada.md     # veredito unico: a pior das cinco
    |-- 04-pendencias-do-prd.md     # marcadores e consistencia interna
    +-- 04b-kpis-check.md           # coerencia KPI <-> evento <-> tela
\end{lstlisting}

\Needspace*{8\baselineskip}
\section{As 25 seções do PRD padrão ouro}
\label{ap:prd}

O \cod{PRD.md} segue um gabarito único, que sintetiza métodos de especificação de produto e de engenharia --- EARS para requisitos, Gherkin para critérios de aceite, C4 para arquitetura, INVEST para \emph{stories} --- e os requisitos do art.~14 da LGPD.

\begin{center}
\begin{longtable}{>{\centering\arraybackslash\sffamily\bfseries\footnotesize\color{unLavanda}}p{0.7cm}>{\sffamily\bfseries\footnotesize\color{unNavy}}p{4.5cm}p{9.0cm}}
\toprule
\sffamily\bfseries\footnotesize\color{unNavy}\# & \sffamily\bfseries\footnotesize\color{unNavy}Seção & \sffamily\bfseries\footnotesize\color{unNavy}O que traz \\
\midrule
\endfirsthead
\toprule
\sffamily\bfseries\footnotesize\color{unNavy}\# & \sffamily\bfseries\footnotesize\color{unNavy}Seção & \sffamily\bfseries\footnotesize\color{unNavy}O que traz \\
\midrule
\endhead
1 & Frontmatter & Metadados, status dos gates, versão e data. \\
2 & Resumo executivo & Problema, solução, público e métrica de sucesso. \\
3 & Objetivos e KPIs & Objetivos pedagógicos e de engajamento, mensuráveis. \\
4 & Não-objetivos & Os limites explícitos do escopo. \\
5 & Personas e contexto & Estudante e docente, na sala real. \\
6 & As seis etapas & O arco narrativo, passo a passo. \\
7 & Máquina de estados & Diagrama e tabela de transições. \\
8 & \emph{User stories} & Histórias com critérios de aceite inequívocos. \\
9 & Requisitos funcionais & Sintaxe formal EARS. \\
10 & Requisitos não-funcionais & Desempenho, acessibilidade e disponibilidade. \\
11 & UX e UI & Links do Stitch, do board e dos diagramas; mapeamento de componentes do \emph{design system}. \\
12 & Modelo de dados & Esquema relacional das entidades da missão. \\
13 & Contratos de API & DTOs e \emph{endpoints} documentados. \\
14 & Telemetria e eventos & Dicionário de eventos analíticos e operacionais. \\
15 & Alinhamento à BNCC & Habilidades mapeadas e rubrica formativa. \\
16 & Dependências & Bibliotecas e serviços necessários. \\
17 & Arquitetura da solução & Componentes e comunicação entre a missão e a plataforma. \\
18 & Plano de implementação & Tarefas em caminhos reais do monorepo. \\
19 & Tarefas com verificação & O comando exato que valida cada tarefa. \\
20 & Riscos e pendências & Matriz de riscos e itens de decisão pendente. \\
21 & Decisões arquiteturais & O registro do porquê de cada escolha técnica. \\
22 & Estratégia de testes & Unitário, integração e ponta a ponta. \\
23 & \emph{Definition of done} & Critérios objetivos de conclusão. \\
24 & Plano de \emph{rollout} & Lançamento gradual e \emph{feature flags}. \\
25 & Glossário e apêndices & Terminologia do universo da missão. \\
\bottomrule
\end{longtable}
\end{center}

\section{Glossário}

\begin{description}[leftmargin=3.5cm,style=sameline,font=\sffamily\bfseries\color{unNavy},itemsep=3pt]
  \item[Agente] Uma posição numerada da esteira. Pode ser um modelo com \emph{prompt} próprio ou um passo determinístico em Python.
  \item[Anatomia do ``ahá''] Os seis passos que um ``ahá'' genuíno tem: crença, aposta, consequência, virada, nomeação e transferência. É tabela obrigatória do núcleo narrativo.
  \item[Ambiente simulado] A casca da missão: a seção do \emph{design system} que reproduz Instagram, YouTube, TikTok, WhatsApp ou ChatGPT.
  \item[Banca] Os cinco avaliadores paralelos da etapa de Dev --- produto, design, pedagogia, dados/LGPD e arquitetura.
  \item[CAA] Comunicação aumentativa e alternativa: símbolos, pranchas e modos de resposta para estudantes que não respondem por texto.
  \item[Conceito limiar] O conceito do material do bimestre que, uma vez compreendido, muda como o estudante vê o resto. Sai do mapa conceitual e tem de viver no Decidir ou no ``ahá''.
  \item[Descoberta central] O que o estudante sabe ao fim e não sabia no início --- e que a missão jamais pode enunciar antes de ele agir.
  \item[Envelope JSON] O bloco padronizado que fecha todo relatório de portão, e que o orquestrador lê para decidir o fluxo.
  \item[Esteira estrita] Modo padrão: reprovação marca herança degradada e o PRD sai bloqueado, em vez de virar ressalva.
  \item[Grau de medida] \cod{scored}, \cod{indicative}, \cod{zero} ou \cod{none}. Define o que um número sobre uma criança pode afirmar --- e a distinção entre \cod{zero} e \cod{none} é regra de proteção, não de contabilidade.
  \item[Gate, portão] O passo que julga o trabalho de outro agente. Determinístico (código) ou avaliador (rubrica).
  \item[\emph{Gap} de informação] A lacuna deliberada entre o que o estudante sabe e o que a missão vai revelar. É o que faz a virada funcionar.
  \item[Impedimento] Achado que impede construir. Distinto de ressalva (risco aceito) e de decisão humana (fora da autoridade da esteira).
  \item[Micro-gate] O portão do agente 02, que audita o núcleo narrativo antes de qualquer construção.
  \item[Ofício] O papel que o estudante assume na simulação, e que justifica as ferramentas que ele usa.
  \item[\emph{One-pager}] O documento canônico de 18 seções (Produto) e, no fim da esteira, o documento de validação humana (agente 49).
  \item[Passagem] O momento entre duas etapas: semáforo dos gates, parecer do dono, recebimento pelo dono seguinte. Registrada no histórico da missão.
  \item[Pergunta de produto] Pergunta que, respondida, muda uma decisão nossa. Tem três partes --- a pergunta, a decisão que dispara e o corte. Sem decisão, é curiosidade.
  \item[Orquestrador] O código que decide a ordem, os ciclos e as repetições. Não é um agente e não usa modelo de linguagem.
  \item[PRD] \emph{Product requirements document}: o entregável de 25 seções que alimenta quem codifica.
  \item[\cod{[PROVISÓRIO]}] Marcador de documento construído sobre algo que um portão reprovou.
  \item[Recorte da aula] O que a professora parametriza: disciplina, tema, assuntos, contexto e foco. Não é ``um tema'' --- é o que torna a mesma missão utilizável em Português, Ciências ou História.
  \item[Re-\emph{spec}] Pedido formal de Design para que Produto redecida algo que a \emph{spec} decidiu mal --- em vez de Design decidir por conta própria.
  \item[\emph{Spec} da missão] O contrato de 12 blocos (mais dois acrescentados em Design) contra o qual tudo é construído e auditado.
\end{description}

\section{Fontes}

\begin{thebibliography}{99}\footnotesize
\setlength{\itemsep}{2pt}

\bibitem{anderson2001}
ANDERSON, L. W.; KRATHWOHL, D. R. \textit{A Taxonomy for Learning, Teaching, and Assessing}. Longman, 2001.

\bibitem{betts2006}
BETTS, J. et al. The development of sustained attention in children. \textit{Child Neuropsychology}, v. 12, n. 3, p. 205--221, 2006.

\bibitem{biggs1982}
BIGGS, J.; COLLIS, K. \textit{Evaluating the Quality of Learning: the SOLO Taxonomy}. Academic Press, 1982.

\bibitem{bommasani2021}
BOMMASANI, R. et al. \textit{On the Opportunities and Risks of Foundation Models}. arXiv:2108.07258, 2021.

\bibitem{eca1990}
BRASIL. \textit{Lei n.\textsuperscript{o} 8.069, de 13 de julho de 1990} --- Estatuto da Criança e do Adolescente.

\bibitem{lgpd2018}
BRASIL. \textit{Lei n.\textsuperscript{o} 13.709, de 14 de agosto de 2018} --- Lei Geral de Proteção de Dados Pessoais.

\bibitem{cne2022}
BRASIL. Conselho Nacional de Educação. \textit{Resolução CNE/CEB n.\textsuperscript{o} 1, de 4 de outubro de 2022}: normas sobre Computação na Educação Básica --- complemento à BNCC. Brasília: MEC, 2022.

\bibitem{bruner1961}
BRUNER, J. S. The act of discovery. \textit{Harvard Educational Review}, v. 31, p. 21--32, 1961.

\bibitem{bruns2018}
BRUNS, B.; COSTA, L.; CUNHA, N. Through the looking glass. \textit{Economics of Education Review}, v. 64, p. 214--250, 2018.

\bibitem{buckingham2007}
BUCKINGHAM, D. \textit{Digital Media Literacies: Rethinking Media Education in the Age of the Internet}. 2007.

\bibitem{castelo2023}
CASTELO, N.; KIM, S.; CARLSON, S. Choice overload em crianças. \textit{Frontiers in Developmental Psychology}, 2023. DOI 10.3389/fdpys.2023.1317426

\bibitem{cemri2025}
CEMRI, M.; PAN, M. Z.; YANG, S. et al. \textit{Why Do Multi-Agent LLM Systems Fail?} NeurIPS 2025, Datasets and Benchmarks Track. arXiv:2503.13657.

\bibitem{cowan2016}
COWAN, N. Working memory maturation. \textit{Perspectives on Psychological Science}, v. 11, n. 2, p. 239--264, 2016.

\bibitem{dehallucinator2024}
\textit{De-Hallucinator}: recuperação iterativa de referências de API do próprio projeto para reduzir alucinação em geração de código. arXiv:2401.01701.

\bibitem{demetriou2022}
DEMETRIOU, A. et al. Cognitive developmental cycles. \textit{Frontiers in Psychology}, v. 13, 954971, 2022.

\bibitem{dewey1938}
DEWEY, J. \textit{Experience and Education}. Kappa Delta Pi, 1938.

\bibitem{doyle1977}
DOYLE, W.; PONDER, G. The practicality ethic in teacher decision-making. \textit{Interchange}, v. 8, n. 3, p. 1--12, 1977.

\bibitem{fischer1980}
FISCHER, K. A theory of cognitive development. \textit{Psychological Review}, v. 87, n. 6, p. 477--531, 1980.

\bibitem{gathercole2004}
GATHERCOLE, S. et al. The structure of working memory from 4 to 15 years of age. \textit{Developmental Psychology}, v. 40, n. 2, p. 177--190, 2004.

\bibitem{godwin2016}
GODWIN, K. et al. Off-task behavior in elementary school children. \textit{Learning and Instruction}, v. 44, p. 128--143, 2016.

\bibitem{hagger2016}
HAGGER, M. et al. A multilab preregistered replication of the ego-depletion effect. \textit{Perspectives on Psychological Science}, v. 11, n. 4, 2016.

\bibitem{cbam2015}
HALL, G.; HORD, S. \textit{Implementing Change: Patterns, Principles and Potholes} --- Concerns-Based Adoption Model. 2015.

\bibitem{isbell2015}
ISBELL, E. et al. Visual working memory continues to develop through adolescence. \textit{Frontiers in Psychology}, v. 6, 696, 2015.

\bibitem{ifscale2025}
JAROSLAWICZ, D.; WHITING, B.; SHAH, P.; MAAMARI, K. \textit{How Many Instructions Can LLMs Follow at Once?} arXiv:2507.11538, 2025.

\bibitem{jonsson2007}
JONSSON, A.; SVINGBY, G. The use of scoring rubrics: reliability, validity and educational consequences. \textit{Educational Research Review}, v. 2, n. 2, p. 130--144, 2007.

\bibitem{kalyuga2003}
KALYUGA, S. et al. The expertise reversal effect. \textit{Educational Psychologist}, v. 38, n. 1, p. 23--31, 2003.

\bibitem{knight1986}
KNIGHT, J. C.; LEVESON, N. G. An experimental evaluation of the assumption of independence in multiversion programming. \textit{IEEE Transactions on Software Engineering}, v. SE-12, n. 1, p. 96--109, 1986.

\bibitem{knoop2018}
KNOOP-VAN CAMPEN, C.; SEGERS, E.; VERHOEVEN, L. The modality and redundancy effects in dyslexia. \textit{Dyslexia}, v. 24, n. 2, p. 140--155, 2018.

\bibitem{livingstone2017}
LIVINGSTONE, S.; THIRD, A. Children and young people's rights in the digital age. \textit{New Media \& Society}, v. 19, n. 5, 2017.

\bibitem{liu2024}
LIU, N. F.; LIN, K.; HEWITT, J. et al. Lost in the middle: how language models use long contexts. \textit{TACL}, v. 12, p. 157--173, 2024.

\bibitem{loewenstein1994}
LOEWENSTEIN, G. The psychology of curiosity. \textit{Psychological Bulletin}, v. 116, n. 1, p. 75--98, 1994.

\bibitem{merrill2002}
MERRILL, M. D. First principles of instruction. \textit{ETR\&D}, v. 50, n. 3, p. 43--59, 2002.

\bibitem{moreno2002}
MORENO, R.; MAYER, R. Verbal redundancy in multimedia learning. \textit{Journal of Educational Psychology}, v. 94, n. 1, p. 156--163, 2002.

\bibitem{osterhaus2021}
OSTERHAUS, C.; KOERBER, S. The development of advanced theory of mind. \textit{Child Development}, v. 92, n. 5, p. 1872--1888, 2021.

\bibitem{papert1980}
PAPERT, S. \textit{Mindstorms: Children, Computers, and Powerful Ideas}. Basic Books, 1980.

\bibitem{parnas1972}
PARNAS, D. L. On the criteria to be used in decomposing systems into modules. \textit{Communications of the ACM}, v. 15, n. 12, p. 1053--1058, 1972.

\bibitem{rozendaal2011}
ROZENDAAL, E. et al. Children's understanding of advertisers' persuasive tactics. \textit{International Journal of Advertising}, v. 30, n. 2, p. 329--350, 2011.

\bibitem{speckit2026}
\textit{Spec Kit Agents}: grounding e validação contra o repositório em desenvolvimento dirigido por especificação. arXiv:2604.05278.

\bibitem{starks2024}
STARKS, D.; REICH, J. ``What about data?'': youth understanding of algorithmic systems. \textit{Information and Learning Sciences}, v. 125, n. 9, p. 673--692, 2024.

\bibitem{sweller1988}
SWELLER, J. Cognitive load during problem solving: effects on learning. \textit{Cognitive Science}, v. 12, n. 2, p. 257--285, 1988.

\bibitem{sweller2019}
SWELLER, J.; VAN MERRIËNBOER, J.; PAAS, F. Cognitive architecture and instructional design: 20 years later. \textit{Educational Psychology Review}, v. 31, n. 2, p. 261--292, 2019.

\bibitem{tran2026}
TRAN, D.; KIELA, D. \textit{Single-Agent LLMs Outperform Multi-Agent Systems on Multi-Hop Reasoning Under Equal Thinking Token Budgets}. arXiv:2604.02460, 2026.

\bibitem{unesco2018}
UNESCO. \textit{A Global Framework of Reference on Digital Literacy Skills for Indicator 4.4.2}. UNESCO Institute for Statistics, 2018.

\bibitem{vygotsky1978}
VYGOTSKY, L. S. \textit{Mind in Society}. Harvard University Press, 1978.

\bibitem{zheng2023}
ZHENG, L.; CHIANG, W.-L.; SHENG, Y. et al. Judging LLM-as-a-judge with MT-Bench and Chatbot Arena. NeurIPS 2023. arXiv:2306.05685.

\end{thebibliography}

\clearpage
\thispagestyle{empty}
\null\vfill
\begin{center}
\includegraphics[width=4.6cm]{marca/logotipo-navy.png}\\[1.4em]
{\color{unLavanda}\rule{3.0cm}{1.2pt}}\\[1.2em]
{\sffamily\color{unNavy}\small Manual do Gerador de Missões}\\[0.3em]
{\color{unLavanda}\footnotesize Setembro de 2026 · documento de trabalho, uso interno}\\[0.3em]
{\color{unLavanda}\footnotesize Unooria Educação e Tecnologia}
\end{center}
\vfill

\end{document}
